\documentclass[11pt,letterpaper]{article}

% ---------------------------------------------------------------------------
% Nixie Clock -- Design Analysis
%
% Derives the board's key electrical characteristics from first principles
% and from component datasheets on file in docs/datasheets/. This document
% works at the schematic level only; it is the source from which the
% eventual Hardware User Guide's tables are extracted.
% ---------------------------------------------------------------------------

\usepackage[letterpaper,margin=1in]{geometry}

% Project-local font (PT Serif: Latin + Cyrillic, OFL-licensed, committed
% under fonts/ so the build has no system-font dependency). Requires
% XeLaTeX or LuaLaTeX -- fontspec does not work under pdfLaTeX.
\usepackage{fontspec}
\setmainfont{PTSerif}[
  Path = fonts/,
  Extension = .ttf,
  UprightFont = *-Regular,
  BoldFont = *-Bold,
  ItalicFont = *-Italic,
  BoldItalicFont = *-BoldItalic,
]

\usepackage{amsmath,amssymb}
\usepackage{siunitx}
\usepackage{booktabs}
\usepackage{longtable}
\usepackage{array}
\usepackage{graphicx}
\usepackage{enumitem}
\usepackage[colorlinks=true,linkcolor=blue,citecolor=blue,urlcolor=blue]{hyperref}
\usepackage{cleveref}

\sisetup{
  per-mode = symbol,
  detect-all = true,
  table-align-text-post = false,
}

% Consistent units used throughout
\DeclareSIUnit\decade{decade}

% Shorthand for a datasheet citation inside table cells and prose.
% Usage: \ds{SN74141\_TI.pdf}{2}
\newcommand{\ds}[2]{\textit{#1}, p.~#2}

\newcolumntype{L}[1]{>{\raggedright\arraybackslash}p{#1}}

% \DesignRev is generated by scripts/make_release.py --sync-doc-rev from
% the KiCad project's REV text variable -- do not hand-edit revision.tex.
% Generated by scripts/make_release.py --sync-doc-rev -- do not hand-edit.
\newcommand{\RevFace}{B}
\newcommand{\RevHv}{C}
\newcommand{\RevMain}{B}
\newcommand{\DesignRev}{face~B / hv~C / main~B}


\title{Nixie Clock --- Design Analysis}
\author{}
\date{Rev \DesignRev{} --- compiled \today}

\begin{document}
\maketitle

\begin{abstract}
\noindent
This document derives the key electrical characteristics of the Nixie
Clock board from the datasheets of its major components and from
first-principles circuit analysis. The analysis stays at the
schematic level: values follow from component ratings and circuit
topology rather than from a specific PCB footprint, except where a
part is inseparable from one particular package. Every numeric claim
about a named, design-driving component traces to a page in a
datasheet under \texttt{docs/datasheets/}; anything not covered by a
datasheet there is sized to standard engineering practice instead. The
document proceeds subsystem by subsystem, working backward from the
load: the Nixie tubes and their binary-coded-decimal (BCD)
decoder/drivers set the
requirements for the 170~V anode supply and for the 5~V rail that
powers those same drivers, the ESP32 MCU sets the requirements for the
3.3~V rail, and all three converters together set the requirements
for the USB-PD front end that supplies them.
\end{abstract}

\tableofcontents

\section{Introduction}
\label{sec:introduction}

The board exists to drive six cold-cathode gas-discharge indicator
tubes from a USB-C Power Delivery source: four IN-12B numeric tubes
forming a HH:MM digit display, and two IN-6 tubes forming a colon
separator between the hour and minute pairs. Each numeric tube is
paired one-to-one with a K155ID1 BCD decoder/driver, which
selects one of ten cathodes per digit under control of an MCU. Every
other subsystem on the board --- the 170~V anode supply, the 5~V rail
for the decoder/drivers, the 3.3~V rail for the MCU, and the USB-PD
front end that supplies all three --- exists only to put a stable,
correctly proportioned high voltage across these tubes and to switch
it through these driver chips. The tubes and drivers are therefore the
correct starting point for a design analysis: their datasheet ratings
set the requirements that every upstream converter stage must satisfy,
rather than the other way around.

This document derives those requirements first (\cref{sec:tubes,sec:bcd-driver}),
then works backward through the 170~V flyback converter that meets
them (\cref{sec:hv-components,sec:hv-converter}), the 5~V rail that
powers the decoder/drivers and the 3.3~V rail that powers the MCU
(\cref{sec:lv-components,sec:5v-converter,sec:3v3-converter}), and
finally the USB-PD input stage sized to supply all three downstream
converters.

\paragraph{Method.} Every parameter attributed to a named,
design-driving component --- a Nixie tube, the BCD decoder/driver, the
flyback controller, the flyback transformer --- is transcribed from a
datasheet on file in \texttt{docs/datasheets/} and cited by filename
and page. Where a datasheet gives a range, a tolerance, or a specific
test condition, that context is preserved rather than collapsed to a
single number, because the converter design in
\cref{sec:hv-converter} depends on knowing which figures are
guaranteed worst cases and which are typical or test-condition values.
\texttt{docs/datasheets/} is the boundary of that treatment: a part
gets a cited, traceable figure if and only if its datasheet lives
there, and everything else --- resistors, capacitors, standard diodes,
and any other component the design doesn't hinge on --- is sized to
ordinary engineering practice instead. Package and hand-soldering
constraints enter only where a single critical part is available in
one specific package (for example, the flyback transformer or the
load-switch dual MOSFET); otherwise the analysis stays at the
schematic level.

\paragraph{A note on the BCD driver substitution.} No datasheet for
the K155ID1 itself --- a Soviet-made functional clone of the Texas
Instruments SN74141 --- could be located. By explicit decision, this
document treats the K155ID1 as electrically equivalent to a generic
TI SN74141 and cites that datasheet in its place
(\cref{sec:bcd-driver}). This is a documented substitution, not a
verified characteristic of the part actually populated on the board,
and it is flagged again at the point where it matters most: the
driver's off-state output voltage rating.

\section{Nixie tube characteristics}
\label{sec:tubes}

The design carries six cold-cathode gas-discharge indicator tubes, forming
an HH:MM digit display and a colon separator between the hour and minute
pairs. This document assigns the designators used throughout to refer to
them:

\begin{table}[htbp]
\centering
\caption{Nixie tube designators}\label{tab:tube-designators}
\begin{tabular}{|l|l|}
\hline
\textbf{Designator} & \textbf{Component} \\
\hline
NX1--NX4 & IN-12B digit tubes (HH:MM display) \\
\hline
NX5--NX6 & INS-1 colon tubes (colon separator) \\
\hline
\end{tabular}
\end{table}

Both are Soviet-made parts documented only in scanned factory sheets;
the figures below are transcribed directly from those scans, with the
original Russian parameter names noted where a term's meaning is not
self-evident from context. Neither sheet carries guaranteed
min/typ/max columns throughout. Several entries below are single-sided
limits (``not more than'' or ``not less than''), which bears on how
conservatively the anode supply must be sized.

\subsection{IN-12B (digit tubes, NX1--NX4)}
\label{sec:in12b}

Source: \texttt{IN-12A\_IN-12B.pdf}, referenced below as [RU] (original
Soviet factory sheet, covering both IN-12A and IN-12B; IN-12B adds a
decimal-point cathode that IN-12A lacks), and
\texttt{IN-12\_english\_gra-afch.pdf}, referenced below as [EN] (an
English-language summary sheet, used as a cross-check).

{\footnotesize
\begin{longtable}{|L{3.6cm}|c|c|c|c|L{4.4cm}|c|}
\caption{IN-12B electrical and physical characteristics}\label{tab:in12b}\\
\hline
\textbf{Parameter} & \textbf{Min} & \textbf{Typ} & \textbf{Max} & \textbf{Units} & \textbf{Conditions} & \textbf{Source} \\
\hline
\endfirsthead
\caption[]{IN-12B electrical and physical characteristics (continued)}\\
\hline
\textbf{Parameter} & \textbf{Min} & \textbf{Typ} & \textbf{Max} & \textbf{Units} & \textbf{Conditions} & \textbf{Source} \\
\hline
\endhead
Maintaining (plate) voltage & & \num{145} & & \si{\volt} & ``Typical Plate Voltage'' & [EN] p.1 \\
\hline
Ignition voltage & & & \num{170} & \si{\volt} & discharge-onset voltage (``Напряжение возникновения разряда''); not split by variant & [RU] p.4 \\
\hline
Cathode (digit) current, DC operating & \num{2.5} & & \num{3.5} & \si{\milli\ampere} & ``Ток рабочий, постоянный, для цифр'' & [RU] p.4 \\
\hline
Cathode (digit) current, average (unfiltered half-wave AC) & \num{1} & & \num{2} & \si{\milli\ampere} & ``Ток рабочий, средний\ldots для цифр'' & [RU] p.4 \\
\hline
Cathode (digit) overload current & & & \num{5} & \si{\milli\ampere} & Time-limited; see overload duration below & [RU] p.4 \\
\hline
Decimal-point cathode current, DC operating & \num{0.3} & & \num{0.5} & \si{\milli\ampere} & IN-12B only (pin 12); IN-12A pin 12 unconnected & [RU] p.4 \\
\hline
Decimal-point overload current & & & \num{0.7} & \si{\milli\ampere} & & [RU] p.4 \\
\hline
Overload duration limit, IN-12A & & & \num{30} & min & Applies to the \SI{5}{\milli\ampere} digit overload above & [RU] p.5 \\
\hline
Overload duration limit, IN-12B & & & \num{33} & min & Applies to the \SI{5}{\milli\ampere} digit / \SI{0.7}{\milli\ampere} decimal-point overload above & [RU] p.5 \\
\hline
Anode supply voltage & \num{200} & \num{250} & & \si{\volt} & DC or half-wave-rectified AC; the \num{250}~V recommended value reduces discharge-onset delay time & [RU] p.4, p.5 \\
\hline
Conditioning current, pulsed & \num{1} & & \num{2} & \si{\milli\ampere} & $\geq$\SI{1}{\minute}/cathode, applied once after storage & [RU] p.5 \\
\hline
Conditioning current, DC & \num{2.5} & & \num{3} & \si{\milli\ampere} & $\geq$\SI{1}{\minute}/cathode, applied once after storage & [RU] p.5 \\
\hline
Continuous single-cathode life rating & & & \num{1000} & h & ``Longevity'' mode, one cathode continuously lit & [RU] p.5 \\
\hline
Digit height & & \num{18} & & mm & & [EN] p.1 \\
\hline
Envelope diameter & \num{29} & & \num{31} & mm & $-2$~mm tolerance [RU]; [EN] gives \SI{32}{\milli\meter} & [RU] p.3, [EN] p.1 \\
\hline
\end{longtable}
}

\paragraph{Not specified in either source.} A combined (total) anode
current limit distinct from the per-cathode figures; an interelectrode
or anode-cathode capacitance; a manufacturer-recommended ballast
resistor value or target anode current within the permissible 2.5--3.5~mA
band.

\paragraph{Open question: anode supply voltage.} The datasheet's
minimum anode supply figure (\SI{200}{\volt}, recommended
\SI{250}{\volt}) sits \emph{above} the \SIrange{164.8}{191.5}{\volt}
anode-rail range this document derives (\cref{sec:hv-converter}), so
the design does not reach even the stated minimum. The likely
explanation is that the factory figure is written for the tube's other
stated supply mode, unfiltered half-wave-rectified AC
(\S\ref{sec:in12b}, ``Cathode current, average''). An unsmoothed
half-wave supply spends roughly half of every mains cycle near zero
volts, so its \emph{peak} must be pushed well above the
\SI{145}{\volt} typical maintaining voltage to sustain adequate
\emph{average} current and reliable strike. A continuously present,
regulated DC rail carries no such requirement. The source does not
explain the discrepancy, so this reading is unconfirmed.
\Cref{sec:hv-converter} carries it forward as an open item; the
bench-tuning procedure (raising the trim potentiometer until every
digit of the actual tube stock strikes reliably, then setting the rail
roughly \SI{10}{\volt} above that point) tests the uncertainty
directly, whichever reading is correct.

\subsection{INS-1 (colon tubes, NX5--NX6)}
\label{sec:ins1}

Source: \texttt{INS-1.pdf} (original Soviet factory label sheet,
ODO.334.095 TU). The tube has two electrode leads, counted from the
indicator dot: lead 1 is the anode (A, the small cylinder) and lead 2
the cathode (K, the large cylinder). This assignment follows the
datasheet's electrode table. There is no auxiliary or priming
electrode.

{\footnotesize
\begin{longtable}{|L{4.6cm}|c|c|c|c|L{3.6cm}|c|}
\caption{INS-1 electrical characteristics}\label{tab:ins1}\\
\hline
\textbf{Parameter} & \textbf{Min} & \textbf{Typ} & \textbf{Max} & \textbf{Units} & \textbf{Conditions} & \textbf{Page} \\
\hline
\endfirsthead
\caption[]{INS-1 electrical characteristics (continued)}\\
\hline
\textbf{Parameter} & \textbf{Min} & \textbf{Typ} & \textbf{Max} & \textbf{Units} & \textbf{Conditions} & \textbf{Page} \\
\hline
\endhead
Striking voltage & \num{65} & & \num{90} & \si{\volt} & Discharge-onset voltage (``Напряжение возникновения разряда''); no test current given & 2 \\
\hline
Maintaining voltage & & & \num{55} & \si{\volt} & ``Напряжение поддержания разряда''; single-sided limit, no minimum given & 2 \\
\hline
Optimal indication current & & \num{0.5} & & \si{\milli\ampere} & ``Оптимальным током, обеспечивающим достаточную яркость свечения''; a brightness recommendation, with no band around it & 3 \\
\hline
\end{longtable}
}

\paragraph{Conditioning after storage.} The datasheet directs that an
indicator taken out of long storage be run at its nominal discharge
current for one to two minutes before normal service, and refers
operation to GOST 11163-84 and OST 11~339.003-75. The IN-12B carries a
comparable instruction (\cref{sec:in12b}).

\paragraph{Not specified.} The sheet gives no minimum for the
maintaining voltage, no rated operating-current band around the
\SI{0.5}{\milli\ampere} brightness figure, no maximum discharge
switching frequency, no interelectrode capacitance, no example ballast
value, and no envelope dimensions. Two consequences follow for the
derivations downstream. \Cref{sec:hv-load-colon} sizes the ballast
against the single \SI{0.5}{\milli\ampere} figure and reports the
spread the trim range produces around it, having no band to check
against. And because \SI{55}{\volt} is a ceiling with no floor beneath
it, a tube maintaining below that figure draws proportionally more
current through a fixed ballast, which that section treats as the open
side of the result.

\section{BCD decoder/driver characteristics}
\label{sec:bcd-driver}

Each IN-12B digit tube (\textsc{nx1}--\textsc{nx4}) is paired
one-to-one with a K155ID1 BCD decoder/driver
(\textsc{u1}--\textsc{u4}), a Soviet-made functional clone of the
Texas Instruments SN74141. Each of the driver's ten outputs is an
open-collector transistor wired directly to one cathode of its tube;
the selected (lit) output sinks current, pulling that cathode low,
while the other nine sit in a high-impedance off state. There is no
interposing high-voltage buffer transistor between driver and tube in
this design --- the driver's own output stage must tolerate whatever
voltage an unselected cathode presents.

\paragraph{Substitution.} No datasheet for the K155ID1 itself could be
located in \texttt{docs/datasheets/}. By explicit decision recorded
for this document, the K155ID1 is treated as electrically equivalent
to a generic TI SN74141, and the figures below are transcribed from
that part's datasheet (\texttt{SN74141\_TI.pdf}, Texas Instruments,
originally published 1972, revised 1983) instead. Every value in this
section is therefore a documented substitution, not a verified
characteristic of the part actually populated on the board.

\begin{longtable}{L{4.4cm}L{2.0cm}L{2.6cm}L{4.0cm}L{1.2cm}}
\toprule
Parameter & Value & Category & Conditions & Page \\
\midrule
\endhead
Supply voltage \(V_{CC}\) & \SI{7}{\volt} & Abs.\ max & w.r.t.\ network ground & 2 \\
Input voltage & \SI{5.5}{\volt} & Abs.\ max & & 2 \\
Current into any output (off-state) & \SI{2}{\milli\ampere} & Abs.\ max & & 2 \\
Operating free-air temperature & \SIrange{0}{70}{\celsius} & Abs.\ max & & 2 \\
Storage temperature & \SIrange{-65}{150}{\celsius} & Abs.\ max & & 2 \\
Supply voltage \(V_{CC}\), Min/Nom/Max & 4.75~/~5~/~5.25~V & Rec.\ operating & & 2 \\
\textbf{Off-state output voltage} & \textbf{\SI{60}{\volt} max} & \textbf{Rec.\ operating} & & \textbf{2} \\
\(V_{IH}\), high-level input voltage & \SI{2}{\volt} min & Electrical char.\ & & 2 \\
\(V_{IL}\), low-level input voltage & \SI{0.8}{\volt} max & Electrical char.\ & & 2 \\
\(V_{IK}\), input clamp voltage & \SI{-1.5}{\volt} max & Electrical char.\ & \(V_{CC}=\)min, \(I_I=-5\)~mA & 2 \\
\(V_{O(on)}\), on-state output voltage & \SI{2.5}{\volt} max & Electrical char.\ & \(V_{CC}=\)min, \(I_O=7\)~mA & 2 \\
\textbf{\(V_{O(off)}\), off-state output voltage, counts 0--9} & \textbf{\SI{60}{\volt} min} & \textbf{Electrical char.\ (guaranteed)} & \(V_{CC}=\)max, \(I_O=0.5\)~mA & \textbf{2} \\
\(I_{O(off)}\), off-state reverse current & \SI{50}{\micro\ampere} max & Electrical char.\ & \(V_{CC}=\)max, \(V_O=55\)~V, \(T_A=55\)\si{\celsius} & 2 \\
\(I_{O(off)}\), off-state reverse current, counts 10--15 & 5 (typ) / 15 (max) \si{\micro\ampere} & Electrical char.\ & \(V_O=30\)~V, \(T_A=70\)\si{\celsius} & 2 \\
\(I_I\), input current at max input voltage & \SI{1}{\milli\ampere} max & Electrical char.\ & \(V_{CC}=\)max, \(V_I=5.5\)~V & 2 \\
\(I_{IH}\), A input & \SI{40}{\micro\ampere} max & Electrical char.\ & \(V_{CC}=\)max, \(V_I=2.4\)~V & 2 \\
\(I_{IH}\), B/C/D input & \SI{80}{\micro\ampere} max & Electrical char.\ & \(V_{CC}=\)max, \(V_I=2.4\)~V & 2 \\
\(I_{IL}\), A input & \SI{-1.6}{\milli\ampere} max & Electrical char.\ & \(V_{CC}=\)max, \(V_I=0.4\)~V & 2 \\
\(I_{IL}\), B/C/D input & \SI{-3.2}{\milli\ampere} max & Electrical char.\ & \(V_{CC}=\)max, \(V_I=0.4\)~V & 2 \\
\(I_{CC}\), supply current & 16 (typ) / 25 (max) \si{\milli\ampere} & Electrical char.\ & \(V_{CC}=\)max, inputs grounded, outputs open & 2 \\
\bottomrule
\end{longtable}

Pin mapping (16-pin DIP): BCD inputs A, B, C, D on pins 3, 6, 7, 4;
\(V_{CC}\) on pin 5, ground on pin 12; decimal outputs 0--9 on pins
16, 15, 8, 9, 13, 14, 11, 10, 1, 2 respectively. Not specified in this
datasheet: package thermal resistance, maximum junction temperature,
ESD rating, and propagation delay.

\subsection{The off-state voltage question}
\label{sec:bcd-voltage-risk}

The datasheet's Absolute Maximum Ratings table gives no distinct
off-state output voltage limit --- it bounds only off-state output
\emph{current} (\SI{2}{\milli\ampere}). The only off-state voltage
figure anywhere in the datasheet is \SI{60}{\volt}, and it is a
\emph{guaranteed operating} point (Recommended Operating Conditions,
and corroborated by the \(V_{O(off)}\) electrical-characteristics row,
which promises no breakdown up to \SI{60}{\volt} under a
\SI{0.5}{\milli\ampere} test load) rather than a stated destructive
limit; the datasheet does not say what happens above it, only that
correct operation is guaranteed up to it. Two separate leakage-current
test points exist below that ceiling (\SI{50}{\micro\ampere} max at
\(V_O=\SI{55}{\volt}\), and \SI{5}{}/\SI{15}{\micro\ampere}
typ/max at \(V_O=\SI{30}{\volt}\) with all ten outputs off) and should
not be mistaken for the guaranteed standoff voltage itself.

This board's anode rail runs \SIrange{160}{196}{\volt}
(\cref{sec:hv-converter}), and with no buffer transistor between
driver and tube, an unselected cathode is exposed to something on the
order of the tube's maintaining voltage up to the full rail voltage
--- \SIrange{145}{196}{\volt} by the figures in \cref{sec:tubes} ---
2.5--3.3$\times$ the datasheet's guaranteed \SI{60}{\volt} standoff.

This is treated in this document as an accepted, community-validated
condition rather than an unresolved defect: direct cathode drive from
a 7441/74141/K155ID1-class open-collector decoder, run well above its
guaranteed off-state rating, is the standard topology used across
essentially all amateur IN-12/IN-14-class Nixie clocks, with a long
track record of field reliability at rail voltages in this same range.
Two effects plausibly explain the margin, neither confirmed from a
cited source: guaranteed datasheet minimums are worst-case
production-lot figures, not the actual breakdown point of any given
part, which in practice typically exceeds the guarantee by a wide
margin; and an unselected cathode inside a lit tube is not necessarily
pulled to the full anode potential the way a simple resistive
off-state model would suggest, since it floats within the tube's own
glow-discharge field rather than being driven there by an external
low-impedance source. This document records the substitution and the
margin explicitly rather than silently assuming the outputs are
safe: it is a documented, precedent-based judgment call, not a
datasheet guarantee, and it remains open pending a datasheet for the
actual populated part (K155ID1) that might state a higher verified
rating.

\section{170~V converter: driving component characteristics}
\label{sec:hv-components}

The 170~V anode rail is produced by a current-mode flyback built
around a TI LM5155 controller (\textsc{u201}) and a Coilcraft
DA2032-AL flyback transformer (\textsc{t201}). Both are named,
design-driving parts, so their characteristics are verified from
datasheet before use in the converter derivation of
\cref{sec:hv-converter}.

\subsection{LM5155 flyback controller}
\label{sec:lm5155}

Source: \texttt{lm5155.pdf} (Texas Instruments, SNVSB75E, Rev.\ E,
Aug.\ 2023). Figures apply identically to the LM5155 and LM51551
variants unless noted.

\begin{longtable}{L{4.6cm}L{3.4cm}L{5.4cm}L{1.2cm}}
\toprule
Parameter & Value & Conditions & Page \\
\midrule
\endhead
FB reference voltage \(V_{REF}\) & 0.99 / 1.00 / 1.01~V (min/typ/max) & \(T_J=-40\) to \SI{125}{\celsius} & 7 \\
Current-limit threshold \(V_{CLTH}\) (CS to PGND) & 93 / 100 / 107~mV (min/typ/max) & \(R_T=9.09\)~k\(\Omega\), \(V_{BIAS}=12\)~V; fixed DC threshold & 7 \\
Peak slope-comp current \(I_{SLOPE}\) (optional, external \(R_{SL}\) only) & 22.5 / 30 / 37.5~\si{\micro\ampere} & \(R_T=220\)~k\(\Omega\); \(I_{SLOPE}=30\,\si{\micro\ampere}\times f_{RT}/f_{SYNC}\) & 7, 20 \\
RT resistor equation & \(R_T = \dfrac{2.21\times10^{10}}{f_{RT}\,[\mathrm{Hz}]} - 955\)~\(\Omega\) & Verified against both characterized points: 220~k\(\Omega\)\(\to\)100~kHz, 9.09~k\(\Omega\)\(\to\)2.2~MHz & 16 \\
Soft-start current \(I_{SS}\) & 9 / 10 / 11~\si{\micro\ampere} (min/typ/max) & & 7 \\
Soft-start time, boost topology & \(t_{SS} = \dfrac{C_{SS}}{I_{SS}}\left(1-\dfrac{V_{SUPPLY}}{V_{LOAD}}\right)V_{REF}\) & No flyback-specific equation given (see note) & 16 \\
Enable (EN) threshold, rising/falling & 0.52 / 0.49~V (typ) & Same physical pin as UVLO/SYNC (pin 12) & 6--7 \\
UVLO/SYNC threshold, rising/falling & 1.5 / 1.45~V (typ) & Separate comparator, same pin as EN & 7 \\
UVLO hysteresis current \(I_{UVLO}\) & 4 / 5 / 6~\si{\micro\ampere} (min/typ/max) & \(V_{UVLO}=1.6\)~V; sets external divider hysteresis & 7, 13 \\
BIAS absolute maximum & \(-0.3\) to \SI{50}{\volt} & & 5 \\
BIAS recommended operating & 3.5--45~V & When VCC is supplied from the internal regulator & 6, 14 \\
VCC (internal gate-drive rail) regulation & 6.5 / 6.85 / 7~V (min/typ/max) & \(V_{BIAS}=8\)~V, no load & 6 \\
VCC regulation, loaded & \(\geq\)6.5~V & \(V_{BIAS}=8\)~V, \(I_{VCC}=35\)~mA & 6 \\
VCC absolute maximum & \(-0.3\) to \SI{18}{\volt}, or \(V_{BIAS}+0.3\)~V, whichever is lower & & 5 \\
PGOOD absolute maximum & \(-0.3\) to \SI{18}{\volt} & Open-drain; sink limited to 1~mA when \(V_{PGOOD}>V_{BIAS}\) & 5 \\
PGOOD recommended pull-up & \(\geq\)10~k\(\Omega\) & & 22 \\
PGOOD assert condition & FB $>$ \(V_{UVTH}\) AND VCC $>$ VCC-UVLO AND UVLO/EN above threshold & 25~\si{\micro\second} deglitch & 22 \\
PGOOD undervoltage threshold \(V_{UVTH}\) & 87 / 90 / 93\% of \(V_{REF}\) (falling) & & 7 \\
Minimum on-time \(t_{ON(MIN)}\) & \SI{50}{\nano\second} (typ) & At \(R_T=9.09\)~k\(\Omega\) (2.2~MHz); \(\approx\)165~ns implied near 150~kHz & 7, 21 \\
Maximum duty cycle \(D_{MAX}\) & 80--90\% (at \(R_T=9.09\)~k\(\Omega\)); 90--96\% (at \(R_T=220\)~k\(\Omega\)) & Frequency-dependent; \(\approx\)91--92\% typ interpolated near 150~kHz & 7 \\
GATE absolute maximum (negative) & \(-1\)~V/100~ns, \(-2\)~V/50~ns transients & No explicit positive max beyond VCC rail & 5 \\
CS absolute maximum & \(-0.3\) to \SI{0.3}{\volt} DC; \(-1\)~V/100~ns, \(-2\)~V/20~ns transients & & 5 \\
RT, FB, SS pin absolute maximum & \(-0.3\) to \SI{3.8}{\volt} & Not intended for external drive & 5 \\
\bottomrule
\end{longtable}

\paragraph{Current limit has no slope compensation as designed.} The
100~mV current-limit threshold is a fixed DC level at the CS pin; the
datasheet's peak-current-limit equation only subtracts a
slope-compensation term if the optional external \(R_{SL}\) resistor
is populated, and TI recommends omitting it for accurate current
limiting versus \(V_{IN}\). A separate, always-on slope ramp
(\(V_{SLOPE}=40\,\text{mV}\times f_{RT}/f_{SYNC}\)) is summed only into
the regulation (PWM/COMP) comparator, to prevent subharmonic
oscillation --- a different comparator from the current-limit path.
This distinction matters in \cref{sec:hv-converter}: the current-sense
resistor should be sized against the flat 100~mV threshold, not
against a slope-adjusted figure.

\paragraph{No flyback-specific soft-start formula.} The datasheet
gives soft-start timing equations for boost and SEPIC topologies only.
Its own non-isolated flyback reference design (Fig.\ 10--12) wires
FB/SS/COMP the same way as the boost case, but the boost formula
assumes an input-to-output pass-through path that a flyback
transformer does not provide, so a soft-start time computed from the
boost equation is an approximation to be verified on the bench, not a
guaranteed figure.

\paragraph{EN and UVLO are distinct comparators on one pin.} Pin 12
(``UVLO/SYNC'') carries two separate thresholds: EN (\(\approx\)0.5~V)
gates shutdown$\leftrightarrow$standby, and UVLO
(\(\approx\)1.5~V/1.45~V) gates standby$\leftrightarrow$run. An
external resistor divider from the input rail to this pin, sized
against the UVLO threshold and hysteresis current, is what sets the
converter's input undervoltage lockout point --- not the EN threshold.

\paragraph{Maximum duty cycle bounds the achievable turns ratio.} The
frequency-dependent \(D_{MAX}\) limit (interpolated to roughly
91--92\% typical near 150~kHz) caps how much the transformer's
reflected voltage can dominate the switching waveform at minimum input
voltage; \cref{sec:hv-converter} checks the design's actual duty cycle
against this limit rather than assuming it is unconstrained.

\subsection{DA2032-AL flyback transformer}
\label{sec:da2032}

Source: \texttt{da2032.pdf} (Coilcraft).

\begin{longtable}{L{4.8cm}L{2.6cm}L{5.0cm}L{1.2cm}}
\toprule
Parameter & Value & Conditions & Page \\
\midrule
\endhead
Turns ratio & 1:10 & Only configuration offered; requires the two primary windings connected in parallel & 1 \\
Primary inductance \(L_P\), at 0~A & \SI{10.0}{\micro\henry} \(\pm10\%\) & 100~kHz, 0.1~V\textsubscript{rms}, 0~A\textsubscript{dc}; primaries in parallel & 1 \\
Primary inductance at \(I_{pk}\) & \SI{9.0}{\micro\henry} min & At \(I_{pk}=3.0\)~A & 1 \\
Leakage inductance, max & \SI{0.150}{\micro\henry} & Referred to primary; primaries in parallel, secondary shorted & 1 \\
Peak primary current & \SI{3.0}{\ampere} & Rating point for the \(L_P\)-at-\(I_{pk}\) spec; no separate \(I_{sat}\) given & 1 \\
DCR, primary (max) & \SI{0.013}{\ohm} & \SI{25}{\celsius}; primaries in parallel & 1 \\
DCR, secondary (max) & \SI{1.60}{\ohm} & \SI{25}{\celsius} & 1 \\
Interwinding isolation & \SI{1500}{\volt_{rms}} & 1-minute withstand & 1 \\
Winding-to-core isolation & \SI{750}{\volt_{rms}} & 1-minute withstand & 1 \\
Interwinding capacitance & \SI{67}{\pico\farad} & 100~kHz, 0.1~V\textsubscript{rms}, pin 3 to pin 1 & 1 \\
Ambient/storage temperature & \(-40\) to \SI{125}{\celsius} (component); \(-40\) to \SI{80}{\celsius} (tape \& reel storage) & & 1 \\
\bottomrule
\end{longtable}

Pin assignment: secondary on pins 1 and 12, primary windings on pins
3--10; dot marks pin 1. No recommended operating frequency range is
published beyond the 100~kHz condition used for the \(L_P\) and
interwinding-capacitance measurements --- the 150~kHz operating point
used in this design is not independently characterized in the
datasheet, and \cref{sec:hv-converter} notes this as an extrapolation.
No direct saturation current is published either; the datasheet's
proxy is the ``inductance at \(I_{pk}\)'' figure (\SI{9.0}{\micro\henry}
min at \SI{3.0}{\ampere}), which this document treats as the effective
peak-current rating. The isolation ratings (1500~V\textsubscript{rms}
winding-to-winding, 750~V\textsubscript{rms} winding-to-core) presume
an isolated application; the datasheet gives no caution against this
design's non-isolated, shared-ground use, since the windings remain
galvanically separate regardless of how the secondary return is
externally referenced --- the design simply forgoes the isolation the
part is capable of providing.

\section{170~V flyback converter derivation}\label{sec:hv-converter}

The 170~V supply uses an LM51561H flyback controller
(\cref{sec:lm5156h}). This section specifies its circuit requirements and
component values. \Cref{sec:tubes} defines NX1 to NX6.
\Cref{tab:hv-designators} defines the remaining designators.

{\footnotesize
\begin{table}[htbp]
\centering
\caption{170~V flyback converter: component designators}\label{tab:hv-designators}
\begin{tabular}{|l|l|}
\hline
\textbf{Designator} & \textbf{Function} \\
\hline
\des{T}{HV} & Flyback transformer \\

\hline
\des{R}{BALDIG} & Digit-tube anode ballast, one per tube (NX1 to NX4) \\

\hline
\des{R}{BALCOL} & Colon-tube anode ballast, one per tube (NX5 and NX6) \\

\hline
\des{Q}{HVSW} & Primary-side switching FET \\

\hline
\des{R}{HVCS} & Current-sense resistor \\

\hline
\des{R}{HVCSF}, \des{C}{HVCSF} & CS-pin filter \\

\hline
\des{R}{HVD1}, \des{R}{HVD2} & Feedback divider, upper (fixed) \\

\hline
\des{R}{HVD3}, \des{RV}{HVTRIM} & Feedback divider, lower (fixed \& trim) \\

\hline
\des{R}{HVRT} & RT (switching-frequency-set) resistor \\

\hline
\des{R}{HVCOMP}, \des{C}{HVC1}, \des{C}{HVC2} & COMP-pin compensation network \\

\hline
\des{R}{HVGON} & Gate-drive resistor (turn-on) \\

\hline
\des{R}{HVGOFF}, \des{D}{HVGOFF} & Turn-off speedup network (bypasses \des{R}{HVGON} on turn-off only) \\

\hline
\des{D}{HVRECT} & Output rectifier diode \\

\hline
\des{R}{HVCL}, \des{C}{HVCL}, \des{D}{HVCL} & RCD clamp (snubber) network \\

\hline
\des{R}{HVBLEED} & Bleeder resistor \\

\hline
\des{C}{HVSS} & Soft-start capacitor \\

\hline
\des{R}{HVUV1}, \des{R}{HVUV2} & Input UVLO divider \\

\hline
\des{Q}{HVINH}, \des{R}{HVINHPU}, \des{R}{HVINHPD} & Display inhibit shunt and its gate divider \\

\hline
\des{Q}{HVREL}, \des{R}{HVRELPD} & Inhibit release switch and its gate pull-down \\

\hline
\des{R}{HVPG} & PGOOD pull-up resistor \\

\hline
\des{C}{HVVCC} & VCC regulator bypass capacitor \\

\hline
\des{C}{HVPUMP}, \des{R}{HVPUMP} & Bias charge-pump coupling capacitor and series resistor \\

\hline
\des{D}{HVPUMP}, \des{D}{HVRESET} & Charge-pump delivery and reset diodes \\
\hline
\des{D}{HVSTART} & BIAS start-up diode from $V_{\mathrm{BUS}}$ \\
\hline
\des{D}{HVZ} & BIAS clamp zener \\
\hline
\des{C}{HVBIAS} & BIAS reservoir capacitor \\
\hline
\des{C}{HVI1}, \des{C}{HVI2} & Input bulk / high-frequency capacitance \\
\hline
\des{C}{HVO1}, \des{C}{HVO2} & Output capacitance \\

\hline
\end{tabular}
\end{table}
}

\begin{figure}[htbp]
  \centering
  \includegraphics[width=\textwidth]{hv}
  \caption{170~V flyback converter: reference schematic}\label{fig:hv}
\end{figure}

\Cref{fig:hv} shows the connections. It does not show the tubes or ballast
resistors \des{R}{BALDIG} and \des{R}{BALCOL}. \Cref{sec:hv-load} derives these
loads on $V_{\mathrm{HV}}$.

The input $V_{\mathrm{BUS}}$ spans \SIrange{4.75}{12.6}{\volt}
(\cref{sec:pd-classes}). The lower limit is the \SI{5}{\volt} contract at
$-5\,\%$. The upper limit is the \SI{12}{\volt} contract at $+5\,\%$. The
nominal input is \SI{9}{\volt}.

Operation at \SI{4.75}{\volt} adds three requirements. The converter needs the
auxiliary bias supply in \cref{sec:hv-house-pump}. It needs the UVLO threshold
in \cref{sec:hv-uvlo}. It also needs the display inhibit in
\cref{sec:hv-inhibit}. The inhibit uses the negotiated current limit because
the contract voltage does not identify the available current.

\subsection{Preliminaries}\label{sec:hv-prelim}

\subsubsection{Switching frequency}\label{sec:hv-fsw-choice}

In DCM, $L_p f_{\mathrm{sw}} I_{\mathrm{pk}}^2 = 2 P_{\mathrm{in}}$. The
transformer sets $L_p$ (\cref{sec:da2032}). The load sets $P_{\mathrm{in}}$
(\cref{sec:hv-budget}). Thus, the switching frequency sets the peak current.
It also sets the margin below the \SI{3.0}{\ampere} saturation rating. Two
constraints limit the frequency. Both scale with $\sqrt{f_{\mathrm{sw}}}$.

Both $t_{\mathrm{on}}$ and $t_{\mathrm{dem}}$ are proportional to
$f_{\mathrm{sw}}^{-1/2}$. The period is proportional to
$f_{\mathrm{sw}}^{-1}$. Therefore, the conduction fraction increases with
$\sqrt{f_{\mathrm{sw}}}$. At the worst condition in \cref{sec:hv-dcm}, DCM
ends at \SI{185}{\kilo\hertz}. The idle fraction falls to $5\,\%$ at
\SI{168}{\kilo\hertz}.

At low frequency, the permitted current-sense resistance decreases
(\cref{sec:hv-cs-sense}). Core protection requires \des{R}{HVCS} to be at least
\SI{37.8}{\milli\ohm}. The power-delivery limit is
$\SI{93}{\milli\volt}/I_{\mathrm{pk}}^{(L\min)}$. This limit increases with
$\sqrt{f_{\mathrm{sw}}}$. The two limits are equal at
\SI{132}{\kilo\hertz}. The specified \SI{39}{\milli\ohm} resistor meets the
power-delivery limit above \SI{141}{\kilo\hertz}.

The loss model does not add a tighter frequency constraint. The design uses
\SI{150}{\kilo\hertz}. This value is inside both permitted ranges. The implemented frequency
is \SI{149.4}{\kilo\hertz}. At this frequency, the commanded peak current is
\SI{2.18}{\ampere}, or $73\,\%$ of $I_{\mathrm{sat}}$.

\subsubsection{RT resistor}\label{sec:hv-fsw}

For the \SI{150}{\kilo\hertz} target of \cref{sec:hv-fsw-choice}, the RT
equation of \cref{sec:lm5156h} gives



\begin{align*}
  R_T &= \frac{2.21\times10^{10}}{f_{RT}} - 955
       = \frac{2.21\times10^{10}}{1.5\times10^{5}} - 955
       = \SI{146.4}{\kilo\ohm}.
\end{align*}
The nearest standard value is $\des{R}{HVRT} = \SI{147}{\kilo\ohm}$. It gives
\begin{align*}
  f_{RT} &= \frac{2.21\times10^{10}}{147\,000 + 955} = \SI{149.4}{\kilo\hertz}
\end{align*}
All subsequent arithmetic uses \SI{149.4}{\kilo\hertz}.


The DA2032-AL is characterized only at \SI{100}{\kilo\hertz}
(\cref{sec:da2032}). The analysis applies its inductance and saturation data at
\SI{149.4}{\kilo\hertz} (inferred). It does not apply the core-loss data at
this frequency. \Cref{sec:hv-op} includes a core-loss allowance.

\subsubsection{Input UVLO}\label{sec:hv-uvlo}

UVLO must release below the \SI{4.75}{\volt} minimum input. The release voltage
must include the LM51561H threshold tolerance. A \SI{4.75}{\volt} release value
could prevent operation at the specified minimum. The design target is
\SI{4.4}{\volt}.

An external divider sets the UVLO/SYNC thresholds. The typical thresholds are
\SI{1.5}{\volt} rising and \SI{1.45}{\volt} falling
(\cref{sec:lm5156h}). The required divider ratio is:
\begin{align*}
  \frac{\des{R}{HVUV1}+\des{R}{HVUV2}}{\des{R}{HVUV2}}
  &= \frac{4.4}{1.5} = 2.93.
\end{align*}

The divider ratio does not set \des{R}{HVUV2}. A value of tens of kilohms gives
low quiescent current. It also gives useful hysteresis from the
\SI{5}{\micro\ampere} pin current (\cref{sec:lm5156h}). The design uses
\SI{32.4}{\kilo\ohm}:
\begin{align*}
  \des{R}{HVUV1} &= \des{R}{HVUV2}\left(\frac{4.4}{1.5} - 1\right)
  = 32.4 \times 1.933 = \SI{62.6}{\kilo\ohm}
  \;\rightarrow\; \SI{63.4}{\kilo\ohm}.
\end{align*}



The same pin has an EN comparator with an approximately
\SI{0.52}{\volt} rising threshold. The divider refers this threshold to
\SI{1.54}{\volt} at $V_{\mathrm{BUS}}$. Thus, EN becomes true before UVLO
releases during a monotonic input ramp. One divider serves both comparators.

\subsubsection{Display inhibit}\label{sec:hv-inhibit}

UVLO prevents operation before the input rail is present. It does not verify
that the source can supply the tube current. The MCU reads the negotiated
USB-PD contract and controls a display-inhibit signal. N-channel MOSFET
\des{Q}{HVINH} is across \des{R}{HVUV2}. It grounds UVLO/SYNC when it is on. The
MCU releases the inhibit. Two requirements set the gate circuit.

\begin{itemize}[nosep]
  \item \textbf{Assert the inhibit before converter operation.} The MCU and
    converter both use $V_{\mathrm{BUS}}$. However, the MCU 3.3~V rail starts
    after $V_{\mathrm{BUS}}$ appears (\cref{sec:3v3-converter}). A gate bias
    from the 3.3~V rail would depend on two converter start times. A gate bias
    from $V_{\mathrm{BUS}}$ rises with the UVLO divider. Resistor ratios then set
    the sequence.
  \item \textbf{Isolate the MCU control pin from the gate node.} The gate bias
    can reach \SI{11.5}{\volt}. The MCU pin limit is \(V_{DD}+0.3\)~V
    (\cref{sec:esp32}). Thus, the control pin and gate cannot use one node.
\end{itemize}

\des{R}{HVINHPU} and \des{R}{HVINHPD} divide $V_{\mathrm{BUS}}$ onto the gate of
\des{Q}{HVINH}. A second N-channel MOSFET, \des{Q}{HVREL}, pulls this gate to
ground when the MCU enables the display. \des{R}{HVRELPD} holds
\des{Q}{HVREL} off when the MCU pin is high impedance. Thus, a high-impedance
MCU pin keeps the display inhibited. \Cref{sec:esp32-usb} gives the same safe
state to the USB data switch.

\paragraph{Gate divider.} $\des{R}{HVINHPU} = \SI{100}{\kilo\ohm}$ and
$\des{R}{HVINHPD} = \SI{1}{\mega\ohm}$ put $0.909\,V_{\mathrm{BUS}}$ on the
gate nominally. The nominal gate voltage spans \SIrange{4.32}{11.5}{\volt}.
Thus, the required
gate rating is $V_{GS} \geq \pm\SI{12}{\volt}$. The pair draws a maximum of
\SI{11.5}{\micro\ampere}.

\paragraph{Ordering.} The earliest UVLO release includes both divider tolerances
and the minimum rising threshold. It occurs at
$V_{\mathrm{BUS}} = \SI{4.16}{\volt}$. The minimum inhibit-gate divider ratio is
0.900. The gate voltage is then \SI{3.74}{\volt}. Therefore, the shunt requires
$V_{GS(th)} \leq \SI{3.74}{\volt}$. The specified
$V_{GS(th)} \leq \SI{2}{\volt}$ meets this limit. It asserts the inhibit by
$V_{\mathrm{BUS}} = \SI{2.22}{\volt}$. The inhibit therefore becomes active at least
\SI{1.94}{\volt} before UVLO can
release. EN can become true first. Until inhibit assertion, the controller stays
in standby and does not switch (\cref{sec:lm5156h}).

\paragraph{Shunt conduction.} At maximum input, \des{Q}{HVINH} sinks
\SI{199}{\micro\ampere}. It must hold UVLO/SYNC below \SI{1.45}{\volt}. Thus,
its channel resistance must be less than \SI{7.3}{\kilo\ohm}. A part with
single-digit ohm resistance at the \SI{3.74}{\volt} minimum required drive meets this
requirement.

\paragraph{Release switch.} \des{Q}{HVREL} blocks a maximum of
\SI{11.5}{\volt} when off. It sinks a maximum of \SI{126}{\micro\ampere} when
on. The MCU supplies a \SI{3.3}{\volt} gate drive. Therefore,
$V_{GS(th)} \leq \SI{2}{\volt}$. One part can serve both switch positions. A
maximum channel resistance of \SI{4.8}{\kilo\ohm} keeps \des{Q}{HVINH} off.
$\des{R}{HVRELPD} = \SI{100}{\kilo\ohm}$ holds \des{Q}{HVREL} off. The MCU
sinks \SI{33}{\micro\ampere} when it turns \des{Q}{HVREL} on.

\subsubsection{Controller housekeeping (BIAS, VCC, PGOOD)}\label{sec:hv-house}

\paragraph{BIAS.}\label{sec:hv-house-bias}
BIAS supplies the internal regulator for the \SI{6.85}{\volt} VCC gate-driver
rail. The datasheet specifies the VCC output at
$V_{\mathrm{BIAS}} = \SI{8}{\volt}$ (\cref{sec:lm5156h}). A direct connection
to $V_{\mathrm{BUS}}$ cannot maintain this condition below approximately
\SI{8}{\volt}. Therefore, the charge pump in \cref{sec:hv-house-pump} supplies
BIAS. A diode-OR connection from $V_{\mathrm{BUS}}$ supplies start-up current.

\paragraph{VCC regulator capacitor.}\label{sec:hv-house-vcc} The datasheet
requires \SIrange{1}{4.7}{\micro\farad} (\ds{lm5156h.pdf}{14}). The design uses
$\des{C}{HVVCC} = \SI{2.2}{\micro\farad}$. This value allows for the DC-bias
derating of a small ceramic capacitor.

\paragraph{PGOOD pull-up.}\label{sec:hv-house-pgood} PGOOD is an open-drain
output (\cref{sec:lm5156h}). The datasheet recommends a pull-up of at least
\SI{10}{\kilo\ohm}. The pull-up must also charge the pin before the
\SI{25}{\micro\second} deglitch interval ends. The design uses
$\des{R}{HVPG} = \SI{100}{\kilo\ohm}$ to the 3.3~V rail. With
\SI{20}{\pico\farad} of stray capacitance, the time constant is approximately
\SI{2}{\micro\second}. The pull-up voltage meets the PGOOD voltage limits. The
low-state current is \SI{33}{\micro\ampere}, below the
\SI{1}{\milli\ampere} limit.

\paragraph{Spread spectrum (DITHOFF).}\label{sec:hv-house-dithoff}
DITHOFF controls the spread-spectrum function (\ds{lm5156h.pdf}{3}). The
modulation span is $0.156 f_{\mathrm{sw}}$, or approximately
\SI{23}{\kilo\hertz}. It combines a triangular sweep with cycle-to-cycle random
modulation (\ds{lm5156h.pdf}{17}).

The dither decreases the DCM idle time. At the top of the dither range,
$t_{\mathrm{on}}+t_{\mathrm{dem}}$ decreases by $3.7\,\%$. The period decreases
by $7.2\,\%$. Thus, the minimum DCM idle fraction decreases from $10.5\,\%$ to
$7.0\,\%$ (\cref{sec:hv-dcm}). The design does not need spread-spectrum
operation. DITHOFF connects to VCC to disable this function.

\subsubsection{Auxiliary bias charge pump}\label{sec:hv-house-pump}

\paragraph{Requirement.} The gate drive above the Miller plateau turns
\des{Q}{HVSW} off. The plateau and \des{D}{HVGOFF} drop subtract
\SI{3.5}{\volt} from VCC. The remaining voltage drives
$\des{R}{HVGOFF}+R_{\mathrm{int}} = \SI{4.8}{\ohm}$. Thus,
$t_{\mathrm{ovl}} = Q_{GD}/I_{\mathrm{gate}}$. Turn-off overlap is the largest
loss term. The commanded peak current must remain below the
\SI{2.36}{\ampere} minimum current limit (\cref{sec:hv-cs-sense}). Solving the
loss iteration limits overlap to \SI{118}{\nano\second}. This condition
requires at least \SI{4.52}{\volt} of gate drive. A direct BIAS feed from
$V_{\mathrm{BUS}}$ gives only \SI{4.20}{\volt} at minimum input. It cannot meet
the requirement.

Therefore, the switch node supplies BIAS through a charge pump. During the off
time, the node is one tenth of the output voltage above $V_{\mathrm{BUS}}$
(\cref{sec:hv-op-converge}). This level is approximately \SI{24}{\volt} at
minimum input. The pump keeps BIAS above \SI{8}{\volt}. At this voltage, the
datasheet specifies the regulated \SI{6.85}{\volt} VCC output.

\paragraph{Topology.} \des{C}{HVPUMP} and \des{R}{HVPUMP} couple the switch
drain to the pump node. \des{D}{HVPUMP} transfers charge to BIAS.
\des{C}{HVBIAS} stores the charge. \des{D}{HVRESET} returns the pump node to
ground during each on time. \des{D}{HVSTART} supplies BIAS from
$V_{\mathrm{BUS}}$ before switching starts.

During the on time, the drain is near ground. \des{D}{HVRESET} holds the pump
node one diode drop below the drain. This condition charges \des{C}{HVPUMP} to
approximately \SI{0.9}{\volt}. At turn-off, the drain rises to
$V_{\mathrm{BUS}} + V_{\mathrm{refl}}$. The pump node rises with it.
\des{D}{HVPUMP} conducts until the node falls to $V_{\mathrm{BIAS}} + V_F$.

The pump needs no regulation loop, because the controller fixes the charge it
must deliver:
\begin{align*}
  Q_{\mathrm{BIAS}} &= Q_G + \frac{I_Q}{f_{\mathrm{sw}}}
  = 40 + \frac{580\,\si{\micro\ampere}}{149.4\,\si{\kilo\hertz}}
  = \SI{43.9}{\nano\coulomb}
  \quad (I_{\mathrm{BIAS}} = \SI{6.56}{\milli\ampere}),
\end{align*}
so BIAS settles wherever the delivered packet equals that demand:
\begin{align*}
  V_{\mathrm{BIAS}} &= (V_{\mathrm{BUS}} + V_{\mathrm{refl}})
  - \SI{0.9}{\volt} - V_{F} - \frac{Q_{\mathrm{BIAS}}}{k\,\des{C}{HVPUMP}},
  &
  k &= 1 - e^{-t_{\mathrm{dem}}/(\des{R}{HVPUMP}\des{C}{HVPUMP})},
\end{align*}
with $k$ the fraction of the packet transferred inside the demagnetization
window. \des{D}{HVSTART} reverse-biases once the pump lifts BIAS past
$V_{\mathrm{BUS}}$, and must be a Schottky: at \SI{4.75}{\volt} a silicon drop
would leave only \SI{0.55}{\volt} above the controller's \SI{3.5}{\volt} BIAS
minimum during start-up.

\paragraph{Sizing.} \des{R}{HVPUMP}'s time constant must be long against the
\SI{10}{\nano\second} leakage spike (\cref{sec:hv-clamp}) and short against
$t_{\mathrm{dem}}$, so the packet completes inside the shelf. $\des{R}{HVPUMP} =
\SI{82}{\ohm}$ gives $\tau = \SI{385}{\nano\second}$, between that spike and a
$\ge\SI{940}{\nano\second}$ window, and limits the current drawn from the switch
node to \SI{118}{\milli\ampere}.

\des{C}{HVPUMP} is set by the worst delivery corner, the bottom of the input
window at the lowest trim setpoint (\SI{164.8}{\volt},
\cref{tab:hv-ballast-divider}), where the reflected shelf and the load are both
smallest, shortening $t_{\mathrm{dem}}$ to \SI{940}{\nano\second} and $k$ to
\num{0.92}. BIAS must stay at or above \SI{8}{\volt}, the condition at which the
datasheet specifies the \SI{6.85}{\volt} VCC output, and $\des{C}{HVPUMP} =
\SI{4.7}{\nano\farad}$ holds \SI{9.09}{\volt} there even at the $-5\,\%$ limit
of a C0G tolerance. The dielectric is specified because the pump meters charge:
a class-II part's voltage coefficient would change the delivered packet with
operating point.

Across the rest of the window BIAS rises with the input, reaching
\SI{20.4}{\volt} at \SI{12.6}{\volt} in, well inside the pin's \SI{60}{\volt}
recommended range. \des{D}{HVZ}, a \SI{30}{\volt} zener, bounds a fault that
drives the reflected voltage above that range and conducts no current in normal
operation. \des{Q}{HVSW} bounds the fault in turn: the drain reaches its
\SI{100}{\volt} rating at roughly \SI{520}{\volt} out, where the shelf is
\SI{65}{\volt} and the pump meters $k\,\des{C}{HVPUMP}(V_{\mathrm{shelf}} -
\SI{1.6}{\volt} - V_Z) = \SI{149}{\nano\coulomb}$ against the
\SI{43.9}{\nano\coulomb} demand. The zener sinks the \SI{105}{\nano\coulomb}
surplus, $30 \times 105\,\si{\nano\coulomb} \times 149.4\,\si{\kilo\hertz} =
\SI{471}{\milli\watt}$, and a \SI{0.5}{\watt} part covers it as a transient:
charging $C_{\mathrm{out}}$ to \SI{520}{\volt} at the current limit takes about
\SI{140}{\milli\second}.

\des{C}{HVBIAS} carries the controller between packets, at a droop of
$I_{\mathrm{BIAS}}/(\des{C}{HVBIAS} f_{\mathrm{sw}}) = \SI{88}{\milli\volt}$ for
\SI{0.5}{\micro\farad}, negligible against the regulator's \SI{2.2}{\volt} of
headroom at the worst corner, so \SI{0.5}{\micro\farad} \emph{effective} is the
floor. A class-II ceramic at \SI{20.4}{\volt} of bias retains well over half its
value on a \SI{50}{\volt} rating, so \SI{1}{\micro\farad} X7R keeps that floor.
The pump stops during pulse skipping at light load. BIAS then draws current
through \des{D}{HVSTART}. The reduced gate voltage increases turn-off loss, but
the peak current is low.

The pump diodes hold the pump node between $-V_F$ and
$V_{\mathrm{BIAS}}+V_F$. Each diode blocks at most \SI{21.1}{\volt} during
normal operation. \des{D}{HVZ} limits the fault stress to
\SI{30.7}{\volt}. With \SI{20}{\percent} derating, require a
\SI{40}{\volt} rating. Each diode must recover before the next switching edge.

\paragraph{Cost.} Charge enters at the shelf potential and leaves at the
\SI{6.85}{\volt} rail. Thus, the loss is
$Q_{\mathrm{BIAS}}f_{\mathrm{sw}}(V_{\mathrm{shelf}}-\SI{6.85}{\volt})$.
Relative to a direct BIAS feed from $V_{\mathrm{BUS}}$, the increment is:
\begin{align*}
  I_{\mathrm{BIAS}} V_{\mathrm{refl}}
  &= \SI{6.56}{\milli\ampere} \times \SI{19.25}{\volt}
  = \SI{126}{\milli\watt},
\end{align*}
constant across the input window because the shelf tracks $V_{\mathrm{BUS}}$.
\des{C}{HVPUMP} redistributes that loss without changing it: a larger value
settles BIAS higher, moving dissipation out of \des{R}{HVPUMP} and into the
controller's internal regulator. At \SI{12.6}{\volt} in, the split is
\SI{89}{\milli\watt} inside the package and the remainder in \des{R}{HVPUMP}.
Gate charge reduces the overhead because the pump supplies it from the shelf
potential. \Cref{sec:hv-summary} includes this effect.

\subsection{Output voltage}\label{sec:hv-setpoint}

\subsubsection{Output setpoint}\label{sec:hv-ignition}

The IN-12B ignition voltage is at most \SI{170}{\volt}
(\cref{sec:in12b}). This value sets the minimum converter output. The INS-1
limit of \SI{90}{\volt} does not constrain the design. A
\SI{15}{\volt} margin gives an adjustment range of
\SIrange{170}{185}{\volt}. One fixed digit ballast cannot keep the current in
range across this full voltage span (\cref{sec:hv-load-digit}). Therefore,
\cref{sec:hv-divider-split} specifies matched ballast and \des{R}{HVD3} values.
The INS-1 has no rated current range. One ballast value serves all output
settings.

\subsubsection{Interpretation of the datasheet's 200/250~V anode-supply figures}\label{sec:hv-anode-figures}

The IN-12B datasheet specifies a minimum anode supply of \SI{200}{\volt}. It
recommends \SI{250}{\volt} DC. Both values exceed the derived range. These
values can apply to unfiltered half-wave AC (inferred). In that mode, the tube
ignites during each mains cycle and conducts near the voltage peak. The supply
then needs more peak voltage to give sufficient average current.

\subsubsection{Operating points used throughout this section}\label{sec:hv-points}

Two operating points recur, bracketing the trim range:
\begin{description}
  \item[Nominal] $V_{\mathrm{BUS}} = \SI{9}{\volt}$, the intermediate USB-PD
  contract; output at \SI{170}{\volt}, the IN-12B's guaranteed ignition ceiling
  (\cref{sec:hv-ignition}); tube currents at their expected values.
  \item[Design-max envelope] $V_{\mathrm{BUS}} = \SI{4.75}{\volt}$ (bottom of
    the input window, worst case for duty cycle and primary conduction), output
    at \SI{191.5}{\volt}, and every digit tube simultaneously at its
    datasheet-maximum \SI{3.5}{\milli\ampere}, which is what the widest-ballast
    row of \cref{tab:hv-ballast-divider} reaches at exactly that current.
\end{description}

In DCM, the energy packet sets peak current:
$I_{\mathrm{pk}} = \sqrt{2P_{\mathrm{in}}/(L_p f_{\mathrm{sw}})}$. Checks that
depend on input voltage use both limits. The \SI{4.75}{\volt} limit maximizes
duty cycle, conduction loss, and on time. The \SI{12.6}{\volt} limit maximizes
commutation voltage, reverse stress, and sense-current slew rate.

\paragraph{Achieved tuning range.}\label{sec:hv-achieved-range}
Select the digit ballast and feedback resistor together from
\cref{tab:hv-ballast-divider}. Each row maps the rated
\SIrange{2.5}{3.5}{\milli\ampere} current range to a different voltage range.
Thus, $V_{\mathrm{HV}}$ can be set from \SI{164.8}{\volt} to
\SI{191.5}{\volt}. Component stress checks use this full range. Functional
checks use the required \SIrange{170}{185}{\volt} range.

\subsection{Ballast sizing}\label{sec:hv-load}

\subsubsection{Digit tubes (NX1 to NX4, ballast \des{R}{BALDIG})}\label{sec:hv-load-digit}

Each IN-12B anode connects to $V_{\mathrm{HV}}$ through its ballast. The
selected cathode returns through a K155ID1 output. The rated DC cathode current
is $I_K \in \SIrange{2.5}{3.5}{\milli\ampere}$ (\cref{sec:in12b}). The ballast
equation is
\begin{align*}
  R_A &= \frac{V_{\mathrm{rail}} - V_M - V_K}{I_K},
\end{align*}
with maintaining voltage $V_M = \SI{145}{\volt}$ (typical, no tabulated spread;
\cref{sec:in12b}) and $V_K$ the K155ID1's on-state drop at the selected output.
The driver specifies a maximum $V_K$ of \SI{2.5}{\volt} at
\SI{7}{\milli\ampere} (\cref{sec:bcd-driver}). The operating current is less
than half this test current. Therefore, the analysis uses an expected
\SI{1.0}{\volt} drop (inferred). At the smallest ballast, a
\SI{1.5}{\volt} error changes current by approximately
\SI{0.20}{\milli\ampere}. This error is less than the
\SI{0.5}{\milli\ampere} half-width of the rated range.

One ballast value cannot keep the cathode current in its rated range across all
output settings. \Cref{sec:hv-divider-split} therefore gives matched ballast
and \des{R}{HVD3} values. For each row, $V_{\mathrm{min}}$ and
$V_{\mathrm{max}}$ give exactly \SI{2.5}{\milli\ampere} and
\SI{3.5}{\milli\ampere}. The trimmer end stops extend past both targets with
worst-case component tolerances.

The installed \des{R}{BALDIG} value is \SI{7.50}{\kilo\ohm}. The table also
defines permissible alternative values through \SI{13.00}{\kilo\ohm}. Each
alternative requires the matching \des{R}{HVD3}. Calculations that depend on the
design envelope cover the full table. Worst-case dissipation grows with $R_A$ at the shared
\SI{3.5}{\milli\ampere} rated maximum, so the largest entry bounds every row:
\begin{align*}
  I_K^2 R_A = 3.5^2 \times 13.00 \approx \SI{159}{\milli\watt},
\end{align*}
still well inside a DIN0207 axial part's $\approx\SI{0.6}{\watt}$ rating
($\approx$27\,\%).

The IN-12B decimal-point cathode has no driver signal and carries no current.

\subsubsection{Colon tubes (NX5 and NX6, ballast \des{R}{BALCOL})}\label{sec:hv-load-colon}

Each INS-1 anode connects to $V_{\mathrm{HV}}$ through its ballast. Its cathode
connects directly to ground. The colon tubes have no driver. They illuminate
when the high-voltage rail is on. The ballast equation has no driver drop. One
resistor value serves both tubes.

The INS-1 maintaining voltage is at most \SI{55}{\volt}
(\cref{sec:ins1}). Use of this limit gives the minimum calculated current. The
datasheet gives no operating-current range. It gives only a
\SI{0.5}{\milli\ampere} optimum brightness current. The ballast gives this
current at the \SI{177.5}{\volt} center of the output range:
\begin{align*}
  R_A &= \frac{\SI{177.5}{\volt} - \SI{55}{\volt}}{\SI{0.5}{\milli\ampere}}
       = \SI{245}{\kilo\ohm}.
\end{align*}
The nearest standard value is $\des{R}{BALCOL} = \SI{240}{\kilo\ohm}$. Current
spans \SIrange{0.479}{0.542}{\milli\ampere} across the required output range.
It spans \SIrange{0.458}{0.569}{\milli\ampere} across the full adjustment
range. These limits are $8.5\,\%$ below and $14\,\%$ above the optimum.

The datasheet does not give a minimum maintaining voltage. A lower maintaining
voltage increases current through a fixed ballast. For example, a
\SI{45}{\volt} maintaining voltage gives \SI{0.61}{\milli\ampere} at maximum
output. The datasheet gives no maximum current for comparison. Therefore, this
upper current is open.

\paragraph{Ballast tolerance.} \des{R}{BALCOL} has a $5\,\%$ tolerance. This
tolerance gives a current range of \SIrange{0.436}{0.599}{\milli\ampere}. The
datasheet gives no rated range that this variation can exceed. Two resistors at
opposite tolerance limits differ in current by \SI{10.5}{\percent}. A $1\,\%$
pair would differ by $2\,\%$. The datasheet does not relate this current
difference to brightness. A \SI{5}{\volt} difference in maintaining voltage
also causes a $4\,\%$ current difference at midrange. The visible colon
ballasts use carbon-film bodies as a construction preference.

\subsection{Feedback divider design}\label{sec:hv-divider}

\subsubsection{Divider current selection}\label{sec:hv-divider-current}

Two limits set the divider current.

\textbf{Power loss.} The divider loads $V_{\mathrm{HV}}$ whenever the converter
operates. A divider current that is 100 times less than the tube current has a
small power cost.

\textbf{Feedback-node impedance.} A high resistance increases error from FB
bias or leakage current. It also increases coupling from switching noise. The
datasheet gives no FB leakage value. The analysis assumes nanoampere-scale
leakage for this device class (inferred). A \SI{100}{\micro\ampere} divider
current makes this leakage negligible.

At the \SI{170}{\volt} reference point that sets a target total divider
resistance of $(170-1)\,\si{\volt}/\SI{100}{\micro\ampere} \approx
\SI{1.69}{\mega\ohm}$.

\subsubsection{Resistor split and trim range}\label{sec:hv-divider-split}

The divider is $R_{\mathrm{upper}}$ to FB to $R_{\mathrm{lower}}$ to ground.
Its equation is $V_{\mathrm{out}} =
V_{\mathrm{REF}}(1+R_{\mathrm{upper}}/R_{\mathrm{lower}})$. At
\SI{170}{\volt}, the required ratio is 169. The lower resistance is less than
$1\,\%$ of the total. It sets the adjustment range but has little effect on
divider current. Two equal upper resistors share the voltage stress
(\cref{sec:hv-divider-stress}). Therefore, $\des{R}{HVD1} =
\des{R}{HVD2} = \SI{845}{\kilo\ohm}$.

$R_{\mathrm{lower}}$ contains fixed resistor \des{R}{HVD3} and trimmer
\des{RV}{HVTRIM} in series. A \SI{2}{\kilo\ohm} trimmer reaches each row's
$V_{\min}$ and $V_{\max}$ with worst-case divider tolerances. Select
\des{R}{HVD3} so that all three targets remain inside the usable trimmer span.
Adjust the output by measurement.

\begin{table}[htbp]
\centering
\caption{Digit-ballast / \des{R}{HVD3} pairing table (component tolerances: \des{R}{HVD1}/\des{R}{HVD2}
    $\pm1\,\%$ each, \des{R}{HVD3} $\pm1\,\%$, $V_{\mathrm{REF}}$ $\pm1\,\%$,
    \des{RV}{HVTRIM} total resistance assumed $\pm10\,\%$, all stacked
    worst-case additive)}
\label{tab:hv-ballast-divider}
\begin{tabular}{|c|c|c|c|c|c|c|}
\hline
\textbf{$R_A$} & \textbf{$V_{\min}$} & \textbf{$V_{\mathrm{typ}}$} &
\textbf{$V_{\max}$} & \textbf{\des{R}{HVD3}} & \textbf{$V_{\mathrm{floor,wc}}$} &

\textbf{$V_{\mathrm{ceil,wc}}$} \\
(\si{\kilo\ohm}) & (\si{\volt}) & (\si{\volt}) & (\si{\volt}) &
(\si{\kilo\ohm}) & (\si{\volt}) & (\si{\volt}) \\
\hline
7.50  & 164.8 & 168.5 & 172.2 & 9.09 & 146.5 & 192.6 \\
\hline
8.20  & 166.5 & 170.6 & 174.7 & 8.87 & 149.4 & 197.3 \\
\hline
9.10  & 168.8 & 173.3 & 177.8 & 8.87 & 149.4 & 197.3 \\
\hline
10.00 & 171.0 & 176.0 & 181.0 & 8.66 & 152.3 & 202.1 \\
\hline
11.00 & 173.5 & 179.0 & 184.5 & 8.45 & 155.3 & 207.1 \\
\hline
12.00 & 176.0 & 182.0 & 188.0 & 8.25 & 158.3 & 212.1 \\
\hline
13.00 & 178.5 & 185.0 & 191.5 & 8.25 & 158.3 & 212.1 \\
\hline
\end{tabular}
\end{table}

$V_{\min}$, $V_{\mathrm{typ}}$, and $V_{\mathrm{max}}$ are adjustment targets.
They give \SIlist{2.5;3.0;3.5}{\milli\ampere} through the applicable ballast.
Measure the output and adjust it to a target. Divider tolerance does not change
these measured targets. $V_{\mathrm{floor,wc}}$ and
$V_{\mathrm{ceil,wc}}$ are calculated end-stop limits. Start with the
smallest-$R_A$ row. Use a larger ballast only if the installed tubes do not
ignite in that row's range.

\paragraph{User-guide requirement.} The hardware user guide must include
\cref{tab:hv-ballast-divider}. It must identify each row as an acceptable
ballast and \des{R}{HVD3} pair. It must require voltage adjustment with the
tubes removed. The end-stop values are informational. Adjustment to an end stop
is not permitted. The guide must also state that the trimmer range compensates
for component tolerance. The end stops are excluded from operating and stress
analyses.

Divider current at each row's nominal point ranges
\SIrange{99.1}{108.9}{\micro\ampere}, matching the \SI{100}{\micro\ampere}
target of \cref{sec:hv-divider-current} across the whole table.

\subsubsection{Divider element stress}\label{sec:hv-divider-stress}

\des{R}{HVD1} and \des{R}{HVD2} split $(V_{\mathrm{out}}-1)$ evenly:
\SI{84.5}{\volt} each at the \SI{170}{\volt} reference point, \SI{95.25}{\volt}
at the \SI{191.5}{\volt} design-max envelope (\cref{sec:hv-achieved-range}),
against the \SI{250}{\volt} limiting element voltage of the specified DIN0207
axial metal film parts. Dissipation is $95.25^2/\SI{845}{\kilo\ohm} =
\SI{10.7}{\milli\watt}$. \des{R}{HVD3} and \des{RV}{HVTRIM} see about
\SI{1}{\volt}.

\subsection{Aggregate current and power budget}\label{sec:hv-budget}

The bleeder \des{R}{HVBLEED} completes the load list at \SI{1.5}{\mega\ohm},
sized in \cref{sec:hv-bleeder} against the output-discharge and minimum-load
requirements.

{\footnotesize
\begin{table}[htbp]
\centering
\caption{$V_{\mathrm{HV}}$: aggregate anode current budget}\label{tab:hv-current-budget}
\begin{tabular}{|l|r|r|}
\hline
\textbf{Load element} & \textbf{Nominal (\SI{170}{\volt})} & \textbf{Design-max (\SI{191.5}{\volt})} \\
\hline
NX1 to NX4 digit cathodes & $4 \times 3.20 = \SI{12.80}{\milli\ampere}$
  & $4 \times 3.50 = \SI{14.00}{\milli\ampere}$ \\
\hline
NX5 and NX6 colon anodes & $2 \times 0.479 = \SI{0.96}{\milli\ampere}$
  & $2 \times 0.569 = \SI{1.14}{\milli\ampere}$ \\
\hline
Bleeder \des{R}{HVBLEED} (\SI{1.5}{\mega\ohm}) & \SI{0.11}{\milli\ampere} & \SI{0.13}{\milli\ampere} \\

\hline
Feedback divider ($\approx\SI{1.70}{\mega\ohm}$) & \SI{0.10}{\milli\ampere} & \SI{0.11}{\milli\ampere} \\
\hline
Total $I_{\mathrm{out}}$ & \SI{13.97}{\milli\ampere} & \SI{15.38}{\milli\ampere} \\
\hline
$P_{\mathrm{out}} = V_{\mathrm{out}} I_{\mathrm{out}}$ & \SI{2.38}{\watt} & \SI{2.94}{\watt} \\
\hline
\end{tabular}
\end{table}
}

The nominal-point digit and colon currents follow from the specified ballast
values (\cref{sec:hv-load-digit,sec:hv-load-colon}) at \SI{170}{\volt}; the
divider's \SI{0.10}{\milli\ampere} is under 1\,\% of total output current at
either operating point, as \cref{sec:hv-divider-current} assumed.

\textbf{The maximum-design point can occur.} With a
\SI{13.00}{\kilo\ohm} ballast and a \SI{191.5}{\volt} output, each digit tube
draws \SI{3.5}{\milli\ampere}. Thus, all four tubes can operate at their rated
maximum current at the same time. The INS-1 has no rated maximum current
(\cref{sec:ins1}).

\subsection{Operating point and efficiency}\label{sec:hv-op}

\subsubsection{Component values for the loss model}\label{sec:hv-op-inputs}

The loss model uses the following values. Their cross-referenced sections
derive them from the converged peak current.
\begin{itemize}[nosep]
  \item \des{Q}{HVSW} is an IRL540N-type commodity through-hole logic-level power MOSFET;
  \item $\des{R}{HVGON} = \SI{10}{\ohm}$, governing turn-on;
  \item \des{D}{HVGOFF} (SB240S-E3/54) and $\des{R}{HVGOFF} = \SI{3.3}{\ohm}$,
    the turn-off speedup network, sized in
    \cref{sec:hv-op-losses,sec:hv-switch};
  \item $\des{R}{HVCS} = \SI{39}{\milli\ohm}$, the current-sense resistor,
    derived in \cref{sec:hv-cs};
  \item $\des{R}{HVCL} = \SI{30}{\kilo\ohm}$, $\des{C}{HVCL} =
    \SI{12}{\nano\farad}$, the RCD clamp network, sized in
    \cref{sec:hv-clamp}.
\end{itemize}

\subsubsection{Loss mechanisms}\label{sec:hv-op-losses}

In DCM, the converter transfers a fixed magnetizing-energy packet in each
cycle:
\begin{align*}
  P_{\mathrm{in}} = \tfrac{1}{2} L_p I_{\mathrm{pk}}^2 f_{\mathrm{sw}}
  \quad\Longrightarrow\quad
  I_{\mathrm{pk}} = \sqrt{\frac{2 P_{\mathrm{in}}}{L_p f_{\mathrm{sw}}}},
  \qquad
  P_{\mathrm{in}} = \frac{P_{\mathrm{out}}}{\eta}.
\end{align*}
Calculate efficiency by iteration. Start with an efficiency value. Calculate
$P_{\mathrm{in}}$ and $I_{\mathrm{pk}}$. Then calculate each loss and update
the efficiency. Repeat these calculations until they converge. The loss models
are as follows.
\begin{description}[style=nextline]
  \item[Primary conduction] Triangular current over $t_{\mathrm{on}}$:
    $I_{\mathrm{rms}} = I_{\mathrm{pk}}\sqrt{D/3}$ through
    $R_{DS(\mathrm{on})} + \des{R}{HVCS} + R_{\mathrm{DCR,pri}}
     = 55 + 39 + 13 = \SI{107}{\milli\ohm}$. The datasheet tabulates
    $R_{DS(\mathrm{on})}$ at only two $V_{GS}$ points, both at
    $T_J=\SI{25}{\celsius}$: \SI{53}{\milli\ohm} at $V_{GS}=\SI{5.0}{\volt}$ and
    \SI{44}{\milli\ohm} at $V_{GS}=\SI{10}{\volt}$
    (\ds{infineon-irl540n-datasheet-en.pdf}{2}), so interpolating to the
    \SI{6.85}{\volt} gate rail gives a cold value of
    $\approx\SI{49.7}{\milli\ohm}$. It then rises with junction temperature
    (\ds{infineon-irl540n-datasheet-en.pdf}{3}, Fig.\,4, read as a linear
    approximation between its $1.0\times$-at-\SI{25}{\celsius} and
    $2.4\times$-at-\SI{175}{\celsius} anchors), which \des{Q}{HVSW}'s own
    dissipation sets. Solved self-consistently against the design-max loss
    budget below (\SI{25}{\celsius} ambient assumed), the converged point is
    $T_J \approx \SI{37}{\celsius}$, $R_{DS(\mathrm{on})} \approx
    \SI{55}{\milli\ohm}$.

  \item[Turn-off overlap]
    The IRL540N datasheet (\ds{infineon-irl540n-datasheet-en.pdf}{2}) specifies
    the gate-drain (Miller) charge only as a guaranteed maximum, $Q_{GD} \le
    \SI{38}{\nano\coulomb}$, with no typical value tabulated; its gate-charge
    curve (\ds{infineon-irl540n-datasheet-en.pdf}{4}, Fig.\,6) shows the Miller
    plateau sitting at $V_{\mathrm{plateau}} \approx \SI{3}{\volt}$. The same
    datasheet's resistive-load switching-time table gives $t_f =
    \SI{62}{\nano\second}$ typical at its own test condition
    ($R_G=\SI{5.0}{\ohm}$, $V_{GS}$ pulse $=\SI{5.0}{\volt}$), which implies a
    typical $Q_{GD} \approx t_f \times (5.0-3.0)/5.0\,\si{\ohm} \approx
    \SI{25}{\nano\coulomb}$ at that test point. The LM51561H's output stage adds
    an internal pull-down resistance in series with $\des{R}{HVGON}$. Its datasheet
    (\ds{lm5156h.pdf}{6}, MOSFET DRIVER) gives the low-state output's voltage
    drop under an externally forced \SI{100}{\milli\ampere} as \SI{0.15}{\volt},
    implying $R_{\mathrm{int}} \approx \SI{1.5}{\ohm}$ for this pull-down path
    (the table's two rows are ambiguous about which pin state pairs with
    which current direction, so this is the best available reading).
    Turn-off overlap dominates the loss budget, so that edge is given a faster
    discharge path of its own. \des{D}{HVGOFF} in series with
    \des{R}{HVGOFF} bypasses $\des{R}{HVGON}$ for the discharge direction only:
    its anode sits at the gate, its cathode at the
    $\des{R}{HVGON}$/driver-output node, so it blocks during turn-on, leaving
    that edge unchanged, and conducts during turn-off. \des{R}{HVGOFF} is sized
    against the driver's \SI{1.5}{\ampere} peak source/sink capability
    (\ds{lm5156h.pdf}{29}), the binding current being the brief pre-plateau
    discharge from the fully-enhanced gate, which exceeds the plateau current:
    \begin{align*}
      I_{\mathrm{gate}}^{\mathrm{peak}} &= \frac{V_{\mathrm{drive}} -
        V_{F(\des{D}{HVGOFF})}}{\des{R}{HVGOFF} + R_{\mathrm{int}}}
        = \frac{6.85 - 0.5}{3.3+1.5} \approx \SI{1.32}{\ampere}
        \quad (88\,\% \text{ of the } \SI{1.5}{\ampere} \text{ rating}),
    \end{align*}
    using $V_{F(\des{D}{HVGOFF})} \approx \SI{0.5}{\volt}$, interpolated below
    the sole tabulated point ($\SI{0.55}{\volt}$ at \SI{2}{\ampere}). The
    \SI{12}{\percent} headroom rests on that interpolation and on the
    $R_{\mathrm{int}}$ estimate; both are derived, neither measured.
    At the lower Miller-plateau voltage, the turn-off gate current is
    $I_{\mathrm{gate}} = (6.85-3.0-0.5)/(3.3+1.5) \approx \SI{698}{\milli\ampere}$,
    giving
    \begin{align*}
      t_{\mathrm{ovl}}^{\mathrm{typ}} &= \frac{Q_{GD}^{\mathrm{typ}}}{I_{\mathrm{gate}}}
        \approx \SI{36}{\nano\second}, &
      t_{\mathrm{ovl}}^{\max} &= \frac{Q_{GD}^{\max}}{I_{\mathrm{gate}}}
        \approx \SI{54}{\nano\second}.
    \end{align*}
    With the commutation voltage $V_{\mathrm{com}} = V_{\mathrm{BUS}} +
    V_{\mathrm{refl}}$, $E_{\mathrm{off}} \approx \tfrac{1}{2} I_{\mathrm{pk}}
    V_{\mathrm{com}} t_{\mathrm{ovl}}$. The convergence below uses
    $t_{\mathrm{ovl}}^{\mathrm{typ}}$; \cref{sec:hv-switch} checks
    \des{Q}{HVSW}'s dissipation at $t_{\mathrm{ovl}}^{\max}$. Turn-on is
    near-lossless in DCM (zero current; $C_{oss}$ discharge contributes a few
    milliwatts) and is governed by $\des{R}{HVGON}$ alone.
    \item[Gate drive and controller]
    $P = (Q_G f_{\mathrm{sw}} + I_Q)(V_{\mathrm{BUS}} + V_{\mathrm{refl}})$
    with $Q_G \approx \SI{40}{\nano\coulomb}$ at the \SI{6.85}{\volt} VCC rail
    (scaled from the \SI{10}{\volt} datasheet point, inferred) and controller
    operating current $I_Q \le \SI{580}{\micro\ampere}$ (\ds{lm5156h.pdf}{5}). Gate
    charge is drawn through the internal VCC regulator from BIAS, which the
    charge pump of \cref{sec:hv-house-pump} lifts from the switch node, so the
    charge is drawn at the reflected shelf, well above $V_{\mathrm{BUS}}$. It is
    the largest single term in the budget below.

  \item[RCD clamp] Leakage energy $\tfrac{1}{2} L_{\mathrm{lk}} I_{\mathrm{pk}}^2$
    per cycle, dissipated in \des{R}{HVCL} with the clamp-voltage multiplier
    derived in \cref{sec:hv-clamp}.
  \item[Rectifier] $V_F \approx \SI{1.0}{\volt}$ (ultrafast-diode class figure at
    $\sim$\SI{0.2}{\ampere}) times $I_{\mathrm{out}}$; recovery loss is negligible
    because in DCM the secondary current reaches zero before the switch turns on.
  \item[Secondary conduction] $I_{s,\mathrm{rms}}^2 \times \SI{1.60}{\ohm}$ DCR.
  \item[Capacitor ESR] Input $\approx\SI{40}{\milli\ohm}$ and output
  $\approx\SI{0.5}{\ohm}$, both the polymer parts' own; \des{C}{HVI2}'s
  impedance stays above the input figure below \SI{850}{\kilo\hertz}, past the
  ripple's dominant harmonics.
  \item[Core] Uncharacterized at this frequency; a \SI{50}{\milli\watt} allowance
    is carried.
\end{description}

\subsubsection{Converged operating point (design-max envelope)}\label{sec:hv-op-converge}

Iterating at the design-max envelope ($P_{\mathrm{out}} = \SI{2.94}{\watt}$,
$V_{\mathrm{BUS}} = \SI{4.75}{\volt}$, $V_{\mathrm{out}} = \SI{191.5}{\volt}$,
reflected voltage $V_{\mathrm{refl}} = (191.5+1)/10 = \SI{19.25}{\volt}$,
$t_{\mathrm{ovl}} = t_{\mathrm{ovl}}^{\mathrm{typ}} = \SI{36}{\nano\second}$),
the loop converges to $\eta = \SI{83.3}{\percent}$:
\begin{align*}
  P_{\mathrm{in}} &= \frac{2.945}{0.833} = \SI{3.54}{\watt}, &
  E_{\mathrm{cyc}} &= \frac{3.54}{149.4\times10^{3}} = \SI{23.7}{\micro\joule},\\
  I_{\mathrm{pk}} &= \sqrt{\frac{2 \times 23.7\,\si{\micro\joule}}
  {\SI{10}{\micro\henry}}} = \SI{2.18}{\ampere}, &
  t_{\mathrm{on}} &= \frac{L_p I_{\mathrm{pk}}}{V_{\mathrm{BUS}}}
                   = \frac{10\,\si{\micro\henry} \times 2.18}{4.75}
                   = \SI{4.58}{\micro\second},\\
                   D &= t_{\mathrm{on}} f_{\mathrm{sw}} = 0.68, &
                   I_{\mathrm{rms,pri}} &= I_{\mathrm{pk}}\sqrt{D/3} = \SI{1.04}{\ampere},\\
                   I_{\mathrm{in}} &= \frac{P_{\mathrm{in}}}{V_{\mathrm{BUS}}}
                   = \SI{744}{\milli\ampere}.
\end{align*}
Converged loss budget at this point:

{\footnotesize
\begin{table}[htbp]
\centering
\caption{170~V flyback converter: loss budget at the design-max envelope}\label{tab:hv-loss-budget}
\begin{tabular}{|l|r|}
\hline
\textbf{Mechanism} & \textbf{Loss} \\
\hline
Bias draw through the charge pump (\cref{sec:hv-house-pump}) & \SI{157}{\milli\watt} \\

\hline
MOSFET turn-off overlap ($t_{\mathrm{ovl}}=\SI{36}{\nano\second}$) & \SI{140}{\milli\watt} \\

\hline
Primary conduction (FET 60, \des{R}{HVCS} 42, DCR 14) & \SI{116}{\milli\watt} \\

\hline
RCD clamp (\des{R}{HVCL}) & \SI{86}{\milli\watt} \\

\hline
Core allowance (uncharacterized) & \SI{50}{\milli\watt} \\
\hline
Capacitor ESR (input + output) & \SI{22}{\milli\watt} \\

\hline
Rectifier \des{D}{HVRECT} conduction & \SI{15}{\milli\watt} \\

\hline
Secondary DCR & \SI{4}{\milli\watt} \\
\hline
Total & \SI{590}{\milli\watt} \\

\hline
$\eta = 2.945/(2.945+0.590)$ & \SI{83.3}{\percent} \\

\hline
\end{tabular}
\end{table}
}

\paragraph{Efficiency curve.}\label{sec:hv-op-curve}
The charge pump holds the gate rail constant (\cref{sec:hv-house-pump}).
Turn-off overlap and bias loss increase with
$V_{\mathrm{BUS}} + V_{\mathrm{refl}}$. Together, they increase from
\SI{297}{\milli\watt} to \SI{400}{\milli\watt}. Primary conduction and
capacitor ESR decrease with duty cycle. Together, they decrease from
\SI{138}{\milli\watt} to \SI{57}{\milli\watt}. The other modeled losses total
\SI{155}{\milli\watt}. The total loss is
$A(V_{\mathrm{BUS}}+V_{\mathrm{refl}})+B/V_{\mathrm{BUS}}+C$. Here,
$A=\SI{12.4}{\milli\watt\per\volt}$ and
$B=\SI{683}{\milli\watt\volt}$. The minimum occurs at
$\sqrt{B/A}=\SI{7.4}{\volt}$:

{\footnotesize
  \begin{table}[htbp]
    \centering
    \caption{Design-max efficiency across the input window}\label{tab:hv-eta-line}
    \begin{tabular}{|l|c|c|c|c|c|}
      \hline
      $V_{\mathrm{BUS}}$ & \SI{4.75}{\volt} & \SI{7}{\volt} & \SI{9}{\volt} &
      \SI{11}{\volt} & \SI{12.6}{\volt} \\
      \hline
      Duty & 0.68 & 0.46 & 0.36 & 0.30 & 0.26 \\
      \hline
      Total loss & \SI{590}{\milli\watt} & \SI{578}{\milli\watt} &
      \SI{583}{\milli\watt} & \SI{596}{\milli\watt} & \SI{610}{\milli\watt} \\
      \hline
      $\eta$ & \SI{83.3}{\percent} & \SI{83.6}{\percent} & \SI{83.5}{\percent} &
      \SI{83.2}{\percent} & \SI{82.8}{\percent} \\
      \hline
    \end{tabular}
  \end{table}
}

The convergence above uses $t_{\mathrm{ovl}}^{\mathrm{typ}}$; across the
\SIrange{36}{54}{\nano\second} spread the design-max efficiency band runs
\SIrange{81.4}{83.3}{\percent}. \Cref{sec:hv-switch} checks \des{Q}{HVSW}'s
dissipation at $t_{\mathrm{ovl}}^{\max}$, and \cref{sec:hv-cs-sense} folds the
same spread into its current-limit margin as a tolerance term.

The two largest losses depend on the \des{Q}{HVSW} gate charge. Bias loss
depends on $Q_G$. Turn-off overlap depends on $Q_{GD}$. Together, these losses
are \SI{297}{\milli\watt} of the \SI{590}{\milli\watt} total. Both values are
derived. Measure converter efficiency to verify this loss model.

\subsubsection{Nominal-point efficiency}\label{sec:hv-op-nominal}

The same procedure at the nominal point ($\SI{2.38}{\watt}$, \SI{170}{\volt})
gives, at each contract:

{\footnotesize
  \begin{table}[htbp]
    \centering
    \caption{Nominal-point ($\SI{2.38}{\watt}$, \SI{170}{\volt}) operating figures}\label{tab:hv-nominal-line}
    \begin{tabular}{|l|c|c|c|}
      \hline
      $V_{\mathrm{BUS}}$ & \SI{4.75}{\volt} & \SI{9}{\volt} & \SI{12.6}{\volt} \\
      \hline
      $\eta$ & \SI{82.7}{\percent} & \SI{82.5}{\percent} & \SI{81.7}{\percent} \\
      \hline
      $P_{\mathrm{in}}$ & \SI{2.87}{\watt} & \SI{2.88}{\watt} & \SI{2.91}{\watt} \\
      \hline
      $I_{\mathrm{pk}}$ & \SI{1.96}{\ampere} & \SI{1.96}{\ampere} & \SI{1.97}{\ampere} \\
      \hline
      Duty & 0.62 & 0.33 & 0.23 \\
      \hline
      $I_{\mathrm{rms,pri}}$ & \SI{0.89}{\ampere} & \SI{0.65}{\ampere} & \SI{0.55}{\ampere} \\
      \hline
      $I_{\mathrm{in}}$ average & \SI{605}{\milli\ampere} & \SI{320}{\milli\ampere} &
      \SI{231}{\milli\ampere} \\
      \hline
    \end{tabular}
  \end{table}
}

Nominal efficiency is lower than maximum-load efficiency. The main losses are
proportional to $\sqrt{P}$. Bias current does not decrease with load. Therefore,
the \SIrange{143}{195}{\milli\watt} bias loss is a larger fraction of the
\SI{2.38}{\watt} nominal output. The nominal efficiency is
0.6 to 1.1 percentage points lower than the maximum-load efficiency.

The $I_{\mathrm{in}}$ row is what \cref{sec:hv-inhibit}'s current test is against.
A \SI{5}{\volt} contract must carry \SI{605}{\milli\ampere} at the nominal point
and \SI{744}{\milli\ampere} at the design-max envelope, so a
\SI{3}{\ampere} \SI{5}{\volt} contract qualifies and a default
\SI{500}{\milli\ampere} port does not.

\subsubsection{Peak current at minimum inductance}\label{sec:hv-op-lmin}

With the minimum-tolerance inductance ($L_p = \SI{9}{\micro\henry}$ at rated peak
current, \cref{sec:da2032}), the same energy packet requires
\begin{align*}
  I_{\mathrm{pk}}^{(L\min)} &= \sqrt{\frac{2 \times 23.7\,\si{\micro\joule}}
                                          {\SI{9}{\micro\henry}}}
                             = \SI{2.29}{\ampere},
\end{align*}
the worst-case commanded peak used in \cref{sec:hv-cs}.

\subsection{DCM verification}\label{sec:hv-dcm}

DCM requires $t_{\mathrm{on}} + t_{\mathrm{dem}} < T$. The remaining time must
include the drain ring-down. For the 1:10 turns ratio, the demagnetization time
is
\begin{align*}
  t_{\mathrm{dem}} &= \frac{n L_p I_{\mathrm{pk}}}{V_{\mathrm{out}} + V_F}.
\end{align*}
The binding corner maximizes $t_{\mathrm{on}} + t_{\mathrm{dem}}$: maximum
inductance ($\SI{11}{\micro\henry}$, $+10\,\%$), minimum input
(\SI{4.75}{\volt}), design-max load:
\begin{align*}
  I_{\mathrm{pk}} &= \sqrt{\frac{2\times23.7\,\si{\micro\joule}}{\SI{11}{\micro\henry}}}
  = \SI{2.07}{\ampere},\\
  t_{\mathrm{on}} &= \frac{11\,\si{\micro\henry}\times 2.07}{4.75}
  = \SI{4.80}{\micro\second}, &
  t_{\mathrm{dem}} &= \frac{10 \times 11\,\si{\micro\henry} \times 2.07}{192.5}
  = \SI{1.19}{\micro\second},\\
  t_{\mathrm{on}} + t_{\mathrm{dem}} &= \SI{5.99}{\micro\second}
  = 0.89\,T, &
  t_{\mathrm{idle}} &= \SI{0.71}{\micro\second}.
\end{align*}

\textbf{DCM has a minimum idle fraction of $10.5\,\%$.} The secondary circuit
sets demagnetization time, which does not depend on input voltage. On time is
inversely proportional to $V_{\mathrm{BUS}}$. Therefore, minimum input gives
the smallest idle fraction. At \SI{12.6}{\volt}, duty falls to 0.26. A lighter
load also increases the idle fraction.

CCM operation invalidates the energy-packet loss model and the single-pole
plant in \cref{sec:hv-comp}. It also invalidates the rectifier-recovery and
slope-compensation conclusions. The \SI{0.71}{\micro\second} idle time can end
during drain ring-down. This timing adds only a few milliwatts because DCM
turn-on current is zero.

The 0.72 duty cycle does not cause subharmonic oscillation in DCM. The
$D>0.5$ threshold applies to CCM. In DCM, the current returns to zero during
each cycle. Therefore, the design does not require $R_{SL}$
(\cref{sec:hv-cs}).

Duty at the worst corner remains under the controller's $\approx$0.91 maximum
near this frequency, and $t_{\mathrm{on}} \ge \SI{4.8}{\micro\second}$ is well
above the $\approx\SI{165}{\nano\second}$ minimum on-time (\cref{sec:lm5156h}).

\subsection{Current sensing and current limit}\label{sec:hv-cs}

\subsubsection{CS filter}\label{sec:hv-cs-filter}

With the \SI{100}{\milli\volt} (93 to 107) current-limit threshold of
\cref{sec:lm5156h} and \des{R}{HVCS} $= \SI{39}{\milli\ohm}$ (declared in
\cref{sec:hv-op-inputs}):
\begin{align*}
  I_{\mathrm{lim}}^{\mathrm{typ}} &= \frac{\SI{100}{\milli\volt}}{\SI{39}{\milli\ohm}}
    = \SI{2.56}{\ampere},\\
  I_{\mathrm{lim}}^{\min} &= \frac{\SI{93}{\milli\volt}}{39\,\si{\milli\ohm}\times1.01}
    = \SI{2.36}{\ampere}, &
  I_{\mathrm{lim}}^{\max} &= \frac{\SI{107}{\milli\volt}}{39\,\si{\milli\ohm}\times0.99}
    = \SI{2.77}{\ampere}.
\end{align*}
The CS pin connects to the sense node through \des{R}{HVCSF}.
\des{C}{HVCSF} connects the pin to ground. The circuit does not use $R_{SL}$.
Thus, the current-limit threshold does not depend on duty cycle. The minimum
on-time and off-time criteria have large margins. Two other requirements set
$\tau = \des{R}{HVCSF}\des{C}{HVCSF}$.

\textbf{Upper limit: current-limit error.} The filter delays the sampled current
by $\Delta I_{\mathrm{delay}} = (V_{\mathrm{BUS}}/L_p)\,\tau$. Current continues
to increase during this delay. Limit the increase to half of the current-limit
tolerance range, or \SI{0.205}{\ampere}:
\begin{align*}
  \tau < \frac{\SI{0.205}{\ampere}}{(V_{\mathrm{BUS}}/L_p)_{\max}}
       = \frac{\SI{0.205}{\ampere}}{\SI{1.40}{\ampere\per\micro\second}}
       = \SI{146}{\nano\second},
\end{align*}
evaluated at the fastest observed slope (\SI{12.6}{\volt}, \SI{9}{\micro\henry}).

\textbf{Lower limit: spike attenuation.} Gate current and $C_{oss}$ discharge
cause a turn-on current spike. The spike duration is tens of nanoseconds
(inferred). The filter time constant must be comparable to this duration.

$\des{R}{HVCSF} = \SI{100}{\ohm}$ and
$\des{C}{HVCSF} = \SI{470}{\pico\farad}$. Both values are in the TI
recommended ranges. Their time constant is
\begin{align*}
  \tau = \des{R}{HVCSF}\des{C}{HVCSF} = 100 \times 470\,\si{\pico\farad} = \SI{47}{\nano\second},
\end{align*}
This value is comparable to the inferred spike duration. It is less than the
\SI{146}{\nano\second} maximum. The resulting current error is approximately
\SI{66}{\milli\ampere}. A \SI{50}{\nano\second} comparator and driver delay
(inferred) adds \SI{70}{\milli\ampere}. The total overshoot is approximately
\SI{0.14}{\ampere}.
\Cref{sec:hv-cs-sense} uses this value.

\subsubsection{Sense resistor and limit thresholds}\label{sec:hv-cs-sense}

\des{R}{HVCS} carries the switch current to ground; with the current-limit
thresholds of \cref{sec:hv-cs-filter} and the filter's overshoot $\Delta I_{\mathrm{ovs}}
\approx \SI{0.14}{\ampere}$, the sense resistor must satisfy two opposing
constraints:
\begin{align*}
  \text{power delivery:}\quad
    R_{S} &< \frac{\SI{93}{\milli\volt}}{1.01 I_{\mathrm{pk}}^{(L\min)}}
          = \frac{93\,\si{\milli\volt}}{1.01\times2.29\,\si{\ampere}} = \SI{40.2}{\milli\ohm},\\
  \text{core protection:}\quad
    R_{S} &> \frac{\SI{107}{\milli\volt}}{0.99(I_{\mathrm{sat}} - \Delta I_{\mathrm{ovs}})}
          = \frac{107\,\si{\milli\volt}}{0.99(3.0 - 0.14)\,\si{\ampere}} = \SI{37.8}{\milli\ohm},
\end{align*}
a \SI{2.4}{\milli\ohm} nominal window; neither bound depends on
$V_{\mathrm{out}}$. \SI{39}{\milli\ohm} (declared in \cref{sec:hv-op-inputs})
sits inside it, \SI{1.2}{\milli\ohm} under the power-delivery ceiling and
\SI{1.2}{\milli\ohm} above the core-protection floor. Both bounds are
worst-case-additive: minimum datasheet threshold, \des{R}{HVCS} at its own
tolerance extreme, and, on the power-delivery side, $L_p$ at its $-10\,\%$
extreme, all stacked simultaneously in the same direction.

\SI{39}{\milli\ohm} is the only standard value in this range.

\textbf{Tolerance combination.} An insufficient margin causes the current
limit to end a switching cycle early. The result is reduced output energy and
display brightness. It does not cause component damage. Therefore, this check
also uses a root-sum-square (RSS) tolerance combination. Turn-off overlap is an
independent part-to-part variation. Include it in the RSS calculation.

The nominal margin is $M_{\mathrm{nom}}=\SI{0.389}{\ampere}$. It uses
$\SI{100}{\milli\volt}/\SI{39}{\milli\ohm}$ and nominal $L_p$. It also uses
the typical turn-off overlap. The individual changes are:
\begin{align*}
  \Delta_{\mathrm{threshold}} &= \SI{-0.180}{\ampere}, &
  \Delta_{R_S} &= \SI{-0.025}{\ampere}, &
  \Delta_{L_p} &= \SI{-0.118}{\ampere}, &
  \Delta_{t_{\mathrm{ovl}}} &= \SI{-0.026}{\ampere}.
\end{align*}
Worst-case-additive sums these directly: $M_{\mathrm{wc}} = 0.389 - 0.180 - 0.025
- 0.118 - 0.026 = \SI{+0.040}{\ampere}$ ($+1.7\,\%$ of $I_{\mathrm{lim,min}}$).
RSS combines the same four in quadrature:
\begin{align*}
  M_{\mathrm{RSS}} &= M_{\mathrm{nom}} - \sqrt{\Delta_{\mathrm{threshold}}^2 +
    \Delta_{R_S}^2 + \Delta_{L_p}^2 + \Delta_{t_{\mathrm{ovl}}}^2}
    = 0.389 - 0.218 = \SI{+0.171}{\ampere}
    \quad (+7.2\,\% \text{ of } I_{\mathrm{lim,min}}),
\end{align*}
The threshold specification dominates the combination. \textbf{The RSS treatment
is scoped to this one check}; every voltage- and thermal-rating margin elsewhere
remains worst-case-additive, because those checks guard against destructive
outcomes.

Both tolerance methods give a positive margin. The additive margin is
$1.7\,\%$. The RSS margin is $7.2\,\%$. The bias current decreases this margin
because it increases input power and commanded peak current.

\paragraph{What hiccup changes about that.}\label{sec:hv-cs-hiccup}
On the LM51561H, sixty-four \emph{consecutive} limited cycles trip the fault
timer (\cref{sec:hv-ss}), after which the converter stops for
\SI{219}{\milli\second} and retries. A board that exhausts this margin at the
top of its trim range on a \SI{5}{\volt} contract blinks at roughly
\SI{4.6}{\hertz}.

The fault requires 64 consecutive limited cycles. Eight unlimited cycles reset
the counter. Thus, intermittent limiting does not cause hiccup. Continuous
limiting requires the maximum load and output setting. A lower output setting
or a larger ballast decreases peak current.

\textbf{The hardware user guide must include this condition.} If the display
blinks at maximum adjustment on a \SI{5}{\volt} supply, decrease the output.
Alternatively, install a larger ballast. Measure current limit with the
selected transformer to verify the calculated margin.

On the other side, the maximum clamped peak plus overshoot reaches
$2.77 + 0.14 = \SI{2.91}{\ampere}$ against the \SI{3.0}{\ampere} transformer
rating: a 3\,\% margin, acceptable because the \SI{3.0}{\ampere} figure is a
soft-saturation point, with $L$ already specified at \SI{9}{\micro\henry}
there.

\des{R}{HVCS} dissipation is $I_{\mathrm{rms}}^2 R_S = 1.04^2 \times 0.039 =
\SI{42}{\milli\watt}$ at the design-max envelope, trivial in a 2512.

\subsection{RCD clamp}\label{sec:hv-clamp}

At turn-off, the leakage inductance carries $I_{\mathrm{pk}}$. This inductance
is not coupled to the secondary. \des{D}{HVCL} directs its current into
\des{C}{HVCL}. \des{R}{HVCL} discharges the capacitor between cycles. The
network returns to $V_{\mathrm{BUS}}$. Thus, \des{C}{HVCL} withstands only
$V_{\mathrm{cl}}$. The resistor also returns the discharged energy to the input
rail.

\subsubsection{Leakage energy and the sizing current}

Per cycle the leakage inductance delivers $E_{\mathrm{lk}} = \tfrac{1}{2}
L_{\mathrm{lk}} I_{\mathrm{pk}}^2$ into the clamp. Two currents matter:
\begin{itemize}[nosep]
  \item \textbf{Normal design-max operation}, $I_{\mathrm{pk}} =
    \SI{2.18}{\ampere}$ (\cref{sec:hv-op-converge}):
    $E_{\mathrm{lk}} = \tfrac{1}{2}\times\SI{0.15}{\micro\henry}\times 2.18^2
    \approx \SI{0.355}{\micro\joule}$, $E_{\mathrm{lk}} f_{\mathrm{sw}}
    \approx \SI{53.0}{\milli\watt}$. This sets the clamp's steady dissipation
    carried in the loss budget (\cref{sec:hv-op-losses}).
  \item \textbf{Worst case}: the current limit bounds the current at turn-off,
    so the sizing case is the guaranteed-maximum
    clamped peak including the CS filter's overshoot, \SI{2.91}{\ampere}
    (\cref{sec:hv-cs-sense}). Hiccup bounds how long the converter holds that
    peak (\cref{sec:hv-ss}) but not the peak itself.
    $E_{\mathrm{lk}} = \tfrac{1}{2}\times\SI{0.15}{\micro\henry}\times 2.91^2
    \approx \SI{0.635}{\micro\joule}$, $E_{\mathrm{lk}} f_{\mathrm{sw}}
    \approx \SI{94.9}{\milli\watt}$.
\end{itemize}
Clamp sizing controls the \des{Q}{HVSW} drain-voltage rating. Thus, it uses the
worst-case current. The loss budget uses the normal maximum-load current.

\subsubsection{Clamp energy multiplier}

The clamp absorbs more than the leakage energy. While \des{D}{HVCL} conducts,
the leakage current decreases linearly from $I_{\mathrm{pk}}$ to zero. The
driving voltage is $V_{\mathrm{cl}}-V_{\mathrm{refl}}$. Here,
$V_{\mathrm{refl}}=\SI{19.25}{\volt}$ (\cref{sec:hv-op-converge}):
\begin{align*}
  t_1 = \frac{L_{\mathrm{lk}} I_{\mathrm{pk}}}{V_{\mathrm{cl}} - V_{\mathrm{refl}}}.
\end{align*}
The delivered charge is $Q_{\mathrm{cl}}=\tfrac{1}{2}I_{\mathrm{pk}}t_1$.
It enters the clamp at $V_{\mathrm{cl}}$:
\begin{align*}
  E_{\mathrm{cl}} = V_{\mathrm{cl}} Q_{\mathrm{cl}}
    = \tfrac{1}{2} L_{\mathrm{lk}} I_{\mathrm{pk}}^2 \cdot
      \frac{V_{\mathrm{cl}}}{V_{\mathrm{cl}} - V_{\mathrm{refl}}}
    = E_{\mathrm{lk}} \cdot
      \frac{V_{\mathrm{cl}}}{V_{\mathrm{cl}} - V_{\mathrm{refl}}}.
\end{align*}
At steady state this equals \des{R}{HVCL}'s dissipation, giving the design
equation the rest of this section solves:
\begin{align*}
  \frac{V_{\mathrm{cl}}^2}{\des{R}{HVCL}}
    = E_{\mathrm{lk}} f_{\mathrm{sw}}
      \frac{V_{\mathrm{cl}}}{V_{\mathrm{cl}} - V_{\mathrm{refl}}}.
\end{align*}

\subsubsection{Target clamp voltage}

Clamp effectiveness sets the minimum $V_{\mathrm{cl}}$. Use at least twice the
reflected voltage (inferred). Switch stress sets the maximum. The total stress
includes $V_{\mathrm{BUS,max}}$, $V_{\mathrm{cl}}$, and the \des{D}{HVCL}
forward drop:
\begin{align*}
  V_{DS,\mathrm{pk}} = V_{\mathrm{BUS,max}} + V_{\mathrm{cl}} + V_{F(\des{D}{HVCL})}.
\end{align*}
This section limits semiconductor voltage stress to $80\,\%$ of the rating.
The analysis assumes a \SI{1.5}{\volt} \des{D}{HVCL} forward drop. This value
is near the high end for a small fast rectifier at ampere-level pulse current
(inferred). With $V_{\mathrm{BUS,max}} = \SI{12.6}{\volt}$:
\begin{align*}
  V_{\mathrm{cl}}^{\max} = (1-m)\times\SI{100}{\volt} - V_{\mathrm{BUS,max}}
    - V_{F(\des{D}{HVCL})} = 80 - 12.6 - 1.5 = \SI{65.9}{\volt}.
\end{align*}
This value is 3.4 times the reflected voltage. Therefore, switch stress sets the
limit. A higher clamp voltage decreases clamp loss and decreases switch margin.
The target is $V_{\mathrm{cl}} = \SI{65.9}{\volt}$.

\subsubsection{\des{R}{HVCL}}

Solving the energy-balance equation for \des{R}{HVCL} at the target, with the
worst-case-corner clamp power:
\begin{align*}
  \des{R}{HVCL} = \frac{V_{\mathrm{cl}}(V_{\mathrm{cl}} - V_{\mathrm{refl}})}
    {E_{\mathrm{lk}} f_{\mathrm{sw}}}
    = \frac{65.9 \times 46.65}{\SI{94.9}{\milli\watt}} \approx \SI{32.4}{\kilo\ohm}.
\end{align*}
Rounding down to the standard value $\des{R}{HVCL} = \SI{30}{\kilo\ohm}$
preserves margin; re-solving the same equation for the resulting clamp voltage
at each of the two operating points,
\begin{align*}
  V_{\mathrm{cl}}^2 - V_{\mathrm{refl}} V_{\mathrm{cl}}
    - E_{\mathrm{lk}} f_{\mathrm{sw}} \des{R}{HVCL} = 0
  \;\Longrightarrow\;
  V_{\mathrm{cl}}^{\mathrm{lim}} = \SI{63.8}{\volt}, \qquad
  V_{\mathrm{cl}}^{\mathrm{dm}} = \SI{50.7}{\volt}
\end{align*}
(sustained current-limit corner and normal design-max operation,
respectively). At the current-limit corner, switch stress is
$V_{DS,\mathrm{pk}} = 12.6 + 63.8 + 1.5 = \SI{77.9}{\volt}$ against the
\SI{100}{\volt} rating. The margin is \SI{22.1}{\percent}. During normal
operation, the clamp voltage is \SI{50.7}{\volt}. This is 2.6 times the
reflected voltage. The clamp dissipates \SI{85.5}{\milli\watt}
(\cref{sec:hv-op-losses}). At the current-limit corner, dissipation is
\SI{136}{\milli\watt}. Hiccup limits this condition to a
\SI{0.2}{\percent} duty cycle (\cref{sec:hv-ss}). Require a
\SI{0.18}{\watt} power rating and an \SI{80}{\volt} voltage rating.

\subsubsection{\des{C}{HVCL}}

The preceding calculations use a constant clamp voltage. Clamp ripple adds to
peak drain stress. Limit ripple to $2\,\%$ of $V_{\mathrm{cl}}$ at the
worst-case current. This limit keeps the stress error small.

The charge delivered during each cycle follows from
$E_{\mathrm{cl}}=V_{\mathrm{cl}}Q_{\mathrm{cl}}$:
\begin{align*}
  Q_{\mathrm{cl}} = \frac{E_{\mathrm{lk}}}{V_{\mathrm{cl}} - V_{\mathrm{refl}}}
    = \frac{\SI{0.635}{\micro\joule}}{44.55} \approx \SI{14.3}{\nano\coulomb}.
\end{align*}
At a $2\,\%$ ripple target ($\Delta V \le \SI{1.28}{\volt}$ at
$V_{\mathrm{cl}}^{\mathrm{lim}} = \SI{63.8}{\volt}$):
\begin{align*}
  \des{C}{HVCL}^{\min} = \frac{Q_{\mathrm{cl}}}{\Delta V}
    = \frac{\SI{14.3}{\nano\coulomb}}{\SI{1.28}{\volt}} \approx \SI{11.2}{\nano\farad}.
\end{align*}
Rounding up to $\des{C}{HVCL} = \SI{12}{\nano\farad}$ gives actual ripple
$Q_{\mathrm{cl}}/\des{C}{HVCL} \approx \SI{1.19}{\volt}$ ($1.9\,\%$) and a
discharge time constant $\tau = \des{R}{HVCL} \des{C}{HVCL} \approx
\SI{360}{\micro\second}$, about $54$ switching periods: the near-DC node the
treatment above assumes. Its voltage rating follows the corner clamp voltage,
$\ge\SI{80}{\volt}$.

\subsubsection{\des{D}{HVCL} requirements}

The design does not specify a part for \des{D}{HVCL}. It must meet these
requirements:
\begin{itemize}[nosep]
  \item \textbf{Reverse voltage.} It must block $V_{\mathrm{BUS,max}} +
    V_{\mathrm{cl}}^{\mathrm{lim}} = 12.6 + 63.8 = \SI{76.4}{\volt}$ during the
    on-time; at the same $20\,\%$ voltage-derating standard used throughout
    this section, that requires a rated reverse voltage of at least
    \SI{95.5}{\volt}, i.e.\ a $\ge\SI{100}{\volt}$-class part.
  \item \textbf{Forward current.} The leakage current decreases from
    \SI{2.91}{\ampere} to zero in approximately \SI{9.8}{\nano\second}. The
    charge is \SI{14.3}{\nano\coulomb}. Its average is approximately
    \SI{2.1}{\milli\ampere}. A rectifier that meets the voltage requirement will
    typically meet this short-pulse requirement (inferred).
  \item \textbf{Reverse recovery.} \des{D}{HVCL} must block again before the
    next turn-on, and the off-interval never falls below
    \SI{4.0}{\micro\second} (\cref{sec:hv-dcm}). A fast
    ($\le\SI{500}{\nano\second}$) or ultrafast ($\le\SI{75}{\nano\second}$) part
    clears that; a standard-recovery part is comparable to the window itself and
    is excluded.
\end{itemize}

\subsection{Switch (\des{Q}{HVSW}) stress and losses}\label{sec:hv-switch}

The clamp sets the maximum drain stress. The worst case is the sustained
current-limit condition (\cref{sec:hv-clamp}):
\begin{align*}
  V_{DS,\mathrm{pk}} &= V_{\mathrm{BUS,max}} + V_{\mathrm{cl}}^{\mathrm{lim}}
    + V_{F(\des{D}{HVCL})} = 12.6 + 63.8 + 1.5 = \SI{77.9}{\volt}
\end{align*}
This stress gives a $22.1\,\%$ margin below the IRL540N
\SI{100}{\volt} rating. It meets the $20\,\%$ derating target. In normal
maximum-load operation, the peak is approximately \SI{64.8}{\volt}. Peak current
spans \SIrange{2.2}{2.9}{\ampere}, below the \SI{36}{\ampere} rating. The
\SI{6.85}{\volt} gate supply fully enhances the logic-level MOSFET. It also
stays inside the $\pm\SI{16}{\volt}$ gate rating.

Total switch dissipation at design-max, at the typical turn-off overlap
($t_{\mathrm{ovl}}^{\mathrm{typ}}=\SI{36}{\nano\second}$, \cref{sec:hv-op-losses}):
conduction \SI{60}{\milli\watt} + turn-off \SI{140}{\milli\watt} + $C_{oss}$ few
\si{\milli\watt} $\approx \SI{0.20}{\watt}$, a rise of $\approx\SI{12}{\kelvin}$
on a free-standing vertical TO-220 at $\theta_{JA} \approx
\SI{62}{\kelvin\per\watt}$. At the maximum \SI{54}{\nano\second} turn-off
overlap, conduction loss is \SI{62}{\milli\watt}. Turn-off loss is
\SI{215}{\milli\watt}. Total dissipation is approximately
\SI{0.28}{\watt}, and the temperature rise is approximately
\SI{17}{\kelvin}.

\textbf{Turn-off overlap is the largest converter loss.}
\des{D}{HVGOFF} and \des{R}{HVGOFF} decrease gate-discharge resistance from
approximately \SI{11.5}{\ohm} to \SI{4.8}{\ohm}. They decrease overlap from
\SIrange{75}{114}{\nano\second} to
\SIrange{36}{54}{\nano\second}. \des{D}{HVGOFF} is reverse-biased during
turn-on. Thus, \des{R}{HVGON} controls turn-on.

\subsection{Output rectifier (\des{D}{HVRECT})}\label{sec:hv-rect}

Reverse stress occurs during the on-time:
\begin{align*}
  V_R &= V_{\mathrm{out}} + n V_{\mathrm{BUS,max}}
       = 191.5 + 10 \times 12.6 = \SI{317.5}{\volt}
\end{align*}
at the maximum output (\cref{sec:hv-achieved-range}). Secondary-side ringing
keeps the total below \SI{345}{\volt}. The UF4005 rating is
\SI{600}{\volt}, which gives at least \SI{42}{\percent} margin. Average forward
current is \SI{15}{\milli\ampere} (\cref{sec:hv-budget}). Peak forward current
is \SI{0.22}{\ampere}. Conduction loss is approximately
\SI{15}{\milli\watt}. In DCM, diode current reaches zero during the idle
interval. Thus, the switch does not hard-commutate the diode.

\subsection{Input and output capacitance}\label{sec:hv-caps}

\subsubsection{Input}

\des{C}{HVI1} is the \SI{100}{\micro\farad} polymer-aluminum bulk capacitor.
\des{C}{HVI2} is a \SI{4.7}{\micro\farad} ceramic capacitor. It closes the
high-frequency primary loop. \des{C}{HVBIAS} bypasses BIAS
(\cref{sec:hv-house-bias}). The primary RMS current is
\SI{1.04}{\ampere}. Its DC component is \SI{0.74}{\ampere}. Thus, the capacitor
bank carries \SI{0.73}{\ampere} RMS. \textbf{Use a polymer-aluminum bulk
capacitor.} A small general-purpose \SI{100}{\micro\farad} aluminum
electrolytic typically has a \SIrange{0.3}{0.5}{\ampere} ripple rating at
\SI{150}{\kilo\hertz} (inferred). This rating is insufficient. The calculated
bus ripple is
\begin{align*}
  \Delta V \approx
  \frac{(I_{\mathrm{pk}}/2 - I_{\mathrm{in}})\,t_{\mathrm{on}}}{C_{\mathrm{in}}}
  + I_{\mathrm{pk}}\,\mathrm{ESR}
  \approx \frac{\SI{1.57}{\micro\coulomb}}{\SI{105}{\micro\farad}} + 2.18\times0.04
  \approx \SI{102}{\milli\volt}\ \text{peak-to-peak},
\end{align*}
unobjectionable on a USB-PD input.

\subsubsection{Output}

$C_{\mathrm{out}} = \des{C}{HVO1}\,(\SI{4.7}{\micro\farad}$, radial polarized, polymer
aluminum$) + \des{C}{HVO2}\,(\SI{100}{\nano\farad}\ \text{film})$.

The tube ballasts make the load approximately constant current. The rail does
not need hold-up energy or a transient reservoir. It needs only
switching-ripple control. Capacitive ripple per cycle is
\begin{align*}
  \Delta V_{\mathrm{cap}} &=
  \frac{I_{\mathrm{out}}\,(T - t_{\mathrm{dem}})}{C_{\mathrm{out}}}
  = \frac{15.4\,\si{\milli\ampere} \times 5.6\,\si{\micro\second}}
         {\SI{4.8}{\micro\farad}}
  = \SI{18}{\milli\volt},
\end{align*}
exceeded by the ESR step $I_{s,\mathrm{pk}}\times\mathrm{ESR} \approx 0.22 \times
0.5 \approx \SI{108}{\milli\volt}$ peak-to-peak ($0.06\,\%$ of the rail) at
\des{C}{HVO1}'s $\approx\SI{0.5}{\ohm}$ ESR, with the film part shunting the
fast edges.
Through a \SI{10}{\kilo\ohm} ballast (\cref{sec:hv-load-digit}) this modulates
tube current by $\sim$\SI{11}{\micro\ampere} at \SI{150}{\kilo\hertz},
optically invisible. The small capacitance is sufficient, and beneficial: it
keeps the stored energy at \SI{69}{\milli\joule}, the start-up rise time at
\SI{17}{\milli\second} (\cref{sec:hv-ss}), and the bleed-down inside seconds
(\cref{sec:hv-bleeder}).

\textbf{\des{C}{HVO1} requires a \SI{250}{\volt} rating.} A
\SI{160}{\volt} part cannot withstand the \SI{191.5}{\volt} maximum setting.
The \SI{250}{\volt} rating gives at least $23\,\%$ margin. The selected part
must also have approximately \SI{0.5}{\ohm} ESR or less.

\subsection{Bleeder (\des{R}{HVBLEED}) and minimum load}\label{sec:hv-bleeder}

\des{R}{HVBLEED} discharges $C_{\mathrm{out}}$ after shutdown. It also supplies
the minimum load that prevents pulse skipping.

\subsubsection{Discharge}\label{sec:hv-bleeder-discharge}

Stored energy is modest, $\tfrac{1}{2} C_{\mathrm{out}} V^2 =
\tfrac{1}{2}\times\SI{4.8}{\micro\farad}\times 170^2 = \SI{69}{\milli\joule}$
(\cref{sec:hv-caps}), but the node must still bleed down after shutdown.

IEC 62368-1 classifies DC voltages at or below \SI{60}{\volt} as ES1. It
requires the stored voltage to reach this limit within \SI{2}{\second} after
disconnection. The tubes help discharge the rail. The digit branches pull the
rail to approximately \SI{148}{\volt} in \SI{28}{\milli\second}. They then stop
conducting. The colon branches continue to conduct toward their maximum
\SI{55}{\volt} maintaining voltage. Thus, the colon tubes can pull the rail
below the ES1 limit. \des{R}{HVBLEED} and the feedback divider set the discharge
time.

Two \SI{240}{\kilo\ohm} ballasts against $\des{R}{HVBLEED} = \SI{1.5}{\mega\ohm}$
in parallel with the feedback divider ($R_{\mathrm{div}} \approx
\SI{1.70}{\mega\ohm}$, \cref{sec:hv-divider}) give
\begin{align*}
  \des{R}{HVBLEED} \parallel R_{\mathrm{div}} &= 1.5\,\si{\mega\ohm} \parallel
    1.70\,\si{\mega\ohm} = \SI{797}{\kilo\ohm}, \\
  \tau &= \frac{C_{\mathrm{out}}}{2/\SI{240}{\kilo\ohm} + 1/\SI{797}{\kilo\ohm}}
       = \SI{0.50}{\second}, &
  V_\infty &= \frac{2\times\SI{55}{\volt}/\SI{240}{\kilo\ohm}}
                   {2/\SI{240}{\kilo\ohm} + 1/\SI{797}{\kilo\ohm}}
            = \SI{47.8}{\volt},\\
  t_{148\rightarrow 60\,\mathrm{V}} &= \tau
    \ln\frac{148 - 47.8}{60 - 47.8} = \SI{1.05}{\second}.
\end{align*}
The total time to \SI{60}{\volt} is approximately \SI{1.08}{\second}. This time
meets the \SI{2}{\second} limit with a \SI{46}{\percent} margin. It applies to a
board with all tubes installed. With the tubes removed, discharge takes
\SI{4.0}{\second}. This condition occurs during high-voltage adjustment.

The minimum-load requirement sets \des{R}{HVBLEED} to no more than
\SI{2.92}{\mega\ohm}. Use \SI{1.5}{\mega\ohm}. It dissipates
\SI{19}{\milli\watt} at the nominal output and \SI{24}{\milli\watt} at the
maximum output. The resistor sees \SI{191.5}{\volt}. With
\SI{20}{\percent} derating, require an element voltage rating of at least
\SI{239}{\volt}. Use a DIN0309 body.

\subsubsection{Minimum load}\label{sec:hv-bleeder-minload}

The controller cannot deliver less than one minimum-width pulse. At
\SI{12.6}{\volt} input and an approximately \SI{165}{\nano\second} minimum on
time:
\begin{align*}
  I_{\mathrm{pk}}^{\min} &= \frac{V_{\mathrm{BUS}} t_{\mathrm{on,min}}}{L_p}
    = \frac{12.6 \times 165\,\si{\nano\second}}{\SI{10}{\micro\henry}}
    = \SI{0.208}{\ampere},\\
  P_{\min} &= \tfrac{1}{2} L_p (I_{\mathrm{pk}}^{\min})^2 f_{\mathrm{sw}}\,\eta
    \approx \tfrac{1}{2}\times\SI{10}{\micro\henry}\times 0.0432 \times
      149.4\,\si{\kilo\hertz} \times 0.833
      \approx \SI{27}{\milli\watt}.
\end{align*}
Below this load, the output increases or the controller skips pulses. The
bleeder supplies \SI{19}{\milli\watt}. The divider supplies
\SI{17}{\milli\watt}. Their \SI{36}{\milli\watt} total is 1.34 times
$P_{\min}$. Thus, the rail regulates before the tubes ignite. The divider alone
is insufficient. The remaining load requires \des{R}{HVBLEED} to be no more
than \SI{2.92}{\mega\ohm}.

\subsection{Loop compensation}\label{sec:hv-comp}

\subsubsection{Plant}

A first-order model of a peak-current-mode DCM flyback is a controlled current
source into the output capacitance and the incremental load. At the nominal
point, the load conductance is
\begin{align*}
  G_L &= \frac{4}{\SI{7.5}{\kilo\ohm}}+
         \frac{2}{\SI{240}{\kilo\ohm}}+
         \frac{1}{\SI{1.5}{\mega\ohm}}+
         \frac{1}{\SI{1.70}{\mega\ohm}}
       = \SI{543}{\micro\siemens}.
\end{align*}
This conductance uses the slope of each load line. The tube maintaining voltages
and driver drop are constant offsets. They do not appear in the incremental
conductance. The output capacitance is $C_{\mathrm{out}} = \des{C}{HVO1} +
\des{C}{HVO2} = \SI{4.8}{\micro\farad}$ (\cref{sec:hv-caps}).

Modulator gain: the controller's PWM comparator sees COMP through its internal
attenuation $G_{\mathrm{COMP}} = 0.142$ and compares it against the sensed current
plus the fixed internal slope ramp of slope $S_e = \SI{40}{\milli\volt}\times
f_{\mathrm{sw}} = \SI{5975}{\volt\per\second}$ (\ds{lm5156h.pdf}{20}). Against
the sensed on-slope $S_n = R_S V_{\mathrm{BUS}}/L_p = 0.039\times9/\SI{10}{\micro\henry}
= \SI{35100}{\volt\per\second}$ (\cref{sec:hv-cs}):
\begin{align*}
  \frac{\partial I_{\mathrm{pk}}}{\partial v_c}
    &= \frac{G_{\mathrm{COMP}}}{R_S\,(1 + S_e/S_n)}
     = \frac{0.142}{0.039 \times 1.17}
     = \SI{3.11}{\ampere\per\volt}.
\end{align*}
Since $I_{\mathrm{out}} = \tfrac{1}{2} L_p I_{\mathrm{pk}}^2 f_{\mathrm{sw}}
\eta / V_{\mathrm{out}}$, the plant transconductance and its output resistance
follow:
\begin{align*}
  \frac{\partial I_{\mathrm{out}}}{\partial I_{\mathrm{pk}}}
    &= \frac{2 I_{\mathrm{out}}}{I_{\mathrm{pk}}}
     = \frac{2 \times 13.97\,\si{\milli\ampere}}{1.96}
     = 14.3\,\si{\milli\ampere\per\ampere},\\
     g_{ps} &= 14.3\,\si{\milli\ampere\per\ampere} \times \SI{3.11}{\ampere\per\volt}
     = \SI{44.3}{\milli\siemens} \text{ (per volt of COMP)},
\end{align*}
and because $I_{\mathrm{out}} \propto 1/V_{\mathrm{out}}$ at fixed peak current,
the converter adds conductance $I_{\mathrm{out}}/V_{\mathrm{out}} =
\SI{82.2}{\micro\siemens}$. The effective output resistance and dominant pole
are therefore
\begin{align*}
  R_{\mathrm{out}} &= \frac{1}{G_L+I_{\mathrm{out}}/V_{\mathrm{out}}}
       = \SI{1.60}{\kilo\ohm},\\
  G_0 &= g_{ps}R_{\mathrm{out}} = 70.9\ (\SI{37.0}{\decibel}), &
  f_p &= \frac{1}{2\pi R_{\mathrm{out}} C_{\mathrm{out}}}
       = \SI{20.7}{\hertz}.
\end{align*}
The DCM right-half-plane zero and the output capacitor's ESR zero
($\approx\SI{67}{\kilo\hertz}$ at \des{C}{HVO1}'s \SI{0.5}{\ohm}) both lie far
above the crossover derived below and are neglected.

\subsubsection{Compensator and loop}

The crossover $f_c$ must be above the \SI{20.7}{\hertz} plant pole. It must also
be below $f_{\mathrm{sw}}/10$, or approximately \SI{15}{\kilo\hertz}. The load
changes slowly because each tube has a dedicated driver. The USB-PD input is
also regulated. Therefore, the design uses a crossover of a few hundred hertz.
The Type-II zero is 0.7 to 0.8 times $f_c$. The high-frequency pole is more
than 10 times $f_c$.

The mid-band gain is $g_m\des{R}{HVCOMP}$. The error-amplifier transconductance
is \SI{2}{\milli\siemens} (\ds{lm5156h.pdf}{6}). The resistor and capacitor
ratios set the zero and high-frequency pole. Use
$\des{R}{HVCOMP}=\SI{10}{\kilo\ohm}$ in series with
$\des{C}{HVC1}=\SI{100}{\nano\farad}$. Shunt this network with
$\des{C}{HVC2}=\SI{220}{\pico\farad}$:
\begin{align*}
  f_z &= \frac{1}{2\pi \des{R}{HVCOMP} \des{C}{HVC1}} = \SI{159}{\hertz},
  &
  f_{hp} &= \frac{1}{2\pi \des{R}{HVCOMP}\,\frac{\des{C}{HVC1}\des{C}{HVC2}}{\des{C}{HVC1}+\des{C}{HVC2}}}
          = \SI{72.5}{\kilo\hertz},
\end{align*}
with integrator behavior below $f_z$ (the amplifier's
$g_m R_O = 2\,\si{\milli\siemens} \times 10\,\si{\mega\ohm} = \SI{86}{\decibel}$
DC gain, \ds{lm5156h.pdf}{25}, makes static regulation error negligible). The
divider contributes $H = V_{\mathrm{REF}}/V_{\mathrm{out}} = 1/170$.

Solving $|T(f)| = 1$ for the product of divider, compensator, and plant
(numerically, with all three break frequencies in place) confirms the target:
\begin{align*}
  f_c &\approx \SI{214}{\hertz} \approx \frac{f_{\mathrm{sw}}}{699},\\
  \varphi_m &= 180^\circ - \arctan\frac{f_c}{f_p} - 90^\circ
    + \arctan\frac{f_c}{f_z} - \arctan\frac{f_c}{f_{hp}}
    = 180 - 84.5 - 90 + 53.3 - 0.2 \approx 59^\circ.
\end{align*}
The zero is $0.74 f_c$. The high-frequency pole is 339 times the crossover
frequency. Both values meet their targets. The phase reaches
$-180^\circ$ only after the loop gain has decreased by tens of decibels. The
datasheet gives no tolerance for $G_{\mathrm{COMP}}$, $S_e$, or $g_m$. Verify
the loop with a transient test or Bode measurement.

\subsection{Soft start}\label{sec:hv-ss}

The datasheet soft-start equations apply only to boost and SEPIC converters.
The boost equation assumes an input-to-output conduction path. A flyback has no
such path. The \SI{10}{\micro\ampere} SS current charges \des{C}{HVSS}. The
error amplifier regulates FB to the lower of $V_{SS}$ and
$V_{\mathrm{REF}}$. Therefore,
$t_{SS} = \des{C}{HVSS} V_{\mathrm{REF}} / I_{SS}$.

The converter must supply enough current for the output to follow the ramp. At
the current limit, it charges $C_{\mathrm{out}}$ from zero to
$V_{\mathrm{out}}$. The tubes draw no current before ignition. The charge time
is
\begin{align*}
  t_{\mathrm{rise}} &\approx
  \frac{\tfrac{1}{2} C_{\mathrm{out}} V_{\mathrm{out}}^2}
       {\tfrac{1}{2} L_p I_{\mathrm{lim}}^2 f_{\mathrm{sw}}\,\eta}
  = \frac{\SI{69}{\milli\joule}}{\SI{4.1}{\watt}}
  \approx \SI{17.0}{\milli\second}.
\end{align*}
$\des{C}{HVSS}$ must be large enough that $t_{SS}$ exceeds $t_{\mathrm{rise}}$;
otherwise the output falls behind the commanded ramp and the error amplifier
saturates into the current limit. Setting $t_{SS} = t_{\mathrm{rise}}$ gives the
minimum:
\begin{align*}
  \des{C}{HVSS}^{\min} &= \frac{t_{\mathrm{rise}} I_{SS}}{V_{\mathrm{REF}}}
  = \frac{\SI{17.0}{\milli\second} \times \SI{10}{\micro\ampere}}{\SI{1.00}{\volt}}
  \approx \SI{170}{\nano\farad}.
\end{align*}
\paragraph{Hiccup constraint.}
The LM51561H specified here enables hiccup mode overload protection, which the
LM5156H option omits (\ds{lm5156h.pdf}{2}). A fault timer counts current-limited
cycles, but only once soft start has finished; 64 of them stop the converter and
pull SS down, it restarts after 32{,}768 cycles, and eight consecutive
unlimited cycles reset the count (\ds{lm5156h.pdf}{6,25}). At
\SI{149.4}{\kilo\hertz} that is a \SI{428}{\micro\second} fault window and a
\SI{219}{\milli\second} retry interval, a fault duty of $0.2\,\%$. The datasheet
states the requirement directly: soft start must be long enough not to trigger
hiccup once it finishes.

The converter charges $C_{\mathrm{out}}$ at its current limit. If the output
does not reach its target before soft start ends, current limiting continues.
Hiccup then operates after \SI{428}{\micro\second}. Therefore, $t_{SS}$ must
exceed the worst-case rise time. The calculation uses maximum output
capacitance, minimum current limit, and minimum efficiency. It gives
$t_{\mathrm{rise}}^{\max} \approx \SI{24.5}{\milli\second}$.

The nominal calculation would select \SI{220}{\nano\farad}. Its minimum
$t_{SS}$ is \SI{17.8}{\milli\second}, which is insufficient.
$\des{C}{HVSS} = \SI{470}{\nano\farad}$ gives a nominal
$t_{SS} = \SI{47.0}{\milli\second}$. Stacking
$\des{C}{HVSS}$'s own tolerance ($\pm10\,\%$, typical for this class of ceramic)
against $I_{SS}$ and $V_{\mathrm{REF}}$'s datasheet ranges
(\SIrange{9}{11}{\micro\ampere} and \SIrange{0.99}{1.01}{\volt},
\ds{lm5156h.pdf}{6}), worst-case-additive:
\begin{align*}
  t_{SS}^{\min} &= \frac{\SI{470}{\nano\farad}\times0.9 \times \SI{0.99}{\volt}}
  {\SI{11}{\micro\ampere}} \approx \SI{38.1}{\milli\second},
\end{align*}
\SI{13.5}{\milli\second} clear of $t_{\mathrm{rise}}^{\max}$. The full range is
\SIrange{38.1}{58.0}{\milli\second}.

Hiccup protection limits the duration of a short circuit. Without it, a shorted
tube socket can hold \des{Q}{HVSW} and \des{T}{HV} at the current limit.

\subsection{Summary}\label{sec:hv-summary}

The analysis uses the following inferred values: a
\SIrange{1.0}{2.5}{\volt} driver on-state drop, typical IRL540N
$R_{DS(\mathrm{on})}$ and $Q_G$, a \SI{0.5}{\volt} \des{D}{HVGOFF} forward
drop, a \SI{1.5}{\volt} \des{D}{HVCL} forward drop, a
\SI{50}{\milli\watt} transformer core-loss allowance, a
\SI{50}{\nano\second} comparator delay, and VCC regulator dropout below its
\SI{8}{\volt} specification point. The analysis does not quantify
\des{Q}{HVSW} output capacitance. The DCM idle margin must include its drain
ring-down interval.

\paragraph{Gate-charge effect on efficiency.}\label{sec:hv-summary-lever}
Bias current and turn-off overlap total
\SIrange{297}{400}{\milli\watt}. Both depend on \des{Q}{HVSW} gate charge. The
charge pump supplies this charge from a \SIrange{24}{32}{\volt} node. Thus,
gate-charge loss is approximately three times the loss from a direct input feed.
A \SI{20}{\nano\coulomb} part with similar on-resistance would increase maximum
load efficiency from \SI{83.3}{\percent} to approximately
\SI{87.1}{\percent}. No suitable through-hole part was identified (inferred).
The design therefore specifies the IRL540N.

{\footnotesize
\begin{longtable}{|p{0.23\textwidth}|p{0.335\textwidth}|p{0.335\textwidth}|}
\caption{170~V flyback converter: specification summary}\label{tab:hv-converter-summary}\\
\hline
\textbf{Parameter} & \textbf{Derived requirement} & \textbf{Specification / result} \\
\hline
\endfirsthead
\caption[]{170~V flyback converter: specification summary (continued)}\\
\hline
\textbf{Parameter} & \textbf{Derived requirement} & \textbf{Specification / result} \\
\hline
\endhead
\hline
\endlastfoot
$R_T$ & \SI{146.4}{\kilo\ohm} for \SI{150}{\kilo\hertz} & \SI{147}{\kilo\ohm}
  $\rightarrow$ \SI{149.4}{\kilo\hertz} \\
\hline
DA2032 frequency & characterized at \SI{100}{\kilo\hertz} only &
  \SI{149.4}{\kilo\hertz}; \SI{50}{\milli\watt} core-loss allowance \\
\hline
UVLO \des{R}{HVUV1}/\des{R}{HVUV2} & enable below the \SI{4.75}{\volt} minimum input & on
\SI{4.44}{\volt} / off \SI{3.97}{\volt}; \SI{63.4}{\kilo\ohm}/\SI{32.4}{\kilo\ohm} \\
\hline
Display inhibit \des{Q}{HVINH}/\des{Q}{HVREL} & default asserted with no
dependency on a rail outside the converter; released only on a contract whose
\emph{current} allowance carries the display (\cref{sec:pd-sequence}) & shunt across \des{R}{HVUV2}, gate
divided from $V_{\mathrm{BUS}}$ through \SI{100}{\kilo\ohm}/\SI{1}{\mega\ohm};
asserted by \SI{2.22}{\volt} in against a \SI{4.16}{\volt} earliest UVLO
release; \SI{199}{\micro\ampere}, $V_{GS(th)}\le\SI{2}{\volt}$ maximum \\
\hline
EN and UVLO & EN satisfied before UVLO releases & \SI{1.54}{\volt} $\ll$
\SI{4.44}{\volt} \\
\hline
PGOOD pull-up \des{R}{HVPG} & $\ge\SI{10}{\kilo\ohm}$, $\le\SI{18}{\volt}$,
  $\le\SI{1}{\milli\ampere}$ & \SI{100}{\kilo\ohm} to +3V3;
  \SI{33}{\micro\ampere} \\
\hline
VCC capacitor \des{C}{HVVCC} & \SIrange{1}{4.7}{\micro\farad} & \SI{2.2}{\micro\farad} \\
\hline
Bias charge pump \des{C}{HVPUMP}/\des{R}{HVPUMP} & hold VCC at \SI{6.85}{\volt} across the
whole input window; $V_{\mathrm{BIAS}}\ge\SI{8}{\volt}$ at the worst delivery corner &
\SI{4.7}{\nano\farad} C0G, \SI{82}{\ohm}; $V_{\mathrm{BIAS}}$ \SI{9.1}{\volt} worst
corner / \SI{20.4}{\volt} at \SI{12.6}{\volt} in; costs \SI{126}{\milli\watt} flat \\
\hline
\des{D}{HVSTART}, \des{D}{HVZ}, \des{C}{HVBIAS} & start-up path, fault clamp, reservoir &
\SI{40}{\volt} Schottky; \SI{30}{\volt} zener (non-conducting in normal operation);
\SI{1}{\micro\farad} \\
\hline
Tuning range & span \SIrange{170}{185}{\volt} (ignition margin) &
  \SIrange{164.8}{191.5}{\volt} nominal targets;
  \SIrange{146.5}{212.1}{\volt} informational stop-to-stop
  (\cref{tab:hv-ballast-divider}); end stops are impermissible setpoints \\
\hline
Digit ballast \des{R}{BALDIG} & rated digit current at every trim setpoint &
installed \SI{7.50}{\kilo\ohm}; permissible alternatives through
\SI{13.00}{\kilo\ohm}, with \des{R}{HVD3} matched per row
  (\cref{tab:hv-ballast-divider}) \\
\hline
Colon ballast \des{R}{BALCOL} & \SI{0.5}{\milli\ampere} brightness optimum, no
  rated band & \SI{240}{\kilo\ohm}; \SIrange{0.458}{0.569}{\milli\ampere}
  across the achieved range \\
\hline
\des{R}{HVD1}/\des{R}{HVD2} element stress & $<\SI{250}{\volt}$ per DIN0207 element &
  \SI{95.25}{\volt} max at design-max \\
\hline
$P_{\mathrm{out}}$ & Not applicable & \SI{2.38}{\watt} nominal / \SI{2.94}{\watt}
  design-max \\
\hline
Efficiency & Not applicable & design-max \SIrange{82.8}{83.6}{\percent} across the window,
peaking near \SI{7.4}{\volt}; nominal \SI{82.7}{\percent} / \SI{82.5}{\percent}
/ \SI{81.7}{\percent} at
\SI{4.75}{\volt} / \SI{9}{\volt} / \SI{12.6}{\volt} \\
\hline
Turn-off speedup \des{D}{HVGOFF}/\des{R}{HVGOFF} & reduce turn-off overlap; driver peak
  $\le\SI{1.5}{\ampere}$ & SB240S-E3/54 + \SI{3.3}{\ohm};
  $t_{\mathrm{ovl}}$ \SIrange{36}{54}{\nano\second}; \SI{1.32}{\ampere} peak \\
\hline
$I_{\mathrm{pk}}$ & Not applicable & \SI{1.96}{\ampere} nominal; \SI{2.18}{\ampere}
design-max; \SI{2.29}{\ampere} at $L_{\min}$ \\
\hline
DCM & $t_{\mathrm{on}}+t_{\mathrm{dem}} < T$ & $0.89\,T$ at the worst corner
(\SI{4.75}{\volt}, $L_p+10\,\%$, design-max): $10.5\,\%$ idle \\
\hline
Sense resistor \des{R}{HVCS} & \SIrange{37.8}{40.2}{\milli\ohm} worst-case-additive &
  \SI{39}{\milli\ohm} \\
\hline
Current-limit margin & Not applicable & $+1.7\,\%$ worst-case-additive; $+7.2\,\%$ RSS \\
\hline
CS filter \des{R}{HVCSF}/\des{C}{HVCSF} & $\tau$ \SIrange{30}{100}{\nano\second} & \SI{100}{\ohm},
  \SI{470}{\pico\farad}; $\tau=\SI{47}{\nano\second}$ \\
\hline
Clamped peak and saturation & Not applicable & \SI{2.91}{\ampere} and
  \SI{3.0}{\ampere} soft-saturation rating \\
\hline
Clamp \des{R}{HVCL}/\des{C}{HVCL} & $V_{\mathrm{cl}} \ge 2\times$ reflected; $20\,\%$ \des{Q}{HVSW}
  voltage derating at the current-limit corner & \SI{30}{\kilo\ohm},
  \SI{12}{\nano\farad}, both $\ge\SI{80}{\volt}$; $V_{\mathrm{cl}}$
  \SI{63.8}{\volt} corner / \SI{50.7}{\volt} design-max \\
\hline
\des{D}{HVCL} & $\ge\SI{100}{\volt}$; \SI{2.91}{\ampere}, \SI{14.3}{\nano\coulomb}
  pulse; fast or ultrafast recovery & Not applicable \\
\hline
\des{R}{HVCL} dissipation & rating $\ge 2\times$ sustained dissipation &
\SI{85.5}{\milli\watt} design-max; the \SI{136}{\milli\watt} corner is
hiccup-limited to a $0.2\,\%$ duty and does not set the rating. Rated
$\ge\SI{0.18}{\watt}$, $\ge\SI{80}{\volt}$ \\
\hline
\des{Q}{HVSW} $V_{DS}$ & $\le\SI{77.9}{\volt}$ at the current-limit corner &
  \SI{100}{\volt} rating; $22.1\,\%$ margin \\
\hline
\des{Q}{HVSW} dissipation & Not applicable & \SI{0.20}{\watt} typical / \SI{0.28}{\watt} maximum
$t_{\mathrm{ovl}}$; $\Delta T \approx \SIrange{12}{17}{\kelvin}$, TO-220, no
  heatsink \\
\hline
\des{D}{HVRECT} & $V_R < \SI{345}{\volt}$; \SI{15}{\milli\ampere} avg & UF4005:
  \SI{600}{\volt}, \SI{1}{\ampere} \\
\hline
Input bulk \des{C}{HVI1} & $\ge\SI{0.73}{\ampere}$ ripple capability &
\SI{100}{\micro\farad} polymer aluminum \\
\hline
Output caps \des{C}{HVO1}/\des{C}{HVO2} & rating $\ge\SI{250}{\volt}$ at the \SI{191.5}{\volt}
  envelope; switching-ripple control only & \SI{4.7}{\micro\farad} polymer +
  \SI{100}{\nano\farad} film, both \SI{250}{\volt} \\
\hline
Bleeder \des{R}{HVBLEED} & preload $\ge\SI{27}{\milli\watt}$ sets the value
  ($\le\SI{2.92}{\mega\ohm}$); $\le\SI{60}{\volt}$ within \SI{2}{\second}
  (IEC 62368-1 ES1) & \SI{1.5}{\mega\ohm}:
  \SI{36}{\milli\watt}, \SI{1.08}{\second} with tubes installed \\
\hline
Compensation \des{R}{HVCOMP}/\des{C}{HVC1}/\des{C}{HVC2} & $\varphi_m \ge 45^\circ$ & \SI{10}{\kilo\ohm},
  \SI{100}{\nano\farad}, \SI{220}{\pico\farad};
  $f_c\approx\SI{214}{\hertz}$, $\varphi_m\approx59^\circ$ \\
  \hline
  Soft start \des{C}{HVSS} & $t_{SS} > t_{\mathrm{rise}}^{\max} \approx
  \SI{24.5}{\milli\second}$, else hiccup trips at every power-up &
  \SI{470}{\nano\farad}; $t_{SS}$ \SI{47.0}{\milli\second} nominal,
  \SIrange{38.1}{58.0}{\milli\second} worst-case-additive \\
  \hline
  Hiccup (LM51561H) & bound the sustained current-limit corner & 64-cycle
  (\SI{428}{\micro\second}) fault window, \SI{219}{\milli\second} retry:
  $0.2\,\%$ fault duty \\
  \hline
  Spread spectrum & Not applicable & DITHOFF to VCC, dithering disabled: it would cost a
  third of the DCM idle margin (\cref{sec:hv-house-dithoff}) \\
  \hline
  Input current & carried by the negotiated contract (\cref{sec:hv-inhibit}) &
  \SI{744}{\milli\ampere} design-max at \SI{4.75}{\volt}; \SI{605}{\milli\ampere}
  nominal at \SI{4.75}{\volt} \\
  \hline
  \end{longtable}
  }

  Select the ballast and \des{R}{HVD3} as a pair from
  \cref{tab:hv-ballast-divider}. The gate speedup network decreases turn-off
  overlap from \SIrange{75}{114}{\nano\second} to
  \SIrange{36}{54}{\nano\second}. The bias pump is also necessary at minimum
  input. Without it, overlap reaches \SI{171}{\nano\second} and the converter
  cannot supply the maximum load. The minimum-input DCM idle margin is
  $10.5\,\%$. The additive current-limit margin is $1.7\,\%$. The bias pump
  dissipates \SI{126}{\milli\watt} because it supplies gate charge from the
  reflected-voltage node.

  Measure gate-discharge current, efficiency, current limit, and start-up
  voltage. These tests verify the main derived parameters.

  \subsection{Component selection}\label{sec:hv-parts}

  \Cref{tab:hv-parts} gives the selection criteria. Use through-hole parts when
  they are available. Use carbon-film bodies for the visible colon ballasts.
  The body column gives a requirement only when package size affects the
  analysis or an exact part is specified. Other parts can use any body that
  meets the electrical requirements.

  Resistor sizes use $50\,\%$ power derating and $20\,\%$ voltage derating. Use
  \SI{0.6}{\watt} and \SI{250}{\volt} for DIN0207 metal-film resistors
  (inferred). Use \SI{0.25}{\watt} and \SI{250}{\volt} for DIN0207 carbon-film
  resistors (inferred). Use a \SI{350}{\volt} DIN0309 voltage rating
  (inferred). Verify these values for the selected resistor series.

  {\footnotesize
  \begin{longtable}{|L{0.11\textwidth}|L{0.13\textwidth}|L{0.10\textwidth}|L{0.33\textwidth}|L{0.13\textwidth}|}
  \caption{170~V flyback converter: component selection}\label{tab:hv-parts}\\
  \hline
  \textbf{Designator} & \textbf{Class} & \textbf{Value} & \textbf{Required ratings and characteristics} & \textbf{Body} \\
  \hline
  \endfirsthead
  \caption[]{170~V flyback converter: component selection (continued)}\\
  \hline
  \textbf{Designator} & \textbf{Class} & \textbf{Value} & \textbf{Required ratings and characteristics} & \textbf{Body} \\
  \hline
  \endhead
  \hline
  \endlastfoot
  Not applicable & Flyback controller & LM51561H & hiccup-enabled variant (\cref{sec:hv-ss}) & HTSSOP-14 (PWP), thermal pad \\
  \hline
  \des{T}{HV} & Flyback transformer & DA2032-AL & 1:10; $L_p = \SI{10}{\micro\henry} \pm 10\,\%$; $I_{\mathrm{sat}} \geq \SI{3.0}{\ampere}$; $L_{\mathrm{lk}} \leq \SI{0.15}{\micro\henry}$; secondary DCR \SI{1.6}{\ohm} (\cref{sec:hv-op-losses}) & through-hole bobbin \\
  \hline
  \des{Q}{HVSW} & N-channel MOSFET & IRL540N & $V_{DS}\geq\SI{100}{\volt}$; $V_{GS}\geq\pm\SI{16}{\volt}$; logic-level, $R_{DS(\mathrm{on})}$ specified at $V_{GS}=\SI{5}{\volt}$ and $\leq\SI{55}{\milli\ohm}$ at \SI{6.85}{\volt}; $Q_{GD}\leq\SI{38}{\nano\coulomb}$; $I_D\geq\SI{3}{\ampere}$; $\theta_{JA}\leq\SI{62}{\kelvin\per\watt}$ (\cref{sec:hv-switch}) & TO-220, vertical \\
  \hline
  \des{Q}{HVINH} & N-channel MOSFET & TN2106N3-G & $V_{GS}\geq\pm\SI{12}{\volt}$; $V_{GS(th)}\leq\SI{2}{\volt}$ maximum (\cref{sec:hv-inhibit}) & TO-92 \\
  \hline
  \des{Q}{HVREL} & N-channel MOSFET & TN2106N3-G & $V_{DS}\geq\SI{12}{\volt}$; $V_{GS(th)}\leq\SI{2}{\volt}$ maximum (\cref{sec:hv-inhibit}) & TO-92 \\
  \hline
  \des{R}{BALDIG} & Fixed resistor & \SI{7.50}{\kilo\ohm} installed & 1\,\%; four; permissible alternatives require the matched \des{R}{HVD3} from \cref{tab:hv-ballast-divider}; $\geq\SI{0.32}{\watt}$ (\cref{sec:hv-load-digit}) & DIN0207 minimum \\
  \hline
  \des{R}{BALCOL} & Fixed resistor & \SI{240}{\kilo\ohm} & 5\,\%; two; $\geq\SI{0.16}{\watt}$; element voltage $\geq\SI{171}{\volt}$ (\cref{sec:hv-load-colon}) & DIN0207 carbon film \\
  \hline
  \des{R}{HVCS} & Current sense & \SI{39}{\milli\ohm} & 1\,\%; $\geq\SI{0.09}{\watt}$; low inductance (\cref{sec:hv-cs-sense}) & 2512 \\
  \hline
  \des{R}{HVCSF} & Fixed resistor & \SI{100}{\ohm} & 5\,\% (\cref{sec:hv-cs-filter}) & Not applicable \\
  \hline
  \des{C}{HVCSF} & Ceramic & \SI{470}{\pico\farad} & C0G; $\geq\SI{50}{\volt}$ (\cref{sec:hv-cs-filter}) & Not applicable \\
  \hline
  \des{R}{HVD1}, \des{R}{HVD2} & Fixed resistor & \SI{845}{\kilo\ohm} & 1\,\%; element voltage $\geq\SI{120}{\volt}$ (\cref{sec:hv-divider-stress}) & DIN0207 minimum \\
  \hline
  \des{R}{HVD3} & Fixed resistor & \SIrange{8.25}{9.09}{\kilo\ohm} & 1\,\%; matched to the \des{R}{BALDIG} row of \cref{tab:hv-ballast-divider} & Not applicable \\
  \hline
  \des{RV}{HVTRIM} & Cermet trimmer & \SI{2}{\kilo\ohm} & multiturn; $\pm10\,\%$ total resistance (\cref{sec:hv-divider-split}) & Not applicable \\
  \hline
  \des{R}{HVRT} & Fixed resistor & \SI{147}{\kilo\ohm} & 1\,\% (\cref{sec:hv-fsw}) & Not applicable \\
  \hline
  \des{R}{HVUV1} & Fixed resistor & \SI{63.4}{\kilo\ohm} & 1\,\% (\cref{sec:hv-uvlo}) & Not applicable \\
  \hline
  \des{R}{HVUV2} & Fixed resistor & \SI{32.4}{\kilo\ohm} & 1\,\% (\cref{sec:hv-uvlo}) & Not applicable \\
  \hline
  \des{R}{HVCOMP} & Fixed resistor & \SI{10}{\kilo\ohm} & 5\,\% (\cref{sec:hv-comp}) & Not applicable \\
  \hline
  \des{C}{HVC1} & Ceramic & \SI{100}{\nano\farad} & X7R; $\geq\SI{16}{\volt}$ (\cref{sec:hv-comp}) & Not applicable \\
  \hline
  \des{C}{HVC2} & Ceramic & \SI{220}{\pico\farad} & C0G; $\geq\SI{16}{\volt}$ (\cref{sec:hv-comp}) & Not applicable \\
  \hline
  \des{R}{HVGON} & Fixed resistor & \SI{10}{\ohm} & 5\,\% (\cref{sec:hv-switch}) & Not applicable \\
  \hline
  \des{R}{HVGOFF} & Fixed resistor & \SI{3.3}{\ohm} & 5\,\%; \SI{1.32}{\ampere} peak for \SI{36}{\nano\second} (\cref{sec:hv-switch}) & Not applicable \\
  \hline
  \des{D}{HVGOFF} & Schottky rectifier & SB240S-E3/54 & $V_R\geq\SI{20}{\volt}$; $V_F\leq\SI{0.5}{\volt}$ at \SI{1.3}{\ampere} (\cref{sec:hv-op-losses}) & DO-41 axial \\
  \hline
  \des{R}{HVCL} & Fixed resistor & \SI{30}{\kilo\ohm} & 5\,\%; $\geq\SI{0.18}{\watt}$; $\geq\SI{80}{\volt}$ (\cref{sec:hv-clamp}) & DIN0207 minimum \\
  \hline
  \des{C}{HVCL} & Metallized polypropylene film & \SI{12}{\nano\farad} & $\pm5\,\%$; $\geq\SI{80}{\volt}$; pulse-rated to \SI{300}{\ampere\per\micro\second}; minimum available lead spacing (\cref{sec:hv-clamp}) & Not applicable \\
  \hline
  \des{D}{HVCL} & Fast rectifier & To be selected & $V_R\geq\SI{100}{\volt}$; recovery $\leq\SI{500}{\nano\second}$; \SI{2.91}{\ampere} repetitive peak, \SI{14.3}{\nano\coulomb} per pulse (\cref{sec:hv-clamp}) & Not applicable \\
  \hline
  \des{D}{HVRECT} & Ultrafast rectifier & UF4005 & $V_R\geq\SI{600}{\volt}$; $I_{F(AV)}\geq\SI{1}{\ampere}$; $I_{F,\mathrm{pk}}\geq\SI{0.22}{\ampere}$ (\cref{sec:hv-rect}) & DO-41 axial \\
  \hline
  \des{C}{HVPUMP} & Ceramic & \SI{4.7}{\nano\farad} & C0G; $\pm5\,\%$; $\geq\SI{100}{\volt}$ (\cref{sec:hv-house-pump}) & Not applicable \\
  \hline
  \des{R}{HVPUMP} & Fixed resistor & \SI{82}{\ohm} & 5\,\%; $\geq\SI{0.21}{\watt}$; \SI{118}{\milli\ampere} peak (\cref{sec:hv-house-pump}) & DIN0207 minimum \\
  \hline
  \des{D}{HVPUMP}, \des{D}{HVRESET} & Fast rectifier & To be selected & $V_R\geq\SI{40}{\volt}$; recovery $\leq\SI{500}{\nano\second}$; \SI{118}{\milli\ampere} peak, \SI{6.6}{\milli\ampere} average (\cref{sec:hv-house-pump}) & Not applicable \\
  \hline
  \des{D}{HVSTART} & Schottky rectifier & To be selected & $V_R\geq\SI{40}{\volt}$; $V_F\leq\SI{0.4}{\volt}$ at \SI{6.6}{\milli\ampere} (\cref{sec:hv-house-pump}) & Not applicable \\
  \hline
  \des{D}{HVZ} & Zener & \SI{30}{\volt} & $\geq\SI{0.5}{\watt}$ (\cref{sec:hv-house-pump}) & Not applicable \\
  \hline
  \des{C}{HVBIAS} & Ceramic & \SI{1}{\micro\farad} & X7R; $\geq\SI{50}{\volt}$; $\geq\SI{0.5}{\micro\farad}$ effective at \SI{21}{\volt} bias (\cref{sec:hv-house-pump}) & Not applicable \\
  \hline
  \des{C}{HVVCC} & Ceramic & \SI{2.2}{\micro\farad} & X7R; \SIrange{1}{4.7}{\micro\farad}; $\geq\SI{16}{\volt}$ (\cref{sec:hv-house-vcc}) & Not applicable \\
  \hline
  \des{C}{HVSS} & Ceramic & \SI{470}{\nano\farad} & X7R; $\pm10\,\%$; $\geq\SI{16}{\volt}$ (\cref{sec:hv-ss}) & Not applicable \\
  \hline
  \des{C}{HVI1} & Polymer aluminum & \SI{100}{\micro\farad} & $\geq\SI{25}{\volt}$; ripple $\geq\SI{0.73}{\ampere}$ at \SI{150}{\kilo\hertz}; ESR $\leq\SI{50}{\milli\ohm}$ (\cref{sec:hv-caps}) & Not applicable \\
  \hline
  \des{C}{HVI2} & Ceramic & \SI{4.7}{\micro\farad} & X7R; $\geq\SI{25}{\volt}$ (\cref{sec:hv-caps}) & Not applicable \\
  \hline
  \des{C}{HVO1} & Polymer aluminum & \SI{4.7}{\micro\farad} & $\geq\SI{250}{\volt}$; ESR $\approx\SI{0.5}{\ohm}$ at \SI{100}{\kilo\hertz} (\cref{sec:hv-caps}) & Not applicable \\
  \hline
  \des{C}{HVO2} & Film & \SI{100}{\nano\farad} & $\geq\SI{250}{\volt}$ (\cref{sec:hv-caps}) & Not applicable \\
  \hline
  \des{R}{HVBLEED} & Fixed resistor & \SI{1.5}{\mega\ohm} & 5\,\%; element voltage $\geq\SI{239}{\volt}$ (\cref{sec:hv-bleeder-discharge}) & DIN0309 axial \\
  \hline
  \des{R}{HVPG} & Fixed resistor & \SI{100}{\kilo\ohm} & 5\,\% (\cref{sec:hv-house-pgood}) & Not applicable \\
\hline
\des{R}{HVINHPU} & Fixed resistor & \SI{100}{\kilo\ohm} & 5\,\% (\cref{sec:hv-inhibit}) & Not applicable \\
\hline
\des{R}{HVINHPD} & Fixed resistor & \SI{1}{\mega\ohm} & 5\,\% (\cref{sec:hv-inhibit}) & Not applicable \\
\hline
\des{R}{HVRELPD} & Fixed resistor & \SI{100}{\kilo\ohm} & 5\,\% (\cref{sec:hv-inhibit}) & Not applicable \\
\hline
\end{longtable}
}

Some parts require more than the nominal value. \des{D}{HVSTART} must be a
Schottky diode. A silicon diode can leave BIAS below its start-up minimum.
\des{C}{HVPUMP}, \des{C}{HVCSF}, and \des{C}{HVCL} require the specified
dielectric. Their capacitance sets charge or timing. \des{R}{HVBLEED} is a
safety-related part. A substitute must meet its element-voltage rating.


\section{5~V and 3.3~V converters: driving component characteristics}
\label{sec:lv-components}

The 5~V rail (for the four K155ID1 BCD decoder/drivers) and the 3.3~V
rail (for the ESP32-S3-WROOM-1 MCU module) are each produced by a
fixed-output synchronous buck built around one IC family: the
AP63205WU (5~V, \textsc{u401}) and the AP63203WU (3.3~V,
\textsc{u411}). Both are named, design-driving parts and are verified
from datasheet here, alongside the MCU module whose supply-current
profile sets the 3.3~V rail's load budget.

\subsection{AP63205WU / AP63203WU synchronous buck converters}
\label{sec:ap632xx}

Source: \texttt{AP632XX.pdf} (Diodes Incorporated, DS41326 Rev.\ 3-2,
Nov.\ 2024). This is a family datasheet covering AP63200, AP63201,
AP63203, AP63205, and AP63212; the ordering-information table (p.~16)
confirms product-version digit ``3'' is the 3.3~V fixed-output part
(AP63203) and digit ``5'' is the 5.0~V fixed-output part (AP63205),
both in the TSOT26 (industry TSOT-23-6) package, both switching at
1.1~MHz. On both fixed variants the FB pin is wired directly to
\(V_{OUT}\); the feedback divider is entirely internal, so the
datasheet's ``\(V_{FB}\)'' accuracy specification \emph{is} the
output-voltage accuracy specification, not a scaled reference.

\begin{longtable}{L{4.6cm}L{3.2cm}L{5.2cm}L{1.2cm}}
\toprule
Parameter & Value & Conditions & Page \\
\midrule
\endhead
Input voltage, absolute maximum & \(-0.3\) to \SI{35}{\volt} (DC); \SI{40}{\volt} (400~ms transient) & Common to family & 4 \\
Input voltage, recommended operating & \SIrange{3.8}{32}{\volt} & \(T_A=-40\) to \SI{85}{\celsius} & 4 \\
Output voltage, AP63203 (3.3~V) & 3.27 / 3.30 / 3.33~V (min/typ/max) & CCM, \(V_{FB}=V_{OUT}\) & 5 \\
Output voltage, AP63205 (5~V) & 4.95 / 5.00 / 5.05~V (min/typ/max) & CCM, \(V_{FB}=V_{OUT}\) & 5 \\
Switching frequency & \SI{1100}{\kilo\hertz} typ, \(\pm6\%\) spread-spectrum dither & Common to family & 5 \\
Minimum on-time & \SI{80}{\nano\second} typ & Common to family & 5 \\
Quiescent current & \SI{22}{\micro\ampere} typ & No load, non-switching, \(V_{EN}\)=open & 5 \\
Shutdown supply current & 1 / 3~\si{\micro\ampere} (typ/max) & \(V_{EN}=0\)~V & 5 \\
HS peak current limit & 2.5 / 2.8 / 3.1~A (min/typ/max) & Cycle-by-cycle, senses high-side FET & 5 \\
LS valley current limit & 2.5 / 3.2 / 3.9~A (min/typ/max) & See note below & 5 \\
EN logic-high threshold & 1.15 / 1.18 / 1.23~V (min/typ/max) & Rising; enables device, starts soft-start & 5 \\
EN logic-low threshold & 1.05 / 1.10 / 1.15~V (min/typ/max) & Falling; disables device & 5 \\
EN input current & 5.5~\si{\micro\ampere} typ @ \(V_{EN}=1.5\)~V; 1.5~\si{\micro\ampere} typ @ \(V_{EN}=1\)~V & Internal 1.5~\si{\micro\ampere} pull-up from internal LDO & 5, 10 \\
Soft-start period & \SI{4}{\milli\second} typ, fixed internal & Not externally programmable & 5, 9, 11 \\
VIN UVLO, rising / falling & 3.50 / \(\approx\)3.06~V typ & Hysteresis \SI{440}{\milli\volt} typ & 5, 11 \\
High-side \(R_{DS(ON)}\) & \SI{125}{\milli\ohm} typ & Common to family & 5 \\
Low-side \(R_{DS(ON)}\) & \SI{68}{\milli\ohm} typ & Common to family & 5 \\
Recommended output capacitor & \SIrange{22}{68}{\micro\farad} ceramic, low ESR & General guidance; BOM tables use \(2\times22\)~\si{\micro\farad} for both variants & 13, 14 \\
Recommended input capacitor & \(>\)\SI{10}{\micro\farad} ceramic & Low ESR; RMS rating \(>\) half max load current; BOM tables use \SI{10}{\micro\farad} & 13 \\
Inductor selection equation & \(L = \dfrac{V_{OUT}(V_{IN}-V_{OUT})}{V_{IN}\,\Delta I_L\,f_{SW}}\) & \(\Delta I_L\) recommended 30--50\% of the 2~A max load & 13 \\
Recommended inductance range & \SIrange{2.2}{10}{\micro\henry} & General guidance & 13 \\
Recommended inductor, AP63203 (3.3~V) BOM table & \SI{3.9}{\micro\henry} & Table 2 & 13 \\
Recommended inductor, AP63205 (5~V) BOM table & \SI{4.7}{\micro\henry} & Table 3 & 13 \\
Inductor saturation-current margin & \(\geq\)35\% above max load current & DCR \(<100\)~m\(\Omega\) recommended for efficiency & 13 \\
Dropout voltage & No numeric spec & Qualitative only: near-100\%-duty ``LDO mode,'' HS FET held on, LS FET pulsed periodically to refresh the BST cap & 14 \\
Thermal shutdown & \SI{160}{\celsius} typ (EC table) & See note on inconsistency below & 5, 11, 12 \\
\(\theta_{JA}\) (TSOT26) & \SI{89}{\celsius\per\watt} & FR-4, single layer, 2~oz Cu, minimum pad & 4 \\
BST capacitor, recommended & \SI{100}{\nano\farad} & Ceramic, SW-to-BST & 2, 14 \\
VSW absolute maximum & \(-1.0\) to \(V_{IN}+0.3\)~V & & 4 \\
VEN absolute maximum & \(-0.3\) to \SI{35}{\volt} & Same as VIN; EN is a high-voltage pin & 4 \\
\bottomrule
\end{longtable}

\paragraph{EN floating or tied to VIN are both valid always-on
configurations,} per the datasheet (an internal 1.5~\si{\micro\ampere}
pull-up auto-enables the part when EN floats, and EN is explicitly
rated as a high-voltage pin safe to tie directly to VIN) --- relevant
because both bucks in this design enable directly from their own input
rail rather than from a sequenced logic signal.

\paragraph{Thermal shutdown threshold is stated three different ways}
across the document: the Electrical Characteristics table gives
\SI{160}{\celsius} typ with \SI{25}{\celsius} hysteresis (implying a
\(\approx\)\SI{135}{\celsius} restart), the narrative Thermal Shutdown
section states \SI{150}{\celsius} shutdown / \SI{130}{\celsius}
restart, and the Power Derating section separately recommends keeping
\(T_J\) below \SI{125}{\celsius} for design margin. The
production-tested table value is treated as authoritative for
compliance; \SI{125}{\celsius} is the more conservative design target
used in \cref{sec:5v-converter,sec:3v3-converter}.

\paragraph{The ``LS valley current limit'' figure (2.5/3.2/3.9~A) is
higher than the ``HS peak current limit'' (2.5/2.8/3.1~A),} which is
unusual naming for a high-side peak-current-mode part: ordinarily a
valley limit sensed on the low-side FET would be expected to trip
\emph{below} the peak limit, not above it. The datasheet gives no
further explanation of what the labeled ``valley'' limit protects
against. \Cref{sec:5v-converter,sec:3v3-converter} therefore treat the
HS peak current limit as the operative cycle-by-cycle limit for margin
calculations, since it is both the lower figure and the one whose
function (peak inductor current limiting) matches the part's stated
control architecture.

\paragraph{No numeric dropout spec.} When \(V_{IN}\) approaches
\(V_{OUT}\), the part enters a qualitatively-described ``LDO mode''
(high-side FET held on nearly continuously, low-side FET pulsed
periodically only to refresh the bootstrap capacitor) rather than
failing to regulate outright. Where dropout behavior matters ---
specifically the 3.3~V rail's alternate UART-supplied input path in
\cref{sec:3v3-converter} --- the resulting output is estimated from
\(R_{DS(ON)}\) and inductor DCR rather than read from a table, since no
such table exists.

\subsection{ESP32-S3-WROOM-1 MCU module}
\label{sec:esp32}

Source: \texttt{esp32-s3-wroom-1\_wroom-1u\_datasheet\_en.pdf}
(Espressif). This is a generic module datasheet covering ten-plus
flash/PSRAM SKUs (N4/N8/N16 flash; no-PSRAM/R2/R8/R16V PSRAM options)
under one part number; the schematic specifies the module generically
with no flash/PSRAM suffix, so the exact installed configuration is
not determinable from the schematic alone.

\begin{longtable}{L{4.8cm}L{2.8cm}L{4.8cm}L{1.2cm}}
\toprule
Parameter & Value & Conditions & Page \\
\midrule
\endhead
VDD33 recommended operating range & \SIrange{3.0}{3.6}{\volt}, 3.3~V typ & & 27 \\
VDD33 absolute maximum & \(-0.3\) to \SI{3.6}{\volt} & Identical to the recommended maximum --- no extra headroom above 3.6~V & 27 \\
External supply current capability required & \(\geq\)\SI{0.5}{\ampere} & Table 6-2, ``current delivered by external power supply'' --- Espressif's own minimum regulator sizing requirement & 27 \\
Wi-Fi TX peak (highest of all published conditions) & \SI{355}{\milli\ampere} & 802.11b, 1~Mbps, @20.5~dBm; rated at 100\% TX duty cycle & 28 \\
Wi-Fi TX, 802.11n HT40 MCS7 @17~dBm & \SI{285}{\milli\ampere} & & 28 \\
Wi-Fi RX & 95 / \SI{97}{\milli\ampere} & HT20 / HT40; peripherals disabled, CPU idle & 28 \\
Bluetooth LE TX peak & \SI{344}{\milli\ampere} & @20.0~dBm (max power) & 28 \\
Bluetooth LE RX & \SI{93}{\milli\ampere} & & 28 \\
Modem-sleep, lowest & \SI{13.2}{\milli\ampere} typ (18.8~mA with peripheral clocks enabled) & WAITI, dual-core idle, 40~MHz & 29 \\
Modem-sleep, highest & \SI{91.7}{\milli\ampere} typ (107.9~mA with peripheral clocks enabled) & Dual-core, 128-bit access, 240~MHz & 29 \\
Light-sleep & \SI{240}{\micro\ampere} typ & Wi-Fi and VDD\_SPI powered down, GPIOs Hi-Z & 30 \\
Deep-sleep (various ULP configurations) & \SIrange{7}{190}{\micro\ampere} typ & See page for the specific ULP/RTC configuration & 30 \\
Power off & \SI{1}{\micro\ampere} typ & EN held low & 30 \\
EN recommended external network & \(R=10\)~k\(\Omega\), \(C=1\)~\si{\micro\farad} & Espressif reference ``Peripheral Schematics'' circuit --- matches this board's R501/C503 exactly & 41 \\
EN reset-release threshold & \(\geq0.75\times V_{DD}\) & \(\approx\)2.475~V at \(V_{DD}=3.3\)~V & 27--28 \\
EN reset threshold & \(\leq0.25\times V_{DD}\) & \(\approx\)0.825~V at \(V_{DD}=3.3\)~V & 27--28 \\
EN/strapping setup and hold time & \(t_{SU}\geq0\), \(t_H\geq3\)~ms & Table 4-2 & 14 \\
GPIO0 (boot strap) default state & Weak pull-up, high = SPI boot (normal) & Internal \(\approx\)45~k\(\Omega\) pull-up; no external pull resistor required; pulling low selects download-mode boot & 13--14, 27 \\
ESD rating & HBM \(\pm\)2000~V, CDM \(\pm\)500~V & & 47 \\
Flash current adder (modem-sleep, 80~Mbit/s, SPI 2-line) & \(+10\)~mA & Added atop the base modem-sleep figures when flash is accessed & 29 \\
PSRAM current adder (light-sleep) & \(+40\) to \(+200\)~\si{\micro\ampere} & Depends on PSRAM size/interface/voltage; no adder is published for active-mode Wi-Fi TX/RX & 30 \\
GPIO \(V_{IH}\) / \(V_{IL}\) & \(\geq0.75\times V_{DD}\) / \(\leq0.25\times V_{DD}\) & Table 6-3, \SI{25}{\celsius} & 27 \\
GPIO \(V_{OH}\) / \(V_{OL}\) & \(\geq0.8\times V_{DD}\) / \(\leq0.1\times V_{DD}\) & High-impedance load & 27 \\
GPIO source/sink current & 40 / \SI{28}{\milli\ampere} typ & \(V_{DD}=3.3\)~V, PAD\_DRIVER=3 & 27 \\
GPIO internal weak pull-up/down & \SI{45}{\kilo\ohm} typ each & & 27 \\
\bottomrule
\end{longtable}

\paragraph{Not specified in this module datasheet.} Brown-out/POR
threshold for VDD33 undervoltage detection, and any GPIO absolute
maximum voltage beyond the recommended \(V_{IH}\) maximum
(\(V_{DD}+0.3\)~V) --- both live in the ESP32-S3 chip-level datasheet,
which this module document references but does not reproduce.

\paragraph{Installed flash/PSRAM configuration is unconfirmed.} The
current-consumption tables explicitly note that PSRAM-equipped SKUs
draw more current than the tabulated baseline: an adder is published
for light-sleep (\(+40\) to \(+200\)~\si{\micro\ampere}) and for
modem-sleep with flash access (\(+10\)~mA), but \emph{no} adder is
given for active-mode Wi-Fi TX/RX --- so if the installed module has
PSRAM, the 285--355~mA active-mode figures above may understate real
draw by an amount Espressif does not quantify. \Cref{sec:3v3-converter}
treats the published figures as the design point and flags this as an
open, unquantified margin item rather than papering over it.

\paragraph{EN network matches Espressif's own reference exactly.} The
board's populated R501 (10~k\(\Omega\)) and C503 (1~\si{\micro\farad})
on the ESP32's EN pin are identical to the values in Espressif's own
recommended application circuit --- confirmed, not merely
plausible.

\paragraph{The regulator must tolerate the full Wi-Fi TX transient, not
just steady-state draw.} Espressif's own stated minimum external
supply-current capability (\(\geq\)0.5~A) is driven directly by the
355~mA Wi-Fi TX peak plus margin; a buck sized only for the module's
idle/modem-sleep current (tens of mA) would brown out during a
transmit burst. \Cref{sec:3v3-converter} sizes the 3.3~V converter
against this transient requirement, not against an averaged current
figure.

\section{5~V converter derivation}\label{sec:5v-converter}

This section derives the required characteristics of the 5~V rail and verifies the
populated Rev-A implementation: U401 (AP63205WU, \cref{sec:ap632xx}) with L401
(\SI{4.7}{\micro\henry}), C401 (\SI{10}{\micro\farad} input), C402 (\SI{100}{\nano\farad}
BST), and C403/C404 ($2\times\SI{22}{\micro\farad}$ output). The exported netlist confirms
the \texttt{+5V} net contains exactly: the converter output components (L401, C403, C404),
the four K155ID1 $V_{CC}$ pins (U1--U4, pin 5), their per-package bypass capacitors
C1--C4 (\SI{100}{\nano\farad} each), and the power-indicator branch R401
(\SI{5.1}{\kilo\ohm}) feeding D401 to ground --- nothing else. The rail powers the BCD
decoder/drivers and its own indicator LED alone; it has no connection to the MCU or any
other subsystem. This is a
far lighter design point than the flyback of \cref{sec:hv-converter}: a fixed-output
integrated buck, a static logic load two decades below the converter's rating, and no
external compensation, sensing, or clamp networks to verify. The derivation is
correspondingly shorter; where a margin is simply large, it is reported as large and left
alone.

\subsection{Load budget}\label{sec:5v-load}

Each K155ID1 draws $I_{CC} = 16$~mA typical, \SI{25}{\milli\ampere} maximum
(\cref{sec:bcd-driver}; specified at $V_{CC}=$ max, inputs grounded, outputs open). The
cathode current sunk by the selected output does \emph{not} add to this: the driver output
is an open-collector TTL stage sinking current that arrives from the 170~V rail through
the tube, into the output pin, and out the ground pin --- it never passes through
$V_{CC}$. The rail load is therefore the four packages' logic supply current only:
\begin{align*}
  I_{\mathrm{load,typ}} &= 4 \times \SI{16}{\milli\ampere} = \SI{64}{\milli\ampere}, &
  I_{\mathrm{load,max}} &= 4 \times \SI{25}{\milli\ampere} = \SI{100}{\milli\ampere},
\end{align*}
\begin{align*}
  P_{\mathrm{out,typ}} &= 5 \times 0.064 = \SI{0.32}{\watt}, &
  P_{\mathrm{out,max}} &= 5 \times 0.100 = \SI{0.50}{\watt}.
\end{align*}
Bipolar TTL supply current is set by the internal bias networks, not by output state, so
this load is essentially static: it does not change with which digit is lit, and the only
input transitions are the once-per-second (and slower) BCD updates. There is no transient
load profile to size against --- the DC maximum is the whole story. The
\SI{100}{\nano\farad} per-package bypass capacitors C1--C4 are the standard one-per-TTL-package
practice, generously adequate at these switching rates, and contribute no DC load.

One further branch loads the rail continuously: the power-indicator LED D401, fed
through R401 (\SI{5.1}{\kilo\ohm}). \emph{Assumption:} forward voltage
$\approx\SI{2.0}{\volt}$ (generic 0805 part, color unrecorded; the value field reads
``+5V,'' a functional label rather than a part choice). Then
\begin{align*}
  I_{\mathrm{D401}} &= \frac{5-2.0}{\SI{5.1}{\kilo\ohm}} = \SI{0.59}{\milli\ampere}
    \quad\text{(continuous, non-switching)},
\end{align*}
raising the rail totals to \SI{65}{\milli\ampere} typical / \SI{101}{\milli\ampere}
maximum. The increment is 0.6\,\% of the driver load --- below the rounding of every
downstream figure and well inside the $\eta=0.85$ conservatism --- so the derivations
that follow carry the driver-only \SI{100}{\milli\ampere} bound unchanged.

\subsection{Operating point and conduction mode}\label{sec:5v-op}

The input is $V_{\mathrm{BUS}}$, the USB-PD-negotiated rail (9 or \SI{12}{\volt} nominal
contracts; analysis window \SIrange{8.5}{13}{\volt}), tied directly to both IN and EN with
no series element. The window sits deep inside the AP63205's recommended
\SIrange{3.8}{32}{\volt} operating range (\cref{sec:ap632xx}), with \SI{22}{\volt} to the
\SI{35}{\volt} absolute maximum. The ideal buck duty cycle $D = V_{OUT}/V_{IN}$ spans
\begin{align*}
  D_{\min} &= \frac{5}{13} = 0.385 \quad (V_{IN}=\SI{13}{\volt}), &
  D_{\max} &= \frac{5}{8.5} = 0.588 \quad (V_{IN}=\SI{8.5}{\volt}).
\end{align*}
Conduction-loss corrections to $D$ are of order $I_{\mathrm{load}} R_{DS(\mathrm{on})}
\approx \SI{13}{\milli\volt}$ against a 5~V output --- below 0.3\,\% --- and are ignored
throughout. Because $D_{\max} > 0.5$ at the bottom of the window, slope compensation is
required for a peak-current-mode part; the AP63205 integrates it internally and no
external action exists to take (\ds{AP632XX.pdf}{9}). Dropout is not remotely in play:
minimum headroom is $8.5 - 5 = \SI{3.5}{\volt}$, so the near-100\%-duty ``LDO mode''
described qualitatively in \cref{sec:ap632xx} is never entered on this rail.

\paragraph{Conduction mode.} With the populated $L = \SI{4.7}{\micro\henry}$ at
$f_{SW} = \SI{1.1}{\mega\hertz}$, the inductor ripple current in fixed-frequency
operation would be
\begin{align*}
  \Delta I_L \big|_{\SI{8.5}{\volt}} &=
    \frac{5\,(8.5-5)}{8.5 \times \SI{4.7}{\micro\henry} \times \SI{1.1}{\mega\hertz}}
    = \frac{17.5}{43.9} = \SI{0.398}{\ampere}, \\
  \Delta I_L \big|_{\SI{13}{\volt}} &=
    \frac{5\,(13-5)}{13 \times \SI{4.7}{\micro\henry} \times \SI{1.1}{\mega\hertz}}
    = \frac{40}{67.2} = \SI{0.595}{\ampere},
\end{align*}
so the CCM boundary load ($\Delta I_L/2$) is \SIrange{199}{298}{\milli\ampere} across the
input window. The actual load, \SIrange{64}{100}{\milli\ampere}, sits below the boundary
at every corner: the converter does not run continuous fixed-frequency PWM here. This is
not the design concern it was for the flyback, where DCM was a \emph{designed-in}
requirement that the compensation, current limit, and transformer reset all depended on
(\cref{sec:hv-dcm}). In the AP63205 the light-load behavior is an automatic, internally
managed PFM mode: the COMP node is clamped at a floor corresponding to a
\SI{450}{\milli\ampere} peak inductor current, zero-cross detection on the low-side FET
blocks reverse current, and the part skips cycles as the load falls, maintaining
regulation continuously through the transition (\ds{AP632XX.pdf}{10}). No component
choice in this design selects or depends on the mode; it is reported here as an operating
fact, not verified as a requirement.

The PFM operating point follows directly from the \SI{450}{\milli\ampere} clamp.
Each pulse ramps the inductor to $I_{\mathrm{pk}} = \SI{0.45}{\ampere}$ and back to zero:
\begin{align*}
  t_{\mathrm{on}} &= \frac{L\,I_{\mathrm{pk}}}{V_{IN}-V_{OUT}}
    = \frac{\SI{2.115}{\micro\volt\second}}{\SI{8}{\volt}} = \SI{264}{\nano\second}
    \;\; (V_{IN}=\SI{13}{\volt}); \quad
    \frac{\SI{2.115}{\micro\volt\second}}{\SI{3.5}{\volt}} = \SI{604}{\nano\second}
    \;\; (V_{IN}=\SI{8.5}{\volt}), \\
  t_{\mathrm{dem}} &= \frac{L\,I_{\mathrm{pk}}}{V_{OUT}}
    = \frac{\SI{2.115}{\micro\volt\second}}{\SI{5}{\volt}} = \SI{423}{\nano\second}
    \quad (\text{both corners}),
\end{align*}
delivering charge $Q = \tfrac{1}{2} I_{\mathrm{pk}} (t_{\mathrm{on}}+t_{\mathrm{dem}})
= \SI{155}{\nano\coulomb}$ (\SI{13}{\volt}) to \SI{231}{\nano\coulomb}
(\SI{8.5}{\volt}) per pulse. The average pulse repetition rate needed to support the load
is $f_p = I_{\mathrm{load}}/Q$:
\begin{align*}
  f_p &= \frac{0.100}{155\,\si{\nano\coulomb}} \approx \SI{650}{\kilo\hertz}
    \;\;(\SI{100}{\milli\ampere}, \SI{13}{\volt})
  &&\text{down to}&
  f_p &= \frac{0.064}{231\,\si{\nano\coulomb}} \approx \SI{277}{\kilo\hertz}
    \;\;(\SI{64}{\milli\ampere}, \SI{8.5}{\volt}),
\end{align*}
i.e., roughly a quarter to two-thirds of the nominal clock, well above audibility and an
idealized average (real skipping is bursty). The datasheet's typical-efficiency curve
($V_{IN}=\SI{12}{\volt}$, $V_{OUT}=\SI{5}{\volt}$) reads approximately 88--90\,\% in the
60--100~mA region (\ds{AP632XX.pdf}{2}); $\eta = 0.85$ is carried below as a conservative
figure. Average input current is then at most
\begin{align*}
  I_{IN} &= \frac{P_{\mathrm{out,max}}}{\eta\,V_{IN}}
          = \frac{0.50}{0.85 \times 8.5} = \SI{69}{\milli\ampere},
\end{align*}
a negligible increment on the $V_{\mathrm{BUS}}$ budget next to the flyback's draw.

\subsection{Minimum on-time}\label{sec:5v-ontime}

For a buck the on-time shrinks as $V_{IN}$ rises, so the binding corner is
$V_{IN}=\SI{13}{\volt}$ --- the opposite corner from the flyback's. In hypothetical
fixed-frequency operation at the +6\,\% spread-spectrum extreme
($f_{SW} = \SI{1.166}{\mega\hertz}$):
\begin{align*}
  t_{\mathrm{on}} &= \frac{D_{\min}}{f_{SW}}
    = \frac{0.385}{\SI{1.166}{\mega\hertz}} = \SI{330}{\nano\second}
    \quad\text{vs.}\quad t_{\mathrm{on,min}} = \SI{80}{\nano\second}
    \;\;(\cref{sec:ap632xx}),
\end{align*}
a $4.1\times$ margin. In the PFM mode actually occupied, the shortest pulse is the
\SI{264}{\nano\second} on-time computed above --- still $3.3\times$ the minimum. For the
constraint ever to bind at this output voltage, the input would need
$D = t_{\mathrm{on,min}} f_{SW} = 0.093$, i.e.\ $V_{IN} = 5/0.093 \approx \SI{54}{\volt}$,
beyond even the \SI{35}{\volt} absolute maximum. Minimum on-time cannot constrain a 5~V
output anywhere in this part's operating space. Match, no further analysis warranted.

\subsection{Inductor}\label{sec:5v-inductor}

The datasheet's selection equation (\cref{sec:ap632xx}) with its prescribed ripple target
--- $\Delta I_L$ equal to 30--50\,\% of the \emph{2~A rated} load, i.e.\
\SIrange{0.6}{1.0}{\ampere}; the reference is the part's rating, not this application's
100~mA --- gives, at the maximum-ripple corner $V_{IN}=\SI{13}{\volt}$:
\begin{align*}
  L\big|_{\Delta I_L = \SI{1.0}{\ampere}} &=
    \frac{5\,(13-5)}{13 \times 1.0 \times \SI{1.1}{\mega\hertz}} = \SI{2.80}{\micro\henry},
  &
  L\big|_{\Delta I_L = \SI{0.6}{\ampere}} &=
    \frac{40}{13 \times 0.6 \times \SI{1.1}{\mega\hertz}} = \SI{4.66}{\micro\henry},
\end{align*}
and \SIrange{1.87}{3.12}{\micro\henry} at \SI{8.5}{\volt}. The populated
\SI{4.7}{\micro\henry} sits at the large-inductance edge of the 13~V window (ripple
29.8\,\% of 2~A at \SI{13}{\volt}, 19.9\,\% at \SI{8.5}{\volt} --- at or slightly below
the band's 30\,\% floor, the benign direction, which the datasheet itself endorses:
``use a larger inductance for improved efficiency under light load,'' precisely this
rail's condition). More directly: \SI{4.7}{\micro\henry} is the datasheet's own Table-3
BOM value for the AP63205 5~V output (\cref{sec:ap632xx}), and the surrounding parts ---
\SI{10}{\micro\farad} input, $2\times\SI{22}{\micro\farad}$ output,
\SI{100}{\nano\farad} BST --- reproduce that table line for line. The populated design is
the manufacturer's reference BOM applied verbatim, not an independent coincidence.
It also sits inside the general \SIrange{2.2}{10}{\micro\henry} recommended range. Match.

\paragraph{Saturation current.} The schematic records only the footprint
(\texttt{L\_APV\_ANR6045}, a $6.0\times6.0\times\SI{4.5}{\milli\metre}$ shielded power
inductor), not a part number, so the rating is checked against the class, with the
assumption stated: \emph{assume} $I_{\mathrm{sat}} \geq \SI{3.0}{\ampere}$ and DCR
$\leq \SI{50}{\milli\ohm}$, typical of \SI{4.7}{\micro\henry} parts in this size class
(commodity 6045-package \SI{4.7}{\micro\henry} shielded inductors commonly rate
3.5--5~A saturation, 20--50~m$\Omega$ DCR). Against the datasheet's application rule
($\geq$35\,\% above maximum load: $1.35 \times 0.100 = \SI{135}{\milli\ampere}$) the
assumed rating is $>20\times$; against the largest current the inductor ever carries in
normal operation --- the \SI{450}{\milli\ampere} PFM peak clamp --- it is $\approx 7\times$;
and it covers even the worst-case fault current, the \SI{3.1}{\ampere} maximum HS peak
limit, at or near its hard-saturation rating, with the converter's hiccup protection
(current-limit events for 2~ms force a 16~ms shutdown, \ds{AP632XX.pdf}{11}) bounding any
fault dwell time. The DCR assumption satisfies the $<\SI{100}{\milli\ohm}$ efficiency
guidance. Match under the stated assumption; the actual part's saturation rating should
be confirmed when the BOM is finalized, the same status as the unrecorded capacitor
series in \cref{sec:hv-caps}.

\subsection{Current-limit margin}\label{sec:5v-ilim}

The operative cycle-by-cycle limit is the HS peak current limit, \SI{2.5}{\ampere}
minimum (\cref{sec:ap632xx}, which also records why the higher ``LS valley'' figure is
not used). The largest peak inductor current the control loop commands at this load is
the \SI{450}{\milli\ampere} PFM clamp:
\begin{align*}
  \frac{I_{\mathrm{lim}}^{\min}}{I_{\mathrm{pk}}} &= \frac{2.5}{0.45} = 5.6\times.
\end{align*}
Even if the load quadrupled and the converter re-entered CCM, the peak
($I_{\mathrm{load}} + \Delta I_L/2 \approx \SI{0.7}{\ampere}$ at \SI{400}{\milli\ampere}
load) would still leave $>3.5\times$. The current limit exists on this rail purely as
fault protection --- a shorted +5~V net lands in hiccup mode rather than anything
thermal --- and no operating condition approaches it. Large margin; nothing further to
derive.

\subsection{Output capacitance and rail accuracy}\label{sec:5v-cout}

The populated $2\times\SI{22}{\micro\farad}$ (C403/C404) equals the datasheet BOM example
and its \SI{44}{\micro\farad} nominal total sits inside the \SIrange{22}{68}{\micro\farad}
recommendation (\cref{sec:ap632xx}). Assumption: the parts are X5R/X7R ceramics rated
$\geq \SI{10}{\volt}$, retaining $\geq 60\,\%$ of nominal capacitance at 5~V DC bias ---
effective $C_{\mathrm{out}} \geq \SI{26}{\micro\farad}$, still above the recommendation's
floor. The datasheet's load-transient sizing bound (\ds{AP632XX.pdf}{14}), with a
permitted overshoot $\Delta V = \SI{50}{\milli\volt}$ and the worst CCM-formula peak:
\begin{align*}
  C_{\mathrm{out}} &> \frac{L\left(I_{\mathrm{out}}+\Delta I_L/2\right)^2}
                          {(\Delta V + V_{OUT})^2 - V_{OUT}^2}
    = \frac{\SI{4.7}{\micro\henry} \times (0.100+0.298)^2\,\si{\square\ampere}}{5.05^2 - 5^2\ \si{\square\volt}}
    = \frac{\SI{0.744}{\micro\joule}}{\SI{0.503}{\square\volt}}
    = \SI{1.5}{\micro\farad},
\end{align*}
met with $17\times$ margin even after derating --- unsurprising, since the load has no
transient profile to speak of. Output ripple in the occupied PFM mode is bounded by one
pulse's full charge dumped into the (derated) capacitance plus the ESR step:
\begin{align*}
  \Delta V_{\mathrm{ripple}} &\lesssim \frac{Q_{\max}}{C_{\mathrm{out}}}
      + I_{\mathrm{pk}}\,\mathrm{ESR}
    = \frac{\SI{231}{\nano\coulomb}}{\SI{26}{\micro\farad}}
      + 0.45 \times \SI{2.5}{\milli\ohm}
    \approx 8.9 + 1.1 \approx \SI{10}{\milli\volt}\ \text{pk-pk},
\end{align*}
(assumed effective ESR \SI{2.5}{\milli\ohm} for two paralleled ceramics; the
fixed-frequency ripple, $\Delta I_L/(8 f_{SW} C) + \Delta I_L\,\mathrm{ESR} \approx
\SI{4}{\milli\volt}$, is smaller still). Stacking the worst case against the load's
requirement: the AP63205 regulates to \SIrange{4.95}{5.05}{\volt} over conditions
(\cref{sec:ap632xx}), and $4.95 - 0.010 = \SI{4.94}{\volt}$ /
$5.05 + 0.010 = \SI{5.06}{\volt}$ sit inside the K155ID1's recommended
\SIrange{4.75}{5.25}{\volt} supply window (\cref{sec:bcd-driver}) with
$\approx\SI{190}{\milli\volt}$ margin on both sides. Match.

\subsection{Input capacitance}\label{sec:5v-cin}

The populated C401 (\SI{10}{\micro\farad} ceramic) equals the datasheet's BOM-table
example; the accompanying prose asks for ``greater than \SI{10}{\micro\farad}''
(\cref{sec:ap632xx}), so the population meets the manufacturer's own example rather than
strictly exceeding the prose figure --- noted for precision, not flagged, because the
requirement scales with load and this rail runs at 1/20th of the part's rating. The
buck's input RMS ripple current at the worst point ($D$ nearest 0.5, i.e.\
$V_{IN}=\SI{8.5}{\volt}$):
\begin{align*}
  I_{\mathrm{RMS}} &\approx I_{\mathrm{out}}\sqrt{D(1-D)}
    = 0.100\sqrt{0.588 \times 0.412} = \SI{49}{\milli\ampere},
\end{align*}
against the guidance that the capacitor's RMS rating exceed half the maximum load
(\SI{50}{\milli\ampere} for this application; ampere-class for any ceramic
\SI{10}{\micro\farad}) --- margin of order $50\times$. Input ripple voltage at nominal
capacitance is $I_{\mathrm{out}} D(1-D)/(f_{SW} C) \approx \SI{2}{\milli\volt}$, still
$\approx\SI{3}{\milli\volt}$ after DC-bias derating at 13~V. The one thing the schematic
does not record is the capacitor's voltage rating: on a 13~V rail it must be
$\geq\SI{16}{\volt}$, and \SI{25}{\volt} is preferred for bias-derating headroom ---
confirm at BOM finalization, same status as the inductor rating above. C402
(\SI{100}{\nano\farad}, SW-to-BST) matches the recommended bootstrap value exactly.
Match throughout.

\subsection{Soft start, enable, and dissipation}\label{sec:5v-ss}

Soft start is the fixed internal \SI{4}{\milli\second} ramp, not externally programmable
(\cref{sec:ap632xx}); there is no component to verify. Start-up charging current is
correspondingly trivial:
\begin{align*}
  I_{\mathrm{chg}} &= C_{\mathrm{out}}\frac{dV}{dt}
    = \SI{44}{\micro\farad} \times \frac{\SI{5}{\volt}}{\SI{4}{\milli\second}}
    = \SI{55}{\milli\ampere},
\end{align*}
which, even summed with the full \SI{100}{\milli\ampere} load, stays a factor of 16 below
the minimum current limit --- no inrush interaction with the limit or with hiccup mode.
EN tied directly to $V_{\mathrm{BUS}}$ is one of the two sanctioned always-on
configurations (\cref{sec:ap632xx}): the converter starts as soon as the upstream load
switch presents $V_{\mathrm{BUS}}$, with the drivers' 5~V arriving about
\SI{4}{\milli\second} later, and no sequencing constraint exists among this rail, the
3.3~V rail, or the flyback at this level. Cross-rail power-up ordering, to the extent it
matters at all, is a system-level topic taken up after the individual rails are
established, not a property of this converter.

Dissipation, using the conservative $\eta = 0.85$ from \cref{sec:5v-op}:
\begin{align*}
  P_{\mathrm{diss}} &= P_{\mathrm{out,max}}\left(\frac{1}{\eta}-1\right)
    = 0.50 \times 0.176 = \SI{88}{\milli\watt}, &
  \Delta T &= \theta_{JA} P_{\mathrm{diss}} = 89 \times 0.088 \approx \SI{8}{\kelvin},
\end{align*}
so $T_J \approx \SI{48}{\celsius}$ at a \SI{40}{\celsius} ambient --- far below the
\SI{125}{\celsius} design target adopted in \cref{sec:ap632xx}. No thermal management
beyond the footprint's own copper is required.

\subsection{Summary}\label{sec:5v-summary}

Assumptions carried through this section: inductor $I_{\mathrm{sat}} \geq \SI{3.0}{\ampere}$
and DCR $\leq \SI{50}{\milli\ohm}$ (generic 6045-class part, no part number recorded);
output ceramics retain $\geq 60\,\%$ capacitance at 5~V bias with effective ESR
\SI{2.5}{\milli\ohm}; converter efficiency 0.85 (conservative floor under the datasheet's
$\approx$0.88--0.90 typical curve); indicator-LED forward voltage \SI{2.0}{\volt};
PFM pulses idealized as clean
\SI{450}{\milli\ampere}-peak triangles for the ripple and pulse-rate estimates. Each
appears at its point of use above.

\begin{longtable}{@{}p{0.26\textwidth}p{0.30\textwidth}p{0.22\textwidth}p{0.14\textwidth}@{}}
\toprule
Parameter & Derived requirement & Populated & Verdict \\
\midrule
\endfirsthead
\toprule
Parameter & Derived requirement & Populated & Verdict \\
\midrule
\endhead
\bottomrule
\endlastfoot
Load current & $4\times I_{CC}$ (64/100~mA) $+$ D401 indicator
  (\SI{0.59}{\milli\ampere} via R401): \SI{65}{\milli\ampere} typ /
  \SI{101}{\milli\ampere} max, static & --- & Derived \\
Input window & inside \SIrange{3.8}{32}{\volt} recommended &
  $V_{\mathrm{BUS}}$ \SIrange{8.5}{13}{\volt}, direct & Match (large margin) \\
Duty cycle & 0.385--0.588; internal slope compensation covers $D>0.5$ & --- &
  Verified \\
Conduction mode & load $<$ CCM boundary (\SIrange{199}{298}{\milli\ampere})
  $\rightarrow$ automatic PFM; regulation maintained & --- & Verified (benign) \\
Minimum on-time & \SI{264}{\nano\second} (PFM) / \SI{330}{\nano\second} (PWM,
  $+6\%$ dither) vs.\ \SI{80}{\nano\second} & --- & Match ($\geq 3.3\times$) \\
Inductance L401 & Eq.-7 window \SIrange{2.8}{4.7}{\micro\henry} at \SI{13}{\volt};
  datasheet 5~V BOM value \SI{4.7}{\micro\henry} & \SI{4.7}{\micro\henry}
  & Match (reference BOM followed) \\
Inductor $I_{\mathrm{sat}}$ & $\geq\SI{0.45}{\ampere}$ normal;
  $\geq\SI{3.1}{\ampere}$ covers fault peak & 6045-class, rating unrecorded
  (assumed $\geq\SI{3}{\ampere}$) & Match under assumption; \textbf{confirm part} \\
Current-limit margin & peak $\leq\SI{0.45}{\ampere}$ vs.\ \SI{2.5}{\ampere} min
  limit & --- & Large margin ($5.6\times$) \\
Output capacitance C403/C404 & \SIrange{22}{68}{\micro\farad}
  ($\geq\SI{1.5}{\micro\farad}$ transient bound) & $2\times\SI{22}{\micro\farad}$
  (= BOM example) & Match \\
Output ripple & $\lesssim\SI{10}{\milli\volt}$ pk-pk (PFM worst case) & --- &
  Derived; $\ll$ load tolerance \\
Rail accuracy at load & K155ID1 requires \SIrange{4.75}{5.25}{\volt} &
  \SIrange{4.94}{5.06}{\volt} incl.\ ripple & Match ($\approx\SI{190}{\milli\volt}$
  both sides) \\
Input capacitor C401 & $\geq\SI{10}{\micro\farad}$ ceramic; RMS
  $\geq\SI{49}{\milli\ampere}$ & \SI{10}{\micro\farad} (= BOM example; rating
  unrecorded) & Match; \textbf{confirm $\geq$16~V rating} \\
BST capacitor C402 & \SI{100}{\nano\farad} SW-to-BST & \SI{100}{\nano\farad} &
  Match \\
Soft start & \SI{4}{\milli\second} fixed internal; inrush
  $\leq\SI{155}{\milli\ampere}$ & internal, no component & Match \\
EN configuration & EN\,=\,VIN valid always-on & EN tied to $V_{\mathrm{BUS}}$ &
  Match \\
Converter dissipation & \SI{88}{\milli\watt}, $\Delta T \approx \SI{8}{\kelvin}$
  vs.\ \SI{125}{\celsius} target & TSOT26, $\theta_{JA}=89$~\si{\celsius\per\watt} &
  Match (large margin) \\
Bypass C1--C4 & \SI{100}{\nano\farad} per TTL package, standard practice &
  \SI{100}{\nano\farad} $\times 4$ & Match \\
\end{longtable}

No populated value is flagged as a mismatch. The two open items are both unrecorded
ratings rather than wrong values: the actual L401 part's saturation current and the C401
voltage rating ($\geq\SI{16}{\volt}$ on a 13~V rail), both to be confirmed when the BOM
is finalized. The rail otherwise reduces to the manufacturer's reference design operating
at one-twentieth of its rating into a static load --- every margin is wide, and the
section's only substantive finding is the (benign, automatic) PFM operating mode.

\section{3.3~V converter derivation}\label{sec:3v3-converter}

This section derives the required characteristics of the 3.3~V rail and verifies the
populated Rev-A implementation: U411 (AP63203WU, \cref{sec:ap632xx}) with L411
(\SI{4.7}{\micro\henry}), C411 (\SI{10}{\micro\farad} input), C412 (\SI{100}{\nano\farad}
BST), and C413/C414 ($2\times\SI{22}{\micro\farad}$ output). The exported netlist confirms
the \texttt{+3V3} net contains exactly: the converter output components (L411, C413,
C414), the ESP32-S3-WROOM-1 module's supply pin (U501 pin 2) with its local decoupling
C501/C502, U411's FB pin (tied to the output, as on all fixed variants), and five
resistive sources feeding other nets --- R206 (LM5155 PGOOD pull-up), R501 (ESP32 EN
network), R502/R503 (UART activity-LED feeds), and R411 (rail indicator LED). The
converter is the same IC family, footprint set, and support-component topology as the 5~V
rail of \cref{sec:5v-converter}, and the generic buck method established there ---
selection equations, mode-boundary logic, transient bound, RMS estimates --- is applied
here without re-derivation. What is genuinely different on this rail, and where this
section spends its effort, is threefold: the load is large and transient-heavy (Wi-Fi TX
bursts, not static logic current), the converter's input is a two-diode OR of
$V_{\mathrm{BUS}}$ and an externally supplied UART-header voltage rather than a direct
rail connection, and the rail sources a status pull-up (PGOOD) whose behavior must be
consistent with the rail's own sequencing.

\subsection{Load budget}\label{sec:3v3-load}

The dominant load is the MCU module. From \cref{sec:esp32}: the worst published sustained
condition is Wi-Fi TX at \SI{355}{\milli\ampere} (802.11b, \SI{20.5}{dBm}, rated at
100\,\% TX duty cycle --- a sustained worst case, not an instantaneous spike atop a
sustained figure), and Espressif's own regulator-sizing requirement is an external supply
capable of $\geq\SI{0.5}{\ampere}$. Typical operation sits far lower: modem-sleep spans
\SIrange{13.2}{91.7}{\milli\ampere} depending on CPU state and clock, with a
\SI{10}{\milli\ampere} adder during flash access; light- and deep-sleep are in the
\si{\micro\ampere} range. The remaining loads are checked for negligibility rather than
derived at length:

\begin{itemize}
\item \emph{BCD drive lines.} Sixteen GPIOs drive the four K155ID1s' BCD inputs. A TTL
input driven high draws $I_{IH} \leq \SI{40}{\micro\ampere}$ (A input) or
\SI{80}{\micro\ampere} (B/C/D inputs) from the driving pin (\cref{sec:bcd-driver}); a
line driven low \emph{sinks} current sourced from the drivers' 5~V rail and adds nothing
to VDD33. Worst case, all sixteen lines high:
\begin{align*}
  I_{\mathrm{BCD}} &= 4 \times (40 + 3\times80)\,\si{\micro\ampere}
    = \SI{1.12}{\milli\ampere},
\end{align*}
0.3\,\% of the Wi-Fi figure. Negligible, as expected.
\item \emph{Indicator LEDs.} R411 (\SI{1.5}{\kilo\ohm}) feeds D413, the rail's power
indicator, continuously; R502/R503 (\SI{1.5}{\kilo\ohm} each) feed D501/D502, activity
LEDs whose cathodes sit on the RXD0/TXD0 nets and which conduct only while the
corresponding UART line is driven low (the UART idle state is high, so idle draw is
zero). \emph{Assumption:} LED forward voltage $\approx\SI{2.0}{\volt}$ (generic 0805
part, color unrecorded). Then
\begin{align*}
  I_{\mathrm{D413}} &= \frac{3.3-2.0}{\SI{1.5}{\kilo\ohm}} = \SI{0.87}{\milli\ampere}
    \quad\text{(continuous)}, \\
  I_{\mathrm{D501,D502}} &\approx \frac{3.3-2.0-0.1}{\SI{1.5}{\kilo\ohm}}
    = \SI{0.8}{\milli\ampere} \quad\text{(per line, while low).}
\end{align*}
(D413's value field reads ``RX,'' evidently copied from the activity LEDs; electrically
it is the rail indicator --- a cosmetic schematic-field inconsistency, flagged for BOM
hygiene only.)
\item \emph{Pull-ups.} R206 (PGOOD) contributes at most \SI{33}{\micro\ampere} and only
during power-not-good windows; R501 (EN) conducts only transiently. Both are quantified
in \cref{sec:3v3-house}.
\end{itemize}

The design totals are therefore
\begin{align*}
  I_{\mathrm{load,max}} &= 355 + 1.12 + 0.87 + 2\times0.8 + 0.03
    \approx \SI{359}{\milli\ampere} \approx \SI{0.36}{\ampere},\\
  P_{\mathrm{out,max}} &= 3.3 \times 0.359 = \SI{1.19}{\watt},
\end{align*}
the sustained worst case, with Espressif's $\geq\SI{0.5}{\ampere}$ capability requirement adopted as the
converter-sizing floor (margin checks below are run at both \SI{0.36}{\ampere} and
\SI{0.5}{\ampere}). Against the AP63203's \SI{2}{\ampere} rating the capability
requirement is met $4\times$ over. For the UART-powered scenario of
\cref{sec:3v3-diodeor} --- an MCU in download/console mode, radio off --- the applicable
draw is the modem-sleep band plus flash access and LEDs: roughly
\SIrange{30}{100}{\milli\ampere}, with \SI{50}{\milli\ampere} carried as the nominal
estimate. One open item is inherited from \cref{sec:esp32} and not resolvable here:
if the installed module SKU has PSRAM, active-mode draw may exceed the tabulated
\SI{355}{\milli\ampere} by an amount Espressif does not quantify. The $4\times$ headroom
between the \SI{0.5}{\ampere} requirement and the \SI{2}{\ampere} part absorbs any
plausible adder, but the published figures remain the design point.

\subsection{Operating point and conduction mode}\label{sec:3v3-op}

Unlike U401, whose IN/EN tie directly to $V_{\mathrm{BUS}}$, U411's IN and EN share a
common node fed through Schottky diode D411 from $V_{\mathrm{BUS}}$ (the OR topology
itself is treated in \cref{sec:3v3-diodeor}; this subsection covers the normal,
$V_{\mathrm{BUS}}$-fed case). \emph{Assumption} (generic SOD-123 Schottky, no part
number recorded): $V_F \approx \SI{0.3}{\volt}$ at \SI{50}{\milli\ampere},
\SI{0.4}{\volt} at \SIrange{200}{300}{\milli\ampere}, \SI{0.45}{\volt} at
\SI{0.5}{\ampere}; reverse leakage $\leq\SI{100}{\micro\ampere}$ warm at \SI{13}{\volt}.
The converter's input node then spans
\begin{align*}
  V_N^{\min} &= 8.5 - 0.4 = \SI{8.1}{\volt} \quad\text{(full load, minimum bus)}, \\
  V_N^{\max} &= 13 - 0.2 \approx \SI{12.8}{\volt} \quad\text{(light load, maximum bus)},
\end{align*}
inside the \SIrange{3.8}{32}{\volt} recommended range with the same style of margin as
the 5~V rail; EN\,=\,IN is a sanctioned always-on configuration, and both pins' absolute
maximum is \SI{35}{\volt} (\cref{sec:ap632xx}) --- a $2.7\times$ margin over
$V_N^{\max}$. The ideal duty cycle spans
\begin{align*}
  D_{\min} &= \frac{3.3}{12.8} = 0.258, &
  D_{\max} &= \frac{3.3}{8.1} = 0.407,
\end{align*}
never crossing 0.5, so even the internal slope compensation is unstressed. Conduction
corrections ($I R_{DS(\mathrm{on})} \approx \SI{45}{\milli\volt}$ at full load, 1.4\,\%
of $V_{OUT}$) are ignored as in \cref{sec:5v-op}.

\paragraph{Conduction mode.} With the populated \SI{4.7}{\micro\henry} at
$f_{SW}=\SI{1.1}{\mega\hertz}$, the fixed-frequency ripple current is
\begin{align*}
  \Delta I_L \big|_{\SI{12.8}{\volt}} &=
    \frac{3.3\,(12.8-3.3)}{12.8 \times \SI{4.7}{\micro\henry} \times \SI{1.1}{\mega\hertz}}
    = \frac{31.35}{66.2} = \SI{0.474}{\ampere}, \\
  \Delta I_L \big|_{\SI{8.1}{\volt}} &=
    \frac{3.3\,(8.1-3.3)}{8.1 \times \SI{4.7}{\micro\henry} \times \SI{1.1}{\mega\hertz}}
    = \frac{15.84}{41.9} = \SI{0.378}{\ampere},
\end{align*}
so the CCM boundary ($\Delta I_L/2$) is \SIrange{189}{237}{\milli\ampere} across the
input window. Where the 5~V rail's load sits permanently below its boundary, this rail
\emph{crosses} it in operation: modem-sleep currents (\SIrange{13}{92}{\milli\ampere})
sit in PFM, and a Wi-Fi TX burst at \SI{355}{\milli\ampere} sits in CCM. The part
manages the PFM/PWM transition automatically and continuously --- the mechanism, its
\SI{450}{\milli\ampere} PFM peak-current clamp, and the zero-cross blocking are as
described in \cref{sec:5v-op} --- so the crossing imposes no component choice; its one
real consequence, transient droop at the mode transition, is bounded in
\cref{sec:3v3-cout}. The datasheet's typical-efficiency curve for
$V_{IN}=\SI{12}{\volt}$, $V_{OUT}=\SI{3.3}{\volt}$ reads approximately 87--90\,\% over
the \SIrange{50}{500}{\milli\ampere} span (\ds{AP632XX.pdf}{2}); $\eta=0.85$ is carried
as the conservative figure, matching \cref{sec:5v-op}. Average input current at full
load is then
\begin{align*}
  I_{IN} &= \frac{P_{\mathrm{out,max}}}{\eta\,V_N^{\min}}
          = \frac{1.19}{0.85 \times 8.1} = \SI{0.17}{\ampere}
  \quad (\SI{0.24}{\ampere}\ \text{at the \SI{0.5}{\ampere} capability point}),
\end{align*}
drawn from $V_{\mathrm{BUS}}$ through D411 --- the largest of the three logic-rail
loads on the bus budget, and the diode adds $I_{IN} V_F \approx \SI{70}{\milli\watt}$
(\SI{96}{\milli\watt} at the capability point) of series loss, about 5\,\% of
throughput at minimum bus voltage.

\subsection{Minimum on-time}\label{sec:3v3-ontime}

The binding corner is $V_N^{\max}=\SI{12.8}{\volt}$. At the $+6\,\%$ spread-spectrum
extreme ($f_{SW}=\SI{1.166}{\mega\hertz}$):
\begin{align*}
  t_{\mathrm{on}} &= \frac{D_{\min}}{f_{SW}}
    = \frac{0.258}{\SI{1.166}{\mega\hertz}} = \SI{221}{\nano\second}
    \quad\text{vs.}\quad t_{\mathrm{on,min}} = \SI{80}{\nano\second}
    \;\;(\cref{sec:ap632xx}),
\end{align*}
a $2.8\times$ margin; the PFM pulse on-time, $L\times\SI{0.45}{\ampere}/(V_N-V_{OUT}) =
\SI{223}{\nano\second}$ at the same corner, is the same figure. The margin is thinner
than the 5~V rail's $3.3$--$4.1\times$ (\cref{sec:5v-ontime}) because a 3.3~V output
yields a smaller duty cycle at the same input, but for the constraint to bind the input
would need $V_N = 3.3/(t_{\mathrm{on,min}} f_{SW}) \approx \SI{35}{\volt}$ --- the
absolute-maximum rating itself, unreachable from a 13~V bus. Match.

\subsection{Inductor}\label{sec:3v3-inductor}

The datasheet selection equation with its prescribed ripple band ($\Delta I_L$ =
30--50\,\% of the \SI{2}{\ampere} rating, i.e.\ \SIrange{0.6}{1.0}{\ampere};
\cref{sec:ap632xx}), at the maximum-ripple corner $V_N=\SI{12.8}{\volt}$:
\begin{align*}
  L\big|_{\Delta I_L = \SI{1.0}{\ampere}} &=
    \frac{3.3\,(12.8-3.3)}{12.8 \times 1.0 \times \SI{1.1}{\mega\hertz}}
    = \SI{2.23}{\micro\henry},
  &
  L\big|_{\Delta I_L = \SI{0.6}{\ampere}} &=
    \frac{31.35}{12.8 \times 0.6 \times \SI{1.1}{\mega\hertz}} = \SI{3.71}{\micro\henry},
\end{align*}
and \SIrange{1.78}{2.96}{\micro\henry} at \SI{8.1}{\volt}. The datasheet's own Table-2
BOM value for the AP63203 is \SI{3.9}{\micro\henry} (\cref{sec:ap632xx}). The populated
value is \SI{4.7}{\micro\henry} --- the \emph{AP63205's} Table-3 value, the same part
and footprint as the 5~V rail's L401. Unlike the 5~V rail, then, this converter is not
the reference BOM applied verbatim: the inductor deviates from the variant-specific
recommendation, evidently to share one line item across both rails. The consequences
all fall in the benign direction. Ripple is \SI{0.474}{\ampere} (23.7\,\% of
\SI{2}{\ampere}) at \SI{12.8}{\volt} and \SI{0.378}{\ampere} (18.9\,\%) at
\SI{8.1}{\volt} --- below the 30--50\,\% band, which lowers the CCM peak, lowers the
output ripple, and improves light-load efficiency (the datasheet's own endorsement of
larger inductance, apt for a clock that idles at modem-sleep currents for essentially
its whole life). The value remains inside the general \SIrange{2.2}{10}{\micro\henry}
recommendation. The cost is a slightly slower current slew into a load step, which the
transient check of \cref{sec:3v3-cout} absorbs with wide margin. Verdict: acceptable ---
a deliberate cross-rail reuse, noted as a deviation from the per-variant BOM table
rather than a mismatch.

\paragraph{Saturation current.} The peak inductor current in normal operation is
\begin{align*}
  I_{\mathrm{pk}} &= I_{\mathrm{load,max}} + \frac{\Delta I_L}{2}
    = 0.359 + 0.237 = \SI{0.60}{\ampere}
  \quad (\SI{0.74}{\ampere}\ \text{at \SI{0.5}{\ampere}}),
\end{align*}
against the datasheet's $\geq$35\,\% rule: $1.35 \times 0.359 = \SI{0.48}{\ampere}$
($1.35 \times 0.5 = \SI{0.68}{\ampere}$ at the capability point). Under the same
stated assumption as \cref{sec:5v-inductor} --- $I_{\mathrm{sat}} \geq \SI{3.0}{\ampere}$,
DCR $\leq\SI{50}{\milli\ohm}$ for a commodity 6045-class \SI{4.7}{\micro\henry} shielded
part --- the rating exceeds the worst normal peak by $\approx 5\times$ (against a real,
sub-ampere operating peak this time, not a PFM clamp), and covers the
\SI{3.1}{\ampere} maximum fault peak at or near hard saturation with hiccup protection
bounding dwell, exactly as argued there. Match under the stated assumption; the actual
part's rating is a confirm-at-BOM item shared with L401.

\subsection{Current-limit margin}\label{sec:3v3-ilim}

Against the \SI{2.5}{\ampere} minimum HS peak limit (\cref{sec:ap632xx}):
\begin{align*}
  \frac{I_{\mathrm{lim}}^{\min}}{I_{\mathrm{pk}}} &= \frac{2.5}{0.60} = 4.2\times
  \quad\text{(worst sustained load)}, \\
  \frac{I_{\mathrm{lim}}^{\min}}{I_{\mathrm{pk}}}\bigg|_{\SI{0.5}{\ampere}}
    &= \frac{2.5}{0.74} = 3.4\times
  \quad\text{(capability point)}.
\end{align*}
This is the check that was pro-forma on the 5~V rail and is meaningful here: the
operating peak is a quarter of the limit rather than a fiftieth, but $3.4\times$ at the
full Espressif sizing point is still comfortable, and start-up adds no interaction ---
the \SI{4}{\milli\second} internal soft-start charges the output-node capacitance
(\SI{66}{\micro\farad} nominal including the module-local C502) at
\begin{align*}
  I_{\mathrm{chg}} &= \SI{66}{\micro\farad} \times
    \frac{\SI{3.3}{\volt}}{\SI{4}{\milli\second}} = \SI{54}{\milli\ampere},
\end{align*}
which even summed with the full load leaves $6\times$ under the minimum limit. A shorted
rail lands in hiccup mode, as on the 5~V rail. Match.

\subsection{Output capacitance, transient response, and rail accuracy}\label{sec:3v3-cout}

The populated $2\times\SI{22}{\micro\farad}$ (C413/C414) equals the family BOM example
and the \SIrange{22}{68}{\micro\farad} recommendation (\cref{sec:ap632xx}). The same
capacitor assumptions as \cref{sec:5v-cout} are carried --- X5R/X7R ceramics, effective
ESR \SI{2.5}{\milli\ohm} paralleled --- with DC-bias derating milder at 3.3~V:
\emph{assume} $\geq 70\,\%$ retention, i.e.\ effective buck-local
$C_{\mathrm{out}} \geq \SI{30.8}{\micro\farad}$. C501
(\SI{100}{\nano\farad}) and C502 (\SI{22}{\micro\farad}) sit on the same net at the
module: the netlist does not label intent, but their values and placement match standard
practice for a second, load-local decoupling stage at a high-transient-current device,
and they are treated as that --- C502 adds $\approx\SI{15}{\micro\farad}$ effective
bulk at the module's pins beyond the converter's own capacitors.

\paragraph{Load-release overshoot.} The datasheet's transient bound applied in
\cref{sec:5v-cout}, at the worst CCM peak and \SI{50}{\milli\volt} permitted overshoot:
\begin{align*}
  C_{\mathrm{out}} &> \frac{L\,I_{\mathrm{pk}}^2}{(\Delta V + V_{OUT})^2 - V_{OUT}^2}
    = \frac{\SI{4.7}{\micro\henry} \times 0.60^2\,\si{\square\ampere}}
           {3.35^2 - 3.30^2\ \si{\square\volt}}
    = \frac{\SI{1.67}{\micro\joule}}{\SI{0.3325}{\square\volt}}
    = \SI{5.0}{\micro\farad}
\end{align*}
(\SI{7.7}{\micro\farad} at the \SI{0.74}{\ampere} capability-point peak) --- met
$6.1\times$ over by the derated buck-local capacitance alone, $4\times$ at the
capability point. The energy argument run forward gives the hard ceiling: dumping the
full \SI{0.60}{\ampere} inductor current into \SI{30.8}{\micro\farad} raises the rail by
only $\sqrt{3.3^2 + L I^2/C} - 3.3 \approx \SI{8}{\milli\volt}$. This matters more here
than on the 5~V rail because the ESP32's absolute maximum equals its recommended
maximum, \SI{3.6}{\volt}, with no headroom above it (\cref{sec:esp32}): even so, worst
regulation maximum plus worst overshoot, $3.33 + 0.008 \approx \SI{3.34}{\volt}$,
leaves \SI{260}{\milli\volt} to the ceiling. Match.

\paragraph{Load-step droop (Wi-Fi TX burst onset).} The bounding step is
idle-to-TX, $\Delta I = \SI{0.355}{\ampere}$. Two mechanisms stack. The inductor's
physical slew limit, worst at minimum input,
\begin{align*}
  \frac{dI_L}{dt} &= \frac{V_N - V_{OUT}}{L} = \frac{8.1-3.3}{\SI{4.7}{\micro\henry}}
    = \SI{1.02}{\ampere\per\micro\second}
  &&\Rightarrow&
  t_{\mathrm{slew}} &= \frac{0.355}{1.02} = \SI{0.35}{\micro\second},
\end{align*}
costs charge $\tfrac{1}{2}\,\Delta I\, t_{\mathrm{slew}} = \SI{62}{\nano\coulomb}$,
i.e.\ \SI{2}{\milli\volt} on \SI{30.8}{\micro\farad} --- negligible; the
\SI{4.7}{\micro\henry} reuse costs nothing measurable here. The dominant term is
control-loop and PFM$\to$PWM transition latency, for which the datasheet publishes no
bandwidth. \emph{Assumption:} the loop commands full inductor current within
\SI{5}{\micro\second} ($\approx$5~switching periods), a conservative figure for an
internally compensated \SI{1.1}{\mega\hertz} peak-current-mode buck. Then
\begin{align*}
  \Delta V_{\mathrm{droop}} &\approx \frac{\Delta I \cdot t_{\mathrm{resp}}}{C_{\mathrm{out}}}
      + \Delta I \cdot \mathrm{ESR}
    = \frac{0.355 \times \SI{5}{\micro\second}}{\SI{30.8}{\micro\farad}} + 0.355\times\SI{2.5}{\milli\ohm}
    \approx 58 + 1 \approx \SI{59}{\milli\volt}.
\end{align*}
Stacked on the worst regulation minimum: $3.27 - 0.059 = \SI{3.21}{\volt}$, a
\SI{210}{\milli\volt} margin over the ESP32's \SI{3.0}{\volt} floor. The conclusion is
robust to the assumption: even tripling the assumed latency to \SI{15}{\micro\second}
(\SI{173}{\milli\volt} droop) keeps the rail at \SI{3.10}{\volt}, still in
specification, and C502's placement at the module supplies the first microseconds
locally. Match.

\paragraph{Ripple and accuracy stack.} CCM ripple at full load,
$\Delta I_L/(8 f_{SW} C) + \Delta I_L\,\mathrm{ESR} \approx 1.8 + 1.2 \approx
\SI{3}{\milli\volt}$; PFM ripple at idle, bounded by one
\SI{450}{\milli\ampere}-peak pulse's charge (method of \cref{sec:5v-cout};
$Q = \SI{243}{\nano\coulomb}$ at the \SI{8.1}{\volt} corner) at
$Q/C + I_{\mathrm{pk}}\,\mathrm{ESR} \approx 7.9 + 1.1 \approx \SI{9}{\milli\volt}$
pk-pk. The full stack against the load's window:
\begin{align*}
  V_{\mathrm{rail}} &\in [\,3.27 - 0.009,\; 3.33 + 0.009\,]
    = [\,\SI{3.26}{\volt},\ \SI{3.34}{\volt}\,]
  \quad\text{vs.}\quad \text{VDD33} \in [\,\SI{3.0}{\volt},\ \SI{3.6}{\volt}\,],
\end{align*}
$\approx\SI{260}{\milli\volt}$ of margin on both sides (\cref{sec:esp32}), and the
transient droop above consumes under a third of the low-side margin. Match.

\subsection{Input capacitance}\label{sec:3v3-cin}

C411 (\SI{10}{\micro\farad} ceramic) on the OR node equals the family BOM example
against the ``greater than \SI{10}{\micro\farad}'' prose (\cref{sec:ap632xx}), the same
precision-not-flag status as C401. The buck's input RMS ripple current at the worst
point ($D$ nearest 0.5, i.e.\ $V_N = \SI{8.1}{\volt}$):
\begin{align*}
  I_{\mathrm{RMS}} &\approx I_{\mathrm{out}}\sqrt{D(1-D)}
    = 0.359\sqrt{0.407\times0.593} = \SI{0.18}{\ampere},
\end{align*}
equal to the half-of-max-load guidance figure and trivially inside any ceramic
\SI{10}{\micro\farad}'s ampere-class rating. Input ripple voltage,
$I_{\mathrm{out}}D(1-D)/(f_{SW}C) \approx \SI{8}{\milli\volt}$ nominal
($\approx\SI{13}{\milli\volt}$ after DC-bias derating at the 12.8~V corner), is
absorbed at the node --- and, usefully, the buck's pulsed input current circulates
between C411 and the IC, so D411 carries the smoothed average
(\SI{0.24}{\ampere} worst) rather than the switching pulses. As with C401, the
schematic does not record the voltage rating: on a node reaching \SI{12.8}{\volt} it
must be $\geq\SI{16}{\volt}$, \SI{25}{\volt} preferred for derating headroom ---
confirm at BOM finalization. C412 (\SI{100}{\nano\farad}, SW-to-BST) matches the
recommended bootstrap value. Match throughout.

\subsection{Diode-OR input path and UART-powered operation}\label{sec:3v3-diodeor}

The OR node (C411, U411 IN/EN) is fed by two SOD-123 Schottky diodes: D411, anode on
$V_{\mathrm{BUS}}$, and D412, anode on \texttt{VCC\_UART} --- pin~3 of J501, the 6-pin
UART header (1/2 GND, 3 VCC, 4 RXD0, 5 TXD0, 6 unconnected). The netlist confirms
\texttt{VCC\_UART} touches exactly these two points; it reaches nothing else on the
board. The node therefore follows whichever source is higher, less one forward drop,
with the other diode reverse-biased --- a protective OR, not a sum. Both directions are
worked through below; common USB-UART adapters (FTDI/CP2102-class) ship configured for
either \SI{3.3}{\volt} or \SI{5}{\volt} on their VCC pin, often jumper-selectable, so
both sub-cases of the UART-powered scenario are treated.

\paragraph{Normal operation ($V_{\mathrm{BUS}}$ present).} The node sits at
\SIrange{8.1}{12.8}{\volt}, so D412's cathode is \SIrange{3}{13}{\volt} above any
adapter VCC (\SIrange{0}{5}{\volt}): D412 is reverse-biased and no forward path exists
from this board's power onto an attached adapter's VCC pin. The only residual coupling
is D412's reverse leakage --- \si{\micro\ampere}-class under the stated Schottky
assumption --- sourced \emph{into} the adapter's own regulated output, which is
inconsequential. Plainly: the board cannot back-power a UART adapter's VCC pin. (For
completeness on the same header: the activity-LED feeds R502/R503 can source
$\lesssim\SI{1}{\milli\ampere}$, series-limited by \SI{1.5}{\kilo\ohm}, into an
\emph{unpowered} adapter's data pins through their ESD structures --- a negligible
stress, noted only so the ``no backfeed'' statement is scoped to the VCC pin
precisely.)

\paragraph{UART-only operation, adapter set to 3.3~V.} With
$V_{\mathrm{BUS}}$ absent, D411 is reverse-biased (blocking the node from back-feeding
the bus; its \si{\micro\ampere} leakage cannot lift the bus net against the flyback's
UVLO divider) and D412 conducts. The load is the download/console-mode estimate of
\cref{sec:3v3-load}, $\approx\SI{50}{\milli\ampere}$; in any pass-through mode the
buck's input current equals its output current, so $V_F(\mathrm{D412}) \approx
\SI{0.3}{\volt}$ and
\begin{align*}
  V_N &= 3.30 - 0.30 = \SI{3.0}{\volt}.
\end{align*}
This fails before regulation is even in question: the AP63203's VIN UVLO thresholds are
\SI{3.50}{\volt} rising, $\approx\SI{3.06}{\volt}$ falling (\cref{sec:ap632xx}), so at
\SI{3.0}{\volt} the converter never enables and the +3V3 rail stays at zero. No
tolerance rescues it --- even the friendliest corner (adapter at $+5\,\%$,
\SI{3.47}{\volt}; $V_F$ at its \SI{0.25}{\volt} floor) reaches only
\SI{3.22}{\volt}, still below the rising threshold, and the recommended operating
minimum is \SI{3.8}{\volt} regardless. The counterfactual is also checked so the
conclusion does not rest on the UVLO figure alone: \emph{were} the part running, it
would be in the qualitative LDO mode of \cref{sec:ap632xx} (no numeric dropout spec
exists), with the HS FET nearly always on and the output estimated through the
resistive path:
\begin{align*}
  V_{OUT} &\approx V_N - I\,(R_{DS(\mathrm{on}),HS} + \mathrm{DCR})
    = 3.00 - 0.05\,(0.125 + 0.050) = \SI{2.99}{\volt}
\end{align*}
at \SI{50}{\milli\ampere} (\SI{2.98}{\volt} at \SI{100}{\milli\ampere}) --- below the
ESP32's \SI{3.0}{\volt} recommended minimum by \SIrange{10}{20}{\milli\volt}
(negative margin) and \SI{280}{\milli\volt} below the regulation minimum. The
\SI{3.3}{\volt} adapter configuration therefore fails both ways: the realistic outcome
is a dead rail, and even the most favorable counterfactual is out of specification.
Flagged: incompatible.

\paragraph{UART-only operation, adapter set to 5~V.} At the
\SI{50}{\milli\ampere} programming load, $I_{IN} \approx \SI{40}{\milli\ampere}$
($\eta\approx0.88$ at this light load and low ratio), $V_F \approx \SI{0.3}{\volt}$:
\begin{align*}
  V_N &= 5.0 - 0.3 = \SI{4.7}{\volt},
\end{align*}
above the \SI{3.50}{\volt} rising UVLO and the \SI{3.8}{\volt} recommended minimum,
with $D = 3.3/4.7 = 0.70$ --- deep inside the part's near-100\,\% duty capability, so
the buck regulates normally at \SIrange{3.27}{3.33}{\volt}, identical to bus-powered
operation. The margin survives every stacked tolerance: a USB-derived
``\SI{5}{\volt}'' sagging to \SI{4.5}{\volt} still yields $V_N \approx
\SI{4.1}{\volt} > \SI{3.8}{\volt}$, and even firmware that lights the radio on UART
power ($I_{IN} \approx \SI{0.30}{\ampere}$, $V_F \approx \SI{0.4}{\volt}$, $V_N
\approx \SI{4.6}{\volt}$) stays regulated --- the limit in that corner is the
adapter's own current capability, not this converter. The \SI{5}{\volt} setting is
therefore not merely the safer choice; it is the \emph{only} working one. The
intuitive pairing --- ``3.3~V board, 3.3~V adapter'' --- is exactly backwards here,
because the rail is regulated on-board and the regulator needs UVLO clearance and
headroom above its own output.

\paragraph{Conclusion and user-guide requirement.} This board's UART header requires
the adapter's VCC set or wired to \SI{5}{\volt}; a \SI{3.3}{\volt}-VCC adapter leaves
the board dead when USB-C power is absent. This is an operational requirement for the
hardware user guide's flashing/console procedure, stated there as: \emph{use an
adapter supplying \SI{5}{\volt} on VCC with \SI{3.3}{\volt} logic levels on TXD/RXD}
--- the second clause because the data pins connect directly to ESP32 GPIOs whose
recommended input maximum is $V_{DD}+\SI{0.3}{\volt}$ (\cref{sec:esp32}); many
FTDI-style cables offer precisely this 5~V-VCC/3.3~V-I/O combination. Attaching or
detaching the adapter while USB-C power is present is safe in either order: the OR
simply follows the higher source, which is always $V_{\mathrm{BUS}}$.

\paragraph{Component ratings on the OR node.} Across all scenarios the node spans
\SIrange{0}{12.8}{\volt}: U411's IN/EN absolute maximum of \SI{35}{\volt}
(\cref{sec:ap632xx}) is met $2.7\times$ over, and C411's required $\geq\SI{16}{\volt}$
rating is the confirm-at-BOM item of \cref{sec:3v3-cin}. D411 sees at most
$\approx\SI{5}{\volt}$ reverse (5~V adapter, bus absent) and \SI{0.24}{\ampere}
forward average (\SI{96}{\milli\watt} --- inside a SOD-123's typical
\SIrange{400}{500}{\milli\watt} envelope, under the stated assumption of an
$I_F\geq\SI{0.5}{\ampere}$-class part). D412 is the binding case: it blocks the full
node voltage, up to \SI{12.8}{\volt}, whenever the bus is up and no adapter drives
VCC, so the part must be rated $V_R \geq \SI{16}{\volt}$ (\SI{20}{\volt}-plus is the
commodity floor for SOD-123 Schottkys, but the schematic records no part number) ---
confirm at BOM finalization alongside the inductor and capacitor ratings.

\subsection{Pull-up loads, soft start, and dissipation}\label{sec:3v3-house}

\paragraph{PGOOD pull-up (R206).} The LM5155's open-drain PGOOD pin is pulled up by
R206 (\SI{100}{\kilo\ohm}) to this rail, with ESP32 GPIO15 --- an input, contributing
only leakage --- reading the same node. The pin-level checks (sink current
\SI{33}{\micro\ampere} vs.\ the \SI{1}{\milli\ampere} limit,
$\SI{100}{\kilo\ohm} \geq \SI{10}{\kilo\ohm}$ recommended, voltage ratings) were
completed in \cref{sec:hv-house} and stand. What remains for this rail is the load
contribution and its polarity. The pin releases to high impedance when the flyback's
output is good and its VCC and UVLO conditions are met, and pulls low otherwise
(\ds{lm5155.pdf}{22}) --- i.e.\ the pull-up conducts its
$3.3/\SI{100}{\kilo\ohm} = \SI{33}{\micro\ampere}$ only during start-up and fault
windows, and essentially nothing (GPIO leakage, \si{\nano\ampere}) in normal operation.
Either way it is four decades under the design load: negligible. GPIO15 accordingly
reads high for power-good --- a coherent arrangement, since the signal's only consumer
is powered by the same rail that supplies the pull-up, so the signal is valid exactly
when there is something running to read it.

\paragraph{EN pull-up (R501).} The R501/C503 network is Espressif's reference circuit
verbatim (\cref{sec:esp32}). Steady-state current into the EN input is leakage-level;
R501 conducts measurably only while charging C503 at power-up or while the reset
button holds EN low --- \SI{330}{\micro\ampere} peak, decaying with
$\tau = \SI{10}{\milli\second}$, or held during a button press. Zero contribution to
the steady-state budget; negligible.

\paragraph{Soft start and dissipation.} Soft start is the fixed internal
\SI{4}{\milli\second} ramp; the inrush interaction was dismissed in
\cref{sec:3v3-ilim}. Converter dissipation at the conservative $\eta=0.85$:
\begin{align*}
  P_{\mathrm{diss}} &= P_{\mathrm{out,max}}\left(\frac{1}{\eta}-1\right)
    = 1.19 \times 0.176 = \SI{0.21}{\watt}, &
  \Delta T &= \theta_{JA}P_{\mathrm{diss}} = 89 \times 0.21 \approx \SI{19}{\kelvin},
\end{align*}
rising to \SI{0.29}{\watt} and $\Delta T \approx \SI{26}{\kelvin}$ at the
\SI{0.5}{\ampere} capability point --- $T_J \leq \SI{66}{\celsius}$ at a
\SI{40}{\celsius} ambient, far below the \SI{125}{\celsius} design target
(\cref{sec:ap632xx}). Wi-Fi TX at 100\,\% duty is already the sustained rating, so no
duty-cycle relief is needed to make this conclusion. Add D411's \SI{96}{\milli\watt}
and the thermal picture on this rail remains untroubled.

\subsection{Summary}\label{sec:3v3-summary}

Assumptions carried through this section: generic SOD-123 Schottky forward drops
\SIrange{0.3}{0.45}{\volt} over \SIrange{50}{500}{\milli\ampere} with
\si{\micro\ampere}-class leakage; inductor $I_{\mathrm{sat}}\geq\SI{3.0}{\ampere}$ and
DCR $\leq\SI{50}{\milli\ohm}$ (6045-class, shared with L401); output ceramics retaining
$\geq70\,\%$ at 3.3~V bias, effective ESR \SI{2.5}{\milli\ohm}; converter efficiency
0.85 (floor under the datasheet's $\approx$87--90\,\% curve); control-loop response to
a load step within \SI{5}{\micro\second} (conclusion shown robust to $3\times$ that);
LED forward voltage \SI{2.0}{\volt}; download-mode MCU draw
$\approx\SI{50}{\milli\ampere}$ (bounded at \SI{100}{\milli\ampere}). Each appears at
its point of use above.

\begin{longtable}{@{}p{0.26\textwidth}p{0.30\textwidth}p{0.22\textwidth}p{0.14\textwidth}@{}}
\toprule
Parameter & Derived requirement & Populated & Verdict \\
\midrule
\endfirsthead
\toprule
Parameter & Derived requirement & Populated & Verdict \\
\midrule
\endhead
\bottomrule
\endlastfoot
Load current & \SI{355}{\milli\ampere} MCU $+$ $\leq$\SI{4}{\milli\ampere}
  housekeeping $\approx$ \SI{0.36}{\ampere} sustained; capability
  $\geq\SI{0.5}{\ampere}$ (Espressif) & \SI{2}{\ampere}-rated AP63203 &
  Match ($4\times$) \\
Rail vs.\ MCU window & VDD33 \SIrange{3.0}{3.6}{\volt}; abs max $=$ rec max &
  \SIrange{3.26}{3.34}{\volt} incl.\ ripple & Match
  ($\approx\SI{260}{\milli\volt}$ both sides) \\
Input window (normal) & inside \SIrange{3.8}{32}{\volt} recommended & OR node
  \SIrange{8.1}{12.8}{\volt} ($V_{\mathrm{BUS}}-V_F$) & Match \\
Duty cycle & 0.258--0.407; never crosses 0.5 & --- & Verified \\
Conduction mode & CCM boundary \SIrange{189}{237}{\milli\ampere}: PFM at idle,
  CCM in TX bursts; automatic transition & --- & Verified (benign) \\
Minimum on-time & \SI{221}{\nano\second} vs.\ \SI{80}{\nano\second}; binds only
  at $V_N\approx\SI{35}{\volt}$ & --- & Match ($2.8\times$) \\
Inductance L411 & Eq.\ window \SIrange{2.2}{3.7}{\micro\henry}; Table-2 BOM
  \SI{3.9}{\micro\henry} & \SI{4.7}{\micro\henry} (5~V rail's part reused) &
  Acceptable deviation (ripple 19--24\,\%, benign) \\
Inductor $I_{\mathrm{sat}}$ & $\geq\SI{0.68}{\ampere}$ (rule at
  \SI{0.5}{\ampere} point); fault $\leq\SI{3.1}{\ampere}$ & 6045-class, rating
  unrecorded (assumed $\geq\SI{3}{\ampere}$) & Match under assumption;
  \textbf{confirm part} \\
Current-limit margin & peak \SI{0.60}{\ampere} (\SI{0.74}{\ampere} at
  capability point) vs.\ \SI{2.5}{\ampere} min & --- & Match ($\geq3.4\times$) \\
Output capacitance C413/C414 & $\geq\SI{5.0}{\micro\farad}$ transient bound;
  \SIrange{22}{68}{\micro\farad} recommended & $2\times\SI{22}{\micro\farad}$
  ($+$ C502 at module) & Match \\
TX-burst droop & $\lesssim\SI{59}{\milli\volt}$ vs.\ \SI{270}{\milli\volt}
  budget to \SI{3.0}{\volt}; robust to $3\times$ latency & --- & Match \\
Output ripple & $\leq\SI{9}{\milli\volt}$ pk-pk PFM / $\approx\SI{3}{\milli\volt}$
  CCM & --- & Derived; $\ll$ margins \\
Input capacitor C411 & $\geq\SI{10}{\micro\farad}$; RMS $\geq\SI{0.18}{\ampere}$;
  $\geq\SI{16}{\volt}$ rating & \SI{10}{\micro\farad} (= BOM example; rating
  unrecorded) & Match; \textbf{confirm rating} \\
BST capacitor C412 & \SI{100}{\nano\farad} SW-to-BST & \SI{100}{\nano\farad} &
  Match \\
Backfeed to adapter VCC & none permissible & D412 orientation blocks; leakage
  \si{\micro\ampere}-class & Verified (none) \\
UART power, \SI{3.3}{\volt} adapter & node \SI{3.0}{\volt} $<$
  \SI{3.50}{\volt} UVLO $\Rightarrow$ rail dead; counterfactual LDO output
  \SI{2.98}{\volt}, below MCU minimum & --- & \textbf{Flagged: incompatible} \\
UART power, \SI{5}{\volt} adapter & node \SIrange{4.1}{4.7}{\volt}; normal
  regulation & --- & Match; \textbf{required configuration} (user-guide item,
  with 3.3~V logic levels) \\
D412 reverse rating & $\geq\SI{16}{\volt}$ (blocks \SI{12.8}{\volt}) & generic
  Schottky, part unrecorded & Match under assumption; \textbf{confirm part} \\
D411 forward path & $I_F$ avg $\leq\SI{0.24}{\ampere}$, \SI{96}{\milli\watt} &
  generic SOD-123 Schottky & Match under assumption \\
PGOOD pull-up R206 & load \SI{33}{\micro\ampere}, not-good windows only
  (high-Z when good) & \SI{100}{\kilo\ohm}; pin checks in \cref{sec:hv-house} &
  Match; negligible \\
EN pull-up R501 & transient only (\SI{330}{\micro\ampere} pk,
  $\tau=\SI{10}{\milli\second}$) & \SI{10}{\kilo\ohm} (= Espressif reference) &
  Match; negligible \\
Indicator LEDs R411/R502/R503 & \SI{0.87}{\milli\ampere} continuous $+$
  $\leq\SI{1.6}{\milli\ampere}$ UART activity & \SI{1.5}{\kilo\ohm} each
  (D413 value field mislabeled ``RX'') & Negligible; fix label \\
Module decoupling C501/C502 & load-local second stage at a burst load &
  \SI{100}{\nano\farad} $+$ \SI{22}{\micro\farad} & Match (role inferred) \\
Soft start / inrush & \SI{54}{\milli\ampere} charging $+$ load $\ll$ limit &
  internal \SI{4}{\milli\second} & Match \\
Converter dissipation & $\leq\SI{0.29}{\watt}$, $\Delta T \leq \SI{26}{\kelvin}$
  vs.\ \SI{125}{\celsius} target & TSOT26,
  $\theta_{JA}=89$~\si{\celsius\per\watt} & Match (large margin) \\
PSRAM active-mode adder & unquantified above \SI{355}{\milli\ampere}
  (\cref{sec:esp32}) & module SKU unrecorded & Open item (absorbed by
  $4\times$ headroom) \\
\end{longtable}

Three findings distinguish this rail from its 5~V sibling. First, the inductor is a
deliberate deviation from the variant-specific reference BOM (\SI{4.7}{\micro\henry}
where Table~2 lists \SI{3.9}{\micro\henry}) --- acceptable, with every consequence
benign at this rail's loads. Second, and operationally most significant: the UART
header's diode-OR feed works \emph{only} with a \SI{5}{\volt}-VCC adapter; a
\SI{3.3}{\volt}-VCC adapter cannot lift the buck's input past its own UVLO and leaves
the board dead when USB-C power is absent. That requirement (5~V VCC, 3.3~V logic
levels) belongs in the user guide's flashing procedure as a hard prerequisite, not a
footnote. Third, the transient checks that were pro-forma on the 5~V rail are real here
but all pass with room: worst droop uses under a third of the low-side voltage margin,
and the current-limit margin at Espressif's full \SI{0.5}{\ampere} sizing point is
$3.4\times$. The confirm-at-BOM items --- inductor saturation rating, C411 voltage
rating, D412 reverse rating --- are unrecorded ratings on generic parts, the same
status as their counterparts in \cref{sec:5v-converter}, plus one cosmetic fix (D413's
value field) and one inherited open item (the PSRAM active-mode current adder).


\end{document}
